\documentclass[11pt,letterpaper]{article}

\usepackage[spanish,es-noshorthands]{babel}
\usepackage{fontspec}
\setmainfont{Source Serif Pro}
\usepackage{microtype}
\usepackage{amsmath}
\usepackage{geometry}
\usepackage{booktabs}
\usepackage{tabularx}
\usepackage{enumitem}
\usepackage{xcolor}
\usepackage{csquotes}
\usepackage{float}
% El documento tiene 36 cuadros y 24 figuras. Sin relajar los parámetros de
% colocación, la cola de flotantes se desborda y las tablas se acumulan al
% final; la barrera por sección impide además que un flotante cruce de sección.
\usepackage[section]{placeins}
\renewcommand{\topfraction}{0.92}
\renewcommand{\bottomfraction}{0.75}
\renewcommand{\textfraction}{0.06}
\renewcommand{\floatpagefraction}{0.68}
\setcounter{topnumber}{3}
\setcounter{bottomnumber}{2}
\setcounter{totalnumber}{5}
\usepackage{tikz}
\usetikzlibrary{arrows.meta,positioning,shapes.geometric,calc,fit,backgrounds,shadows}
\usepackage{pgfplots}
\pgfplotsset{compat=1.18}
% Los años (1995, 2026) no llevan separador de millar
\pgfplotsset{tick label style={/pgf/number format/1000 sep={}}}
\usepackage{tocloft}
\usepackage{xurl}
\usepackage[hidelinks]{hyperref}

% Índice general compacto: el interlineado de \parskip lo haría desbordar
\renewcommand{\cfttoctitlefont}{\large\bfseries}
\setlength{\cftbeforesecskip}{0.4em}
\setlength{\cftbeforesubsecskip}{0.05em}
\setlength{\cftaftertoctitleskip}{0.9em}
\renewcommand{\cftsecaftersnum}{.}  % "9." como en los títulos de sección
\setlength{\cftsecnumwidth}{1.9em}
\usepackage[backend=biber,style=apa,sorting=nyt]{biblatex}
\DeclareLanguageMapping{spanish}{spanish-apa}

\geometry{margin=2.5cm}
\setlength{\emergencystretch}{1em}
\setlength{\parindent}{0pt}
\setlength{\parskip}{0.65em}
\setlist{nosep}
\newcommand{\pct}{\char`\%}  % \% sin el espacio fino de babel

% Apartado sin numerar, con entrada propia en el índice general
\newcommand{\subsec}[1]{%
  \subsection*{#1}%
  \phantomsection\addcontentsline{toc}{subsection}{#1}}
\addbibresource{referencias.bib}
\AtBeginBibliography{\sloppy}

% --------------------------------------------------------------------------
% Estilos gráficos comunes
% --------------------------------------------------------------------------
\definecolor{capa}{HTML}{2F5D8C}
\definecolor{trans}{HTML}{8C5A2F}
\definecolor{acento}{HTML}{9C2B2B}
\definecolor{jugA}{HTML}{1B7F9E}    % jugador A: Terraform
\definecolor{jugB}{HTML}{6F8F14}    % jugador B: Bicep
\definecolor{jugC}{HTML}{9C2B2B}    % jugador C: Pulumi
\definecolor{propia}{HTML}{2F5D8C}  % ruta propuesta para la empresa
\definecolor{descart}{HTML}{7A7A7A} % opción evaluada y descartada

\definecolor{opmobjc}{HTML}{1E7A1E}   % contorno de los objetos (OPM)
\definecolor{opmprocc}{HTML}{2E3FA8}  % contorno de los procesos (OPM)
\definecolor{dsmarq}{HTML}{2F5D8C}    % marcas de la DSM de arquitectura
\definecolor{dsmpro}{HTML}{8C5A2F}    % marcas de la DSM de cartera

\tikzset{
  % ---------------- Notación OPM (ISO 19450) ----------------
  % Objeto: rectángulo de contorno verde.
  opmobj/.style   = {draw=opmobjc, very thick, rectangle, fill=opmobjc!7,
                     font=\footnotesize\bfseries, align=center,
                     inner sep=4pt, minimum height=8mm, text width=26mm},
  opmobjn/.style  = {opmobj, text width=},
  % Proceso: elipse de contorno azul.
  opmproc/.style  = {draw=opmprocc, very thick, ellipse, fill=opmprocc!7,
                     font=\footnotesize\bfseries, align=center, text=black,
                     inner sep=1pt, minimum height=9mm, minimum width=28mm},
  % Estado: rectángulo de esquinas redondeadas, dentro del objeto.
  opmstate/.style = {draw=opmobjc, thick, rounded corners=3pt, fill=white,
                     font=\scriptsize, inner sep=2pt, minimum height=4.5mm},
  % Cosa física: sombra proyectada. Cosa ambiental: contorno discontinuo.
  opmfis/.style   = {drop shadow={shadow xshift=0.8mm, shadow yshift=-0.8mm,
                     opacity=1, fill=black!45}},
  opmamb/.style   = {densely dashed},
  % Enlaces procedimentales.
  instr/.style    = {thick, -{Circle[open,length=5pt]}},   % instrumento
  agente/.style   = {thick, -{Circle[length=5pt]}},        % agente
  efecto/.style   = {thick, {Stealth[length=5.5pt]}-{Stealth[length=5.5pt]}},
  flecha/.style   = {thick, -{Stealth[length=5pt]}},       % consumo/resultado
  invoca/.style   = {thick, densely dashed, -{Stealth[length=5pt]}},
  % Diagramas de bloques
  bloque/.style   = {draw, thick, rounded corners=2pt, align=center,
                     font=\footnotesize, inner sep=5pt},
  etiq/.style     = {font=\scriptsize\itshape, align=center},
}

% --------------------------------------------------------------------------
% Símbolos estructurales de OPM.
% La base del triángulo se apoya en el lado del todo, del general o del
% exhibidor, y el vértice apunta hacia la parte, la especialización o el
% atributo. #1 = rotación en grados (0 = base arriba, vértice abajo);
% #2 = coordenada previamente definida sobre la que se centra el símbolo.
% --------------------------------------------------------------------------
\newcommand{\opmtri}[4]{% #3 = semiancho, #4 = relleno
  \begin{scope}[rotate around={#1:(#2)}]
    \path[draw,very thick,fill=#4] ([shift={(-#3,#3)}]#2) --
      ([shift={(#3,#3)}]#2) -- ([shift={(0mm,-#3)}]#2) -- cycle;
  \end{scope}}
\newcommand{\opmagg}[2]{\opmtri{#1}{#2}{2.1mm}{black}}   % agregación-participación
\newcommand{\opmgen}[2]{\opmtri{#1}{#2}{2.1mm}{white}}   % generalización-especialización
\newcommand{\opmexh}[2]{% exhibición-caracterización: triángulo dentro de triángulo
  \opmtri{#1}{#2}{2.6mm}{white}%
  \begin{scope}[rotate around={#1:(#2)}]
    \path[fill=black] ([shift={(-1.1mm,0.9mm)}]#2) --
      ([shift={(1.1mm,0.9mm)}]#2) -- ([shift={(0mm,-1.1mm)}]#2) -- cycle;
  \end{scope}}

% --------------------------------------------------------------------------
% Matriz de estructura de dependencia (DSM).
% Rejilla cuadrada de lado #1 celdas, con celda de \dsmu de lado.
% --------------------------------------------------------------------------
\newlength{\dsmu}\setlength{\dsmu}{6.4mm}

\title{\textbf{Hoja de ruta tecnológica de la infraestructura como código}\\
\large Alcance del sistema, matriz de estructura de dependencia, modelo de
objetos y procesos, figuras de mérito y su evolución, posicionamiento frente a
la competencia, modelo técnico y financiero, cartera de proyectos y
declaración de estrategia}
\author{Carlos Luis Rojas Aragonés\\
Gestión del Desarrollo Tecnológico}
\date{7 de septiembre de 2026}

\begin{document}

\maketitle

\begin{abstract}
\noindent\small
La infraestructura como código (\emph{Infrastructure as Code}, IaC) sustituye
la configuración manual de servidores, redes y servicios por una
especificación versionada que un motor automático traduce en operaciones sobre
la interfaz de programación de un proveedor. Este trabajo delimita el alcance
de la tecnología, describe su principio de funcionamiento (la conciliación
continua entre un estado deseado y un estado real), desglosa las siete capas y
los seis módulos transversales de un sistema completo, y ordena sus
interdependencias en dos matrices de estructura de dependencia, una de
arquitectura y otra de cartera de proyectos, cuyo particionado produce la
secuencia de implantación. Formaliza después el conjunto mediante un modelo de
objetos y procesos según la metodología OPM, con su diagrama de sistema, el
despliegue estructural de la especificación, la ampliación en detalle del
proceso principal y el lenguaje OPL de los tres diagramas. Define después ocho
figuras de mérito con su unidad, su procedimiento de medida y sus anclas de
normalización, y aplica sobre ellas los tres modelos de evolución tecnológica:
una curva S que sitúa el punto de inflexión de la tecnología en 2012 y el
estado actual en el 92 \pct{} de su límite, un ajuste al modo de la ley de
Moore que encuentra la superficie declarable duplicándose cada 1,8 a 2,2 años,
y un frente de Pareto en el que dos de las seis posiciones consideradas
resultan dominadas. Sobre las mismas figuras se construye el posicionamiento
frente a la competencia, con un gráfico vectorial que sitúa a Terraform, Bicep
y Pulumi en un plano de variación de valor y de coste. Completan la hoja de
ruta la correspondencia con seis motivaciones estratégicas, un modelo técnico
de cuatro ecuaciones calibradas que reproducen los objetivos comprometidos, un
modelo financiero con un valor actual neto de 743,3 miles de euros y una
exposición máxima de 335,6, una cartera de ocho proyectos con cinco prototipos
y cuatro puertas de decisión, el panorama de publicaciones y patentes
verificadas una a una, la declaración de estrategia y la evaluación de madurez
del propio documento. Todas las figuras son de elaboración propia.
\end{abstract}

\clearpage
\setcounter{tocdepth}{2}
{\setlength{\parskip}{0pt}\tableofcontents}
\clearpage

% ==========================================================================
\section{Objeto y delimitación del alcance}
\label{sec:alcance}
% ==========================================================================

\subsec{01. El encargo}

El documento recoge entregas encadenadas de una misma hoja de ruta
tecnológica, y por eso no se organiza por actividades sino por los elementos
del índice recomendado para una hoja de ruta. El Cuadro~\ref{tab:atra} recoge
esa correspondencia.

La primera entrega, el alcance de la tecnología, cubre los tres primeros
elementos. El elemento 1, la descripción de la tecnología, ocupa las
secciones~\ref{sec:principio} a \ref{sec:modulo}, con el resumen funcional y
los diagramas de capas y de módulos. El elemento 2, la matriz de estructura de
dependencia, ocupa la sección~\ref{sec:dsm} con dos matrices, una de los
componentes del sistema y otra de los proyectos de la cartera, y con el
particionado que convierte a ambas en una secuencia de implantación. El
elemento 3, el modelo de objetos y procesos, ocupa la
sección~\ref{sec:opm}, con el diagrama de sistema, el despliegue estructural
de la especificación, la ampliación en detalle del proceso principal y el OPL
de los tres diagramas.

La segunda entrega pide analizar la competencia mediante un gráfico vectorial
que incluya al menos dos competidores. Corresponde al elemento 6 y ocupa la
sección~\ref{sec:comp}, que compara tres jugadores, Terraform, Bicep y
Pulumi, con la trayectoria propuesta para la empresa. Su sitio natural es
inmediatamente después de la sección~\ref{sec:fom}, porque las coordenadas del
gráfico se construyen a partir de las figuras de mérito definidas allí.

La tercera entrega pide definir las figuras de mérito de la tecnología y
aplicarlas para observar su evolución. Es el elemento 4, y ocupa dos
secciones: la~\ref{sec:fom} define las ocho figuras, sus anclas de
normalización y sus objetivos, y la~\ref{sec:evol} aplica los tres modelos de
evolución tecnológica disponibles, la curva S, el ajuste al modo de la ley de
Moore y el frente de Pareto. Con esta entrega el documento completa además los
seis elementos del índice que quedaban sin atender: las motivaciones
estratégicas (sección~\ref{sec:motiv}), el modelo técnico
(sección~\ref{sec:tecnico}), el modelo financiero
(sección~\ref{sec:financiero}), la cartera de proyectos y prototipos
(sección~\ref{sec:portfolio}), las publicaciones y patentes
(sección~\ref{sec:publica}), la declaración de estrategia
(sección~\ref{sec:estrategia}) y la evaluación de madurez
(sección~\ref{sec:madurez}). Esta última mide con qué solidez está atendido
cada uno, y su conclusión no es complaciente: los doce elementos están
completos en su forma y débilmente sustentados en sus cifras, porque ninguna
de las ocho figuras de mérito está instrumentada todavía.

\begin{table}[htbp]
  \centering
  \footnotesize
  \caption{Correspondencia entre los doce elementos del índice recomendado
  para una hoja de ruta tecnológica y las secciones de este documento. Los
  doce quedan atendidos; la sección~\ref{sec:madurez} evalúa con qué
  solidez.}
  \label{tab:atra}
  \begin{tabularx}{\textwidth}{@{}l>{\raggedright\arraybackslash}X
    >{\raggedright\arraybackslash}p{0.30\textwidth}@{}}
    \toprule
    \textbf{N.\textsuperscript{o}} & \textbf{Elemento} &
      \textbf{Dónde se atiende} \\
    \midrule
    1 & Descripción de la tecnología &
      Secciones \ref{sec:alcance} a \ref{sec:modulo} \\
    2 & Matriz de estructura de dependencia (DSM) &
      Sección \ref{sec:dsm} \\
    3 & Modelo con metodología OPM &
      Sección \ref{sec:opm} \\
    4 & Figuras de mérito (FOM) &
      Secciones \ref{sec:fom} y \ref{sec:evol} \\
    5 & Correspondencia con las motivaciones estratégicas &
      Sección \ref{sec:motiv} \\
    6 & Posicionamiento frente a la competencia &
      Sección \ref{sec:comp} \\
    7 & Modelo técnico & Sección \ref{sec:tecnico} \\
    8 & Modelo financiero & Sección \ref{sec:financiero} \\
    9 & Portfolio de proyectos y prototipos de I+D &
      Sección \ref{sec:portfolio} y Cuadro \ref{tab:dsmproy} \\
    10 & Publicaciones, presentaciones y patentes clave &
      Sección \ref{sec:publica} \\
    11 & Declaración de estrategia tecnológica &
      Sección \ref{sec:estrategia} \\
    12 & Evaluación del nivel de madurez de la hoja de ruta &
      Sección \ref{sec:madurez} \\
    \bottomrule
  \end{tabularx}
\end{table}

\subsec{02. Qué se entiende aquí por sistema completo}

Conviene fijar el término antes de dibujar nada, porque \enquote{infraestructura
como código} se usa con dos alcances muy distintos. En el sentido estrecho
designa una herramienta: un fichero declarativo y un binario que lo aplica. En
el sentido amplio, que es el que aquí interesa, designa un sistema
sociotécnico completo, con su cadena de herramientas, su modelo de estado, sus
controles de gobierno y el equipo que lo opera.

Un sistema completo de IaC es aquel que satisface cuatro propiedades
simultáneas \parencite{morris2021}:

\begin{enumerate}
  \item \textbf{Especificación externalizada.} La configuración deseada de la
    infraestructura vive fuera de la infraestructura, en artefactos de texto
    legibles y versionados, no en la memoria de un operador ni en el estado
    acumulado de las máquinas.
  \item \textbf{Reproducibilidad.} Aplicar la misma especificación sobre un
    sustrato vacío produce un resultado equivalente, con independencia de
    quién la aplique y de cuándo lo haga.
  \item \textbf{Idempotencia.} Aplicar la especificación sobre un sistema que
    ya la cumple no produce ningún cambio, y aplicarla \emph{n} veces tiene el
    mismo efecto que aplicarla una vez.
  \item \textbf{Conciliación.} El sistema es capaz de detectar la divergencia
    entre lo especificado y lo realmente desplegado, y de corregirla.
\end{enumerate}

La literatura académica sobre la materia, todavía escasa en comparación con su
difusión industrial, se ha concentrado sobre todo en la calidad y los defectos
del propio código de infraestructura \parencite{rahman2019}.

Un montaje que cumpla sólo la primera propiedad es un repositorio de guiones,
no un sistema de IaC. La cuarta es la que convierte la automatización en
control, y es también la que con más frecuencia falta en las implantaciones
reales.

\subsec{03. La frontera del sistema}

Delimitar el alcance exige decir también qué queda fuera. El
Cuadro~\ref{tab:frontera} fija la frontera adoptada en este documento.

\begin{table}[htbp]
  \centering
  \footnotesize
  \caption{Frontera del sistema de infraestructura como código.}
  \label{tab:frontera}
  \begin{tabularx}{\textwidth}{@{}>{\raggedright\arraybackslash}X
    >{\raggedright\arraybackslash}X@{}}
    \toprule
    \textbf{Dentro del alcance} & \textbf{Fuera del alcance} \\
    \midrule
    Cómputo, red, almacenamiento, identidad y servicios gestionados
      declarados como código. &
    El sustrato físico del proveedor: centros de datos, cableado,
      refrigeración. \\
    Configuración del sistema operativo y de los agentes instalados en él. &
    El código de la aplicación de negocio y sus pruebas funcionales. \\
    Construcción de imágenes y artefactos inmutables. &
    Los datos de producción y su ciclo de vida, que exigen migraciones y
      copias de seguridad con reglas propias. \\
    Orquestación de cargas y despliegue continuo de manifiestos. &
    La gestión de incidencias, la guardia y los procedimientos de crisis. \\
    Estado, secretos, políticas, pruebas y observabilidad de la propia
      infraestructura. &
    La observabilidad funcional de la aplicación (trazas de negocio,
      experiencia de usuario). \\
    \bottomrule
  \end{tabularx}
\end{table}

La distinción entre infraestructura y datos merece un comentario, porque es la
frontera que más problemas causa en la práctica. Un servidor de base de datos
está dentro del alcance: su tamaño, su versión, su red y sus copias
programadas se declaran como código. El contenido de esa base de datos no lo
está: no es idempotente, no se puede destruir y recrear, y su presencia obliga
a que el sistema de IaC distinga recursos desechables de recursos con estado
persistente, y proteja a los segundos de las operaciones de reemplazo.

% ==========================================================================
\section{El principio de funcionamiento}
\label{sec:principio}
% ==========================================================================

\subsec{01. Tres estados y una función}

Todo el funcionamiento de un sistema de IaC se sostiene sobre tres
representaciones de la misma infraestructura y sobre la función que las
compara.

\begin{description}[style=nextline, leftmargin=0pt, labelwidth=0pt, labelsep=0pt]
  \item[Estado deseado.] Lo que la especificación declara. Es un artefacto de
    texto, versionado, que describe recursos y sus propiedades. No describe
    los pasos para llegar a ellos.
  \item[Estado registrado.] Lo que el motor cree que existe, resultado de la
    última operación que él mismo ejecutó. Se materializa en un fichero o
    servicio de estado que asocia cada recurso declarado con el identificador
    real que el proveedor le asignó \parencite{terraformstate}.
  \item[Estado real.] Lo que efectivamente existe en el proveedor, que sólo se
    conoce consultando su interfaz de programación.
\end{description}

La función central del sistema es la \textbf{conciliación}: calcular la
diferencia entre el estado deseado y el estado real, y emitir el conjunto
mínimo de operaciones que anula esa diferencia. La Figura~\ref{fig:bucle}
representa el bucle.

\begin{figure}[ht]
  \centering
  \begin{tikzpicture}
    \node[bloque, fill=blue!8, text width=30mm, minimum height=12mm] (des)
      {\textbf{Estado deseado}\\\scriptsize especificación versionada};
    \node[bloque, fill=orange!12, text width=30mm, minimum height=12mm,
      right=26mm of des] (comp)
      {\textbf{Comparador}\\\scriptsize cálculo de la diferencia};
    \node[bloque, fill=green!10, text width=30mm, minimum height=12mm,
      right=26mm of comp] (real)
      {\textbf{Estado real}\\\scriptsize recursos del proveedor};
    \node[bloque, fill=gray!10, text width=30mm, minimum height=12mm,
      below=17mm of comp] (act)
      {\textbf{Actuador}\\\scriptsize crear, modificar, destruir};

    \draw[flecha] (des) -- node[above, etiq] {declaración} (comp);
    \draw[flecha] (real) -- node[above, etiq] {observación} (comp);
    \draw[flecha] (comp) -- node[right, etiq, xshift=1mm] {plan} (act);
    \draw[flecha] (act.east) to[out=0, in=-90]
      node[pos=0.55, below right, etiq, xshift=-3mm, yshift=1mm]
      {\parbox{24mm}{\centering operaciones\\sobre la API}} (real.south);
  \end{tikzpicture}
  \caption{Bucle de conciliación. El sistema no ejecuta una secuencia de
  pasos, sino que cierra repetidamente la distancia entre dos
  representaciones. Si la diferencia es vacía, el actuador no ejecuta nada, y
  ésa es exactamente la propiedad de idempotencia. Elaboración propia.}
  \label{fig:bucle}
\end{figure}

\subsec{02. Declarativo frente a imperativo}

La diferencia entre los dos paradigmas no es de sintaxis sino de quién calcula
el camino. En el enfoque imperativo, el autor escribe la secuencia de
operaciones y asume la responsabilidad de que sea correcta desde cualquier
estado de partida. En el enfoque declarativo, el autor escribe el destino y el
motor deriva la secuencia.

La consecuencia práctica es que un guion imperativo sólo es correcto para el
estado de partida que su autor imaginó, mientras que una especificación
declarativa es correcta para cualquiera, a condición de que el motor sepa
observar el estado real. Por eso el paradigma declarativo domina la capa de
aprovisionamiento, donde el espacio de estados de partida es enorme, y
sobrevive el imperativo en tareas de configuración muy locales, donde ese
espacio es pequeño y conocido.

Casi ningún sistema es puro. Las herramientas de configuración expresan los
pasos en orden, pero cada paso es un módulo idempotente que sólo actúa si hace
falta \parencite{ansibledocs}. Es un híbrido: secuencia imperativa de
operaciones declarativas.

\subsec{03. Convergencia y deriva}

Un sistema convergente es el que, aplicado repetidamente, se aproxima al
estado deseado y luego se queda quieto. La \textbf{deriva}
(\emph{drift}) es la aparición de diferencias entre el estado real y el
deseado por causas ajenas al sistema: un cambio manual de urgencia, una acción
automática del proveedor, la caducidad de un certificado, el borrado
accidental de un recurso.

La deriva es inevitable, y la madurez de una implantación se mide por lo que
hace con ella. Hay tres posturas, en orden creciente de rigor:

\begin{enumerate}
  \item \textbf{Ignorarla.} El motor sólo se ejecuta cuando alguien lo lanza.
    Entre ejecuciones, el estado real puede alejarse arbitrariamente.
  \item \textbf{Detectarla.} Una ejecución periódica en modo de sólo lectura
    compara y notifica, sin corregir. Convierte la deriva en una métrica
    observable.
  \item \textbf{Corregirla.} Un agente concilia de forma continua y revierte
    todo cambio no declarado. Es el principio de conciliación continua de
    GitOps \parencite{opengitops} y el modo nativo de los controladores de
    Kubernetes \parencite{k8scontrollers,burns2016}.
\end{enumerate}

% ==========================================================================
\section{El ciclo básico de operación}
\label{sec:ciclo}
% ==========================================================================

Descendiendo un nivel desde el bucle abstracto, la operación cotidiana de un
sistema completo recorre siete pasos. La Figura~\ref{fig:ciclo} los ordena.

\begin{figure}[ht]
  \centering
  \begin{tikzpicture}[node distance=4.5mm and 5.5mm]
    \tikzset{paso/.style={bloque, fill=capa!12, draw=capa, text width=20mm,
      minimum height=13mm, font=\scriptsize}}
    \node[paso] (p1) {\textbf{1. Codificar}\\ escribir la\\especificación};
    \node[paso, right=of p1] (p2) {\textbf{2. Validar}\\ formato, tipos,\\
      políticas};
    \node[paso, right=of p2] (p3) {\textbf{3. Planificar}\\ diferencia\\
      deseado--real};
    \node[paso, right=of p3] (p4) {\textbf{4. Autorizar}\\ revisión y\\
      aprobación};
    \node[paso, right=of p4] (p5) {\textbf{5. Aplicar}\\ ejecutar sobre\\
      la API};
    \node[paso, below=9mm of p4] (p6) {\textbf{6. Verificar}\\ pruebas de\\
      infraestructura};
    \node[paso, below=9mm of p2] (p7) {\textbf{7. Conciliar}\\ vigilar y\\
      corregir la deriva};

    \foreach \a/\b in {p1/p2,p2/p3,p3/p4,p4/p5}
      \draw[flecha, draw=capa] (\a) -- (\b);
    \draw[flecha, draw=capa] (p5.south) |- (p6.east);
    \draw[flecha, draw=capa] (p6) -- (p7);
    \draw[flecha, draw=acento, dashed] (p7.west) -|
      node[pos=0.28, below, etiq, text=acento, yshift=-1mm]
      {deriva detectada} (p1.south);
  \end{tikzpicture}
  \caption{Los siete pasos del ciclo de operación. El retorno discontinuo es
  el que distingue un sistema de IaC de una simple automatización de
  despliegue. Elaboración propia.}
  \label{fig:ciclo}
\end{figure}

Los pasos 2, 4 y 6 son los que suelen faltar en implantaciones incompletas, y
son precisamente los que aportan garantía. Sin validación de políticas, el
sistema automatiza también los errores de configuración; sin autorización,
elimina el punto de control humano sobre cambios irreversibles; sin
verificación, no distingue entre \enquote{el motor terminó sin error} y
\enquote{la infraestructura funciona}.

\subsec{01. El paso crítico: la comparación a tres bandas}

El paso 3 merece detalle porque en él se concentra el mecanismo que hace útil
al sistema. El motor no compara dos cosas sino tres, y lo hace en dos fases
(Figura~\ref{fig:tresbandas}).

\begin{figure}[ht]
  \centering
  \begin{tikzpicture}
    \tikzset{caja/.style={bloque, text width=34mm, minimum height=15mm}}
    \node[caja, fill=blue!8]  (d) {\textbf{Deseado}\\\scriptsize lo que
      declara el código};
    \node[caja, fill=gray!12, right=20mm of d] (r) {\textbf{Registrado}\\
      \scriptsize lo que el motor anotó\\ en la última ejecución};
    \node[caja, fill=green!10, right=20mm of r] (a) {\textbf{Real}\\
      \scriptsize lo que devuelve\\ la API del proveedor};

    \draw[efecto, draw=trans] (r) -- node[above, etiq, text=trans]
      {\textbf{fase 1}} node[below, etiq, text=trans] {refresco} (a);
    \draw[efecto, draw=acento] (d) -- node[above, etiq, text=acento]
      {\textbf{fase 2}} node[below, etiq, text=acento] {plan} (r);

    \node[bloque, fill=acento!10, draw=acento, below=12mm of r,
      text width=0.62\textwidth]
      {\textbf{Plan de cambios}: lista de recursos a crear, modificar,
       reemplazar o destruir, con el valor anterior y el posterior de cada
       propiedad afectada};
    \draw[flecha, draw=acento] ($(d.south)!0.5!(a.south)+(0,-2mm)$) --
      ++(0,-5mm);
  \end{tikzpicture}
  \caption{Comparación a tres bandas. La fase de refresco actualiza el
  registro contra el mundo real y es la que descubre la deriva; la fase de
  plan compara el código contra ese registro ya actualizado y es la que
  produce las operaciones. Elaboración propia.}
  \label{fig:tresbandas}
\end{figure}

El estado registrado no es redundante. Cumple tres funciones que ninguna de
las otras dos representaciones puede cumplir. Primera, traduce el nombre
lógico que el código da a un recurso al identificador opaco que el proveedor
le asignó, sin lo cual el motor no podría reconocer que un recurso existente
le pertenece. Segunda, permite detectar destrucciones: un recurso presente en
el registro y ausente del proveedor sólo puede detectarse comparando contra el
registro, nunca comparando el código con la realidad. Tercera, guarda
propiedades calculadas por el proveedor que el código no declara y que otros
recursos necesitan.

Ese fichero es, en consecuencia, un activo crítico: contiene identificadores y
con frecuencia valores sensibles, y su corrupción o pérdida deja al sistema
sin capacidad de reconocer lo que él mismo creó. De ahí que todo montaje serio
lo almacene en un servicio remoto con cifrado, versionado y bloqueo de
concurrencia.

% ==========================================================================
\section{Las capas de un sistema completo}
\label{sec:capas}
% ==========================================================================

Un sistema completo se organiza en siete capas apiladas, cada una de las
cuales consume el resultado de la anterior y expone una abstracción más alta.
La Figura~\ref{fig:pila} las presenta con sus herramientas representativas.

\begin{figure}[ht]
  \centering
  \begin{tikzpicture}[y=13.6mm]
    \def\W{14.6}
    \foreach \i/\nom/\fun/\tool in {
      0/{L0 · Sustrato}/{cómputo, red, almacenamiento e identidad del
         proveedor}/{API de AWS, Azure o GCP; equipo físico},
      1/{L1 · Aprovisionamiento}/{crea, modifica y destruye los recursos
         declarados}/{Terraform, OpenTofu, Pulumi, CloudFormation, Bicep},
      2/{L2 · Imágenes}/{construye artefactos inmutables ya
         configurados}/{Packer, Dockerfile, \emph{buildpacks}},
      3/{L3 · Configuración}/{ajusta el interior de la máquina ya
         arrancada}/{Ansible, Chef, Puppet, cloud-init},
      4/{L4 · Orquestación}/{programa cargas sobre el conjunto de
         máquinas}/{Kubernetes, Nomad, Helm, Kustomize},
      5/{L5 · Entrega}/{lleva el código al motor y concilia de forma
         continua}/{Argo CD, Flux, canalizaciones de integración},
      6/{L6 · Plataforma}/{expone catálogo y autoservicio al equipo de
         producto}/{Backstage, Crossplane, portales internos}} {
      \node[draw=capa, very thick, fill=capa!\the\numexpr 6+2*\i\relax,
        rounded corners=2pt, minimum width=\W cm, minimum height=12.6mm,
        anchor=south west] (L\i) at (0,\i) {};
      \node[anchor=west, font=\footnotesize\bfseries, text=capa!60!black,
        text width=37mm, align=left] at ([xshift=3.5mm]L\i.west) {\nom};
      \node[anchor=west, font=\scriptsize, text width=50mm, align=left]
        at ([xshift=43mm]L\i.west) {\fun};
      \node[anchor=west, font=\scriptsize\itshape, text width=45mm,
        align=left] at ([xshift=96mm]L\i.west) {\tool};
    }
    \draw[flecha, draw=capa, line width=1pt] (-0.5,0.2) -- (-0.5,6.8)
      node[midway, rotate=90, above, font=\scriptsize, text=capa]
      {abstracción creciente};
  \end{tikzpicture}
  \caption{Las siete capas de un sistema completo de infraestructura como
  código. Cada capa consume la salida de la inferior. Las herramientas citadas
  son representativas, no exhaustivas. Elaboración propia.}
  \label{fig:pila}
\end{figure}

\subsec{01. Qué hace cada capa}

\textbf{L0, sustrato.} No es parte del sistema de IaC sino su objeto. Aporta
la interfaz de programación sobre la que todo lo demás actúa. Su
característica determinante es que ofrece recursos con identidad propia,
creables y destruibles por llamada, lo que es la condición de posibilidad de
todo el resto.

\textbf{L1, aprovisionamiento.} Traduce la especificación en llamadas de
creación, modificación y destrucción. Es la capa que posee el registro de
estado y la que ejecuta la comparación a tres bandas. Trabaja con recursos que
todavía no tienen sistema operativo arrancado, o cuyo interior no le
concierne. Conviene señalar que en 2023 el cambio de licencia de Terraform dio
origen a OpenTofu, una bifurcación acogida por la Linux Foundation
\parencite{opentofu}, de modo que hoy la capa cuenta con dos implementaciones
compatibles del mismo lenguaje: una decisión de dependencia que una hoja de
ruta tecnológica debería registrar de forma explícita.

\textbf{L2, imágenes.} Desplaza el trabajo de configuración desde el momento
de ejecución al momento de construcción. En lugar de arrancar una máquina
genérica y modificarla, construye una imagen ya completa y arranca instancias
idénticas de ella. Es la base de la infraestructura inmutable: en lugar de
modificar un servidor, se reemplaza. Reduce drásticamente la deriva, porque
elimina la ventana temporal en la que un servidor puede divergir.

\textbf{L3, configuración.} Actúa sobre el interior de máquinas ya arrancadas:
paquetes, ficheros, servicios, usuarios. Es la capa históricamente más antigua
y la que más terreno ha cedido a L2 y L4. Sigue siendo necesaria donde hay
máquinas de larga vida, sistemas heredados o equipos físicos.

\textbf{L4, orquestación.} Introduce una indirección importante: deja de
declararse \enquote{esta carga en esta máquina} y pasa a declararse
\enquote{esta carga, tantas réplicas} sobre un conjunto de máquinas
intercambiables. El planificador decide la colocación. Sus controladores son
bucles de conciliación como el de la Figura~\ref{fig:bucle}, pero embebidos y
permanentes \parencite{k8scontrollers}.

\textbf{L5, entrega.} Conecta el repositorio con el motor. Ejecuta las
canalizaciones que validan, planifican y aplican, y, en el modelo de
conciliación continua, aloja el agente que vigila la deriva sin intervención
humana \parencite{opengitops,humble2010}.

\textbf{L6, plataforma.} Envuelve todo lo anterior en una interfaz de
autoservicio, para que un equipo de producto pueda solicitar una base de datos
sin conocer ni la capa L1 ni las políticas de la organización, que quedan
codificadas en el catálogo.

\subsec{02. La tensión entre capas}

Las capas L2, L3 y L4 compiten parcialmente por el mismo territorio. Un mismo
requisito, por ejemplo \enquote{que este servidor tenga la versión 1.24 de un
tiempo de ejecución}, puede resolverse construyendo una imagen que ya la
incluya (L2), instalándola con un módulo de configuración (L3) o empaquetando
la aplicación en un contenedor que la lleve dentro (L4). Las tres soluciones
funcionan y no deben coexistir sobre el mismo recurso: si dos capas gobiernan
la misma propiedad, cada una revierte el trabajo de la otra en cada ciclo de
conciliación, y el sistema oscila.

Decidir qué capa posee cada propiedad es, por tanto, la primera decisión de
arquitectura de una implantación, y conviene documentarla de forma explícita.

% ==========================================================================
\section{Los módulos transversales}
\label{sec:transversales}
% ==========================================================================

Las siete capas describen el flujo principal, pero un sistema completo
necesita además seis módulos que atraviesan todas las capas
(Figura~\ref{fig:transversales}). Ninguno produce infraestructura; todos
condicionan si la que se produce es correcta, segura y auditable.

\begin{figure}[ht]
  \centering
  \begin{tikzpicture}
    \node[draw=capa, very thick, fill=capa!10, rounded corners=3pt,
      minimum width=5.4cm, minimum height=3.7cm, align=center,
      font=\footnotesize\bfseries, text=capa!60!black] (nucleo)
      {Las siete capas\\del flujo principal\\[2pt]
       \scriptsize\mdseries L0 sustrato · L1 aprovisionamiento\\
       \scriptsize\mdseries L2 imágenes · L3 configuración\\
       \scriptsize\mdseries L4 orquestación · L5 entrega\\
       \scriptsize\mdseries L6 plataforma};
    \foreach \ang/\nom/\det in {
      90/{Control de versiones}/{historia, revisión, reversión},
      32/{Gestión de estado}/{registro, bloqueo, cifrado},
      -32/{Gestión de secretos}/{inyección en tiempo de ejecución},
      -90/{Política como código}/{reglas evaluables antes de aplicar},
      -148/{Pruebas}/{estáticas, unitarias, de integración},
      148/{Observabilidad}/{deriva, coste, cumplimiento}} {
      \node[bloque, fill=trans!12, draw=trans, text width=31mm,
        font=\scriptsize, align=center, anchor=center] at (\ang:5.1cm)
        (m\ang) {\textbf{\nom}\\ \det};
      \draw[efecto, draw=trans, shorten >=3pt, shorten <=3pt]
        (m\ang) -- (nucleo);
    }
  \end{tikzpicture}
  \caption{Los seis módulos transversales. Cada uno interactúa con todas las
  capas del flujo principal, no con una en particular. Elaboración propia.}
  \label{fig:transversales}
\end{figure}

El Cuadro~\ref{tab:transversales} resume la función de cada uno y el fallo
característico que aparece cuando falta.

\begin{table}[htbp]
  \centering
  \footnotesize
  \caption{Módulos transversales: función, instrumentos y fallo característico
  por ausencia.}
  \label{tab:transversales}
  \begin{tabularx}{\textwidth}{@{}>{\raggedright\arraybackslash}p{0.145\textwidth}
    >{\raggedright\arraybackslash}X
    >{\raggedright\arraybackslash}p{0.30\textwidth}@{}}
    \toprule
    \textbf{Módulo} & \textbf{Función} & \textbf{Fallo por ausencia} \\
    \midrule
    Control de versiones & Registrar quién cambió qué, cuándo y por qué;
      permitir revertir a un estado anterior conocido. &
      Cambios sin trazabilidad y reversión imposible. \\
    Gestión de estado & Custodiar el registro: almacenamiento remoto, cifrado,
      versionado y bloqueo para impedir ejecuciones concurrentes. &
      Corrupción del registro y recursos huérfanos. \\
    Gestión de secretos & Mantener credenciales fuera del código e inyectarlas
      en el momento de la ejecución, con rotación y ámbito mínimo. Es uno de
      los defectos más frecuentes en código de infraestructura
      \parencite{rahman2019smells}. &
      Credenciales escritas en el repositorio. \\
    Política como código & Expresar las reglas de la organización como
      predicados evaluables sobre el plan, antes de aplicarlo
      \parencite{opa}. &
      El sistema automatiza también los incumplimientos, y a mayor
      velocidad. \\
    Pruebas & Verificar la especificación en varios niveles: análisis
      estático, validación del plan, despliegue efímero y comprobación
      funcional. &
      \enquote{Terminó sin error} se confunde con \enquote{funciona}. \\
    Observabilidad & Medir deriva, coste, tiempo de ciclo y cumplimiento, y
      exponerlos como series temporales. &
      No hay figuras de mérito y la mejora no es demostrable. \\
    \bottomrule
  \end{tabularx}
\end{table}

% ==========================================================================
\section{Modos de operación}
\label{sec:modos}
% ==========================================================================

Sobre la misma arquitectura caben decisiones de operación que cambian de forma
apreciable el comportamiento del sistema. Las tres principales son
independientes entre sí.

\subsec{01. Empuje frente a arrastre}

En el modo de \textbf{empuje} (\emph{push}), un actor externo, típicamente una
canalización de integración, posee las credenciales del entorno y ejecuta el
motor contra él. En el modo de \textbf{arrastre} (\emph{pull}), un agente
residente dentro del entorno consulta el repositorio y se aplica a sí mismo
los cambios. La Figura~\ref{fig:pushpull} contrasta ambos.

\begin{figure}[ht]
  \centering
  \begin{tikzpicture}
    \tikzset{c/.style={bloque, text width=28mm, minimum height=9mm},
             cin/.style={bloque, fill=white, text width=30mm,
               minimum height=7mm, font=\scriptsize},
             ent/.style={draw, thick, rounded corners=2pt, fill=green!10,
               minimum width=52mm}}
    % ---------------- empuje ----------------
    \node[etiq, font=\footnotesize\bfseries] (t1) {Empuje (\emph{push})};
    \node[c, fill=blue!8, below=5mm of t1] (r1) {Repositorio};
    \node[c, fill=orange!12, below=11mm of r1] (ci1)
      {Canalización\\ \scriptsize con las credenciales};
    \node[ent, minimum height=24mm, below=11mm of ci1] (e1) {};
    \node[anchor=north, font=\scriptsize\bfseries]
      at ([yshift=-2mm]e1.north) {Entorno};
    \node[cin, anchor=south] at ([yshift=3mm]e1.south) (i1)
      {Infraestructura};
    \draw[flecha] (r1) -- node[right, etiq, xshift=1mm] {disparo} (ci1);
    \draw[flecha] (ci1) -- node[right, etiq, xshift=1mm] {aplica} (e1.north);
    % ---------------- arrastre ----------------
    \node[etiq, font=\footnotesize\bfseries, right=52mm of t1] (t2)
      {Arrastre (\emph{pull})};
    \node[c, fill=blue!8, below=5mm of t2] (r2) {Repositorio};
    \node[ent, minimum height=34mm, below=26mm of r2] (e2) {};
    \node[anchor=north, font=\scriptsize\bfseries]
      at ([yshift=-1.5mm]e2.north) {Entorno};
    \node[cin, fill=orange!12, anchor=north] at ([yshift=-7mm]e2.north) (ag)
      {Agente\\ \scriptsize con las credenciales};
    \node[cin, anchor=south] at ([yshift=3mm]e2.south) (i2)
      {Infraestructura};
    \draw[flecha] ([xshift=-11mm]e2.north) to[out=95,in=-95]
      node[left, etiq, xshift=-1mm] {consulta} ([xshift=-9mm]r2.south);
    \draw[flecha] ([xshift=9mm]r2.south) to[out=-85,in=85]
      node[right, etiq, xshift=1mm] {\parbox{18mm}{\centering estado\\
      deseado}} ([xshift=11mm]e2.north);
    \draw[flecha, draw=acento] (ag) -- node[right, font=\tiny\itshape,
      text=acento, xshift=0.5mm] {concilia} (i2);
  \end{tikzpicture}
  \caption{Empuje frente a arrastre. En el arrastre, las credenciales del
  entorno nunca salen de él, lo que reduce la superficie de ataque y hace del
  agente el punto natural de conciliación continua. Elaboración propia a
  partir de los principios de \textcite{opengitops}.}
  \label{fig:pushpull}
\end{figure}

El arrastre tiene dos ventajas estructurales. La primera es de seguridad: las
credenciales de producción no viajan a un sistema de integración compartido.
La segunda es que el agente, por estar siempre presente, puede conciliar de
forma continua y no sólo cuando alguien confirma un cambio, que es la cuarta
propiedad exigida en la sección~\ref{sec:alcance}. Su coste es que exige un agente por
entorno y complica los flujos que necesitan orden entre entornos.

\subsec{02. Mutable frente a inmutable}

En el modelo mutable, un recurso desplegado se modifica en su sitio. En el
inmutable, no se modifica nunca: se construye uno nuevo con la configuración
deseada y se sustituye. El modelo inmutable elimina por construcción la deriva
de configuración y hace que el recurso desplegado sea siempre trazable a un
artefacto concreto. A cambio, obliga a que todo estado persistente viva fuera
del recurso reemplazable, lo que es una restricción de diseño exigente.

\subsec{03. El cuadro de decisión}

Las tres decisiones son independientes entre sí, de modo que caben ocho
combinaciones. La más frecuente en implantaciones recientes es la de arrastre,
inmutable y declarativa, pero ninguna de las tres opciones es obligada. El
Cuadro~\ref{tab:modos} resume las consecuencias de cada una.

\begin{table}[!ht]
  \centering
  \footnotesize
  \caption{Tres decisiones de operación y sus consecuencias.}
  \label{tab:modos}
  \begin{tabularx}{\textwidth}{@{}>{\raggedright\arraybackslash}p{0.15\textwidth}
    >{\raggedright\arraybackslash}X>{\raggedright\arraybackslash}X@{}}
    \toprule
    \textbf{Decisión} & \textbf{Opción A} & \textbf{Opción B} \\
    \midrule
    Iniciativa &
      \textbf{Empuje}: control de orden entre entornos; credenciales fuera del
      entorno; conciliación sólo bajo disparo. &
      \textbf{Arrastre}: credenciales confinadas; conciliación continua;
      un agente por entorno. \\
    Mutabilidad &
      \textbf{Mutable}: cambios rápidos y baratos; deriva acumulativa; difícil
      trazar qué versión corre. &
      \textbf{Inmutable}: sin deriva; reemplazo trazable; obliga a externalizar
      todo estado persistente. \\
    Expresión &
      \textbf{Declarativa}: el motor calcula el camino; correcta desde
      cualquier estado de partida. &
      \textbf{Imperativa}: control fino del orden; sólo correcta desde el
      estado de partida previsto. \\
    \bottomrule
  \end{tabularx}
\end{table}

% ==========================================================================
\section{El módulo como unidad de composición}
\label{sec:modulo}
% ==========================================================================

El término \enquote{módulo} tiene en IaC un sentido preciso: es la unidad
reutilizable que encapsula un conjunto de recursos tras una interfaz de
entradas y salidas. Es el mecanismo que permite que una organización no repita
la misma configuración en cada equipo y, sobre todo, que las decisiones
correctas queden incorporadas por omisión.

\begin{figure}[ht]
  \centering
  \begin{tikzpicture}
    \node[draw=capa, very thick, fill=capa!8, rounded corners=3pt,
      minimum width=5.6cm, minimum height=3.2cm] (mod) {};
    \node[anchor=north, font=\footnotesize\bfseries, text=capa!60!black]
      at ([yshift=-3mm]mod.north) {Módulo \texttt{red-privada v2.3.1}};
    \node[bloque, fill=white, font=\scriptsize, text width=44mm,
      anchor=north] at ([yshift=-9mm]mod.north)
      {\textbf{Cuerpo}: recursos declarados,\\ valores por omisión,
       validaciones,\\ etiquetas obligatorias};

    \node[bloque, fill=blue!8, text width=27mm, font=\scriptsize,
      left=16mm of mod] (in)
      {\textbf{Entradas}\\ región, CIDR,\\ número de zonas,\\ etiquetas};
    \node[bloque, fill=green!10, text width=27mm, font=\scriptsize,
      right=16mm of mod] (out)
      {\textbf{Salidas}\\ identificador de red,\\ subredes,\\ tabla de rutas};
    \node[bloque, fill=trans!12, text width=40mm, font=\scriptsize,
      below=9mm of mod] (reg)
      {\textbf{Registro versionado}\\ publica, firma y resuelve versiones};

    \draw[flecha] (in) -- (mod);
    \draw[flecha] (mod) -- (out);
    \draw[flecha] (reg) -- (mod);

    \node[bloque, fill=gray!8, text width=25mm, font=\scriptsize,
      above left=8mm and 8mm of mod.north] (c1) {Composición\\
      \emph{entorno de pruebas}};
    \node[bloque, fill=gray!8, text width=25mm, font=\scriptsize,
      above right=8mm and 8mm of mod.north] (c2) {Composición\\
      \emph{entorno de producción}};
    \draw[flecha, gray] (c1) -- (mod.north west);
    \draw[flecha, gray] (c2) -- (mod.north east);
  \end{tikzpicture}
  \caption{Anatomía de un módulo. Las mismas entradas con valores distintos
  producen entornos comparables; el registro versionado permite que una
  corrección se propague de forma controlada. Elaboración propia.}
  \label{fig:modulo}
\end{figure}

Tres propiedades hacen que un módulo cumpla su función. Debe tener una
\textbf{interfaz estrecha}: cuantas menos entradas expone, más decisiones fija
y más valor aporta. Debe estar \textbf{versionado} con semántica explícita,
para que quien lo consume elija cuándo adoptar un cambio. Y debe delimitar un
\textbf{radio de explosión} razonable: un módulo que agrupa demasiados
recursos obliga a que cada cambio pequeño planifique sobre un conjunto grande,
lo que alarga el ciclo y aumenta el riesgo de cada aplicación.

% ==========================================================================
\section{Matriz de estructura de dependencia (DSM)}
\label{sec:dsm}
% ==========================================================================

Esta sección corresponde al elemento 2 del índice de una hoja de ruta
tecnológica. Las secciones anteriores han enumerado las piezas del sistema,
pero enumerarlas no basta: lo que decide el orden de una implantación, el
riesgo de cada paso y la conveniencia de agrupar unos trabajos y separar otros
son las dependencias entre ellas. La matriz de estructura de dependencia, o
DSM \parencite{steward1981,eppinger2012}, es el instrumento que las hace
visibles y operables, porque permite reordenar filas y columnas hasta que la
propia matriz muestra la secuencia.

Se construyen dos matrices, porque la hoja de ruta necesita responder a dos
preguntas distintas. La primera es de arquitectura: qué componente del sistema
necesita a qué otro, lo que determina qué se puede diseñar y desplegar por
separado. La segunda es de cartera: qué proyecto de la hoja de ruta necesita
que otro esté terminado, lo que determina el calendario y el orden de
inversión.

\subsec{01. Convención de lectura}

Ambas matrices son cuadradas y llevan los mismos elementos en filas y en
columnas, en el mismo orden. Se adopta la convención de entradas por fila:

\begin{quote}
Una marca en la celda de la fila $i$ y la columna $j$ significa que el
elemento $i$ \textbf{necesita algo del} elemento $j$. Se lee, por tanto,
recorriendo la fila $i$ de izquierda a derecha: las marcas de esa fila son
todo lo que $i$ recibe. Recorriendo en cambio la columna $j$ de arriba abajo
se obtiene todo lo que $j$ entrega.
\end{quote}

De esa convención se derivan las tres lecturas que se usan después. La
\textbf{suma de una fila} mide cuántas condiciones previas tiene un elemento,
es decir, su dificultad de arranque. La \textbf{suma de una columna} mide a
cuántos elementos afecta, es decir, el alcance del daño si falla o se retrasa.
Y una marca \textbf{por encima de la diagonal}, una vez ordenada la matriz,
señala una realimentación: un elemento que necesita algo que en el orden
propuesto todavía no existe, lo que obliga a iterar.

La diagonal no se marca: ningún elemento se necesita a sí mismo. Las celdas
diagonales se sombrean únicamente para orientar la vista.

\subsec{02. La DSM de arquitectura}

Los elementos son los trece componentes ya descritos: las siete capas de la
Figura~\ref{fig:pila} y los seis módulos transversales del
Cuadro~\ref{tab:transversales}. El Cuadro~\ref{tab:dsmelem} fija los códigos
que usa la matriz y, para cada componente, el tipo de dependencia que
introduce.

\begin{table}[htbp]
  \centering
  \footnotesize
  \caption{Los trece componentes de la DSM de arquitectura y la dependencia
  que cada uno introduce en el sistema.}
  \label{tab:dsmelem}
  \begin{tabularx}{\textwidth}{@{}l>{\raggedright\arraybackslash}p{0.235\textwidth}
    >{\raggedright\arraybackslash}X@{}}
    \toprule
    \textbf{Cód.} & \textbf{Componente} & \textbf{Qué necesita de los demás} \\
    \midrule
    L0 & Sustrato del proveedor & Nada del sistema: es su objeto. \\
    L1 & Aprovisionamiento & La interfaz de L0, la especificación versionada,
         el registro de estado, las credenciales y el veredicto de la
         política. \\
    L2 & Imágenes & Sustrato donde construir, definiciones versionadas y
         credenciales del registro de artefactos. \\
    L3 & Configuración & Máquinas ya creadas por L1 a partir de imágenes de
         L2. \\
    L4 & Orquestación & El agrupamiento aprovisionado por L1 y las imágenes de
         contenedor de L2. \\
    L5 & Entrega & El motor de L1, el destino de L4, el repositorio, las
         políticas y las pruebas que ejecuta en cada paso. \\
    L6 & Plataforma & El catálogo que ejecuta sobre L1, la canalización de L5,
         las políticas que codifica y las métricas que muestra. \\
    M1 & Control de versiones & Nada: es la raíz del sistema. \\
    M2 & Gestión de estado & Almacenamiento remoto en L0 y claves de cifrado
         del módulo de secretos. \\
    M3 & Gestión de secretos & El servicio de custodia y cifrado de L0. \\
    M4 & Política como código & El plan que produce L1 y el repositorio donde
         viven las reglas. \\
    M5 & Pruebas & Entornos efímeros que crea L1 y el repositorio del que
         parten. \\
    M6 & Observabilidad & Telemetría de L0, ejecuciones de sólo lectura de L1,
         registros de L5 y el estado custodiado por M2. \\
    \bottomrule
  \end{tabularx}
\end{table}

La Figura~\ref{fig:dsmarq} presenta la matriz en el orden natural de
exposición, capas primero y módulos después. Es el orden en que se ha descrito
el sistema, no el orden en que conviene construirlo, y por eso las marcas
aparecen repartidas a ambos lados de la diagonal.

\begin{figure}[ht]
  \centering
  \begin{tikzpicture}[x=6.4mm, y=6.4mm, font=\scriptsize]
    \def\n{13}
    % ---- marcas: fila/columna (la fila necesita algo de la columna)
    \def\marcas{2/1,2/8,2/9,2/10,2/11,
                3/1,3/8,3/10,
                4/2,4/3,4/8,4/10,
                5/2,5/3,5/8,5/10,
                6/2,6/5,6/8,6/10,6/11,6/12,
                7/2,7/6,7/8,7/11,7/13,
                9/1,9/10,
                10/1,
                11/2,11/8,
                12/2,12/8,
                13/1,13/2,13/6,13/9}
    % ---- fondo de la rejilla y diagonal
    \fill[gray!4] (0,0) rectangle (\n,-\n);
    \foreach \i in {1,...,13}
      \fill[gray!35] (\i-1,-\i+1) rectangle (\i,-\i);
    % ---- marcas
    \foreach \r/\c in \marcas
      \fill[dsmarq] (\c-1+0.14,-\r+1-0.14) rectangle (\c-0.14,-\r+0.14);
    % ---- rejilla
    \draw[gray!55, thin, step=1] (0,0) grid (\n,-\n);
    \draw[thick] (0,0) rectangle (\n,-\n);
    % ---- separación entre capas y módulos
    \draw[thick, gray!70] (7,0) -- (7,-\n);
    \draw[thick, gray!70] (0,-7) -- (\n,-7);
    % ---- etiquetas
    \foreach \i/\lab in {1/L0,2/L1,3/L2,4/L3,5/L4,6/L5,7/L6,
                         8/M1,9/M2,10/M3,11/M4,12/M5,13/M6} {
      \node[anchor=east] at (-0.2,-\i+0.5) {\lab};
      \node[anchor=west, rotate=90] at (\i-0.5,0.2) {\lab};
    }
    % ---- sumas
    \foreach \i/\v in {1/0,2/5,3/3,4/4,5/4,6/6,7/5,8/0,9/2,10/1,11/2,12/2,13/4}
      \node[anchor=west, font=\scriptsize\bfseries] at (\n+0.35,-\i+0.5) {\v};
    \foreach \i/\v in {1/5,2/7,3/2,4/0,5/1,6/2,7/0,8/8,9/2,10/6,11/3,12/1,13/1}
      \node[anchor=north, font=\scriptsize\bfseries] at (\i-0.5,-\n-0.25) {\v};
    \node[anchor=west, font=\scriptsize\itshape, align=left]
      at (\n+0.35,0.7) {entradas};
    \node[anchor=north, font=\scriptsize\itshape] at (\n+1.1,-\n-0.25)
      {salidas};
  \end{tikzpicture}
  \caption{DSM de arquitectura en el orden natural de exposición. Una marca en
  la fila $i$ y la columna $j$ indica que el componente $i$ necesita algo del
  componente $j$. La columna de la derecha suma las entradas de cada fila y la
  banda inferior suma las salidas de cada columna. Elaboración propia.}
  \label{fig:dsmarq}
\end{figure}

\subsec{03. Qué dice la matriz sin reordenarla}

Antes de tocar el orden, las sumas ya dicen tres cosas.

\textbf{Dos de los tres concentradores de riesgo son módulos transversales, no
capas.} Las columnas con más salidas son M1, el control de versiones, con
ocho; L1, el aprovisionamiento, con siete; y M3, la gestión de secretos, con
seis. Es un resultado con consecuencias prácticas: una interrupción del
repositorio detiene ocho componentes y una del custodio de secretos detiene
seis, mientras que la caída de la capa de configuración, L3, con cero salidas,
no detiene a ninguno. La disponibilidad que se exige a cada pieza debería
seguir ese orden, y con frecuencia no lo sigue, porque la atención se
concentra en el motor, que es sólo el segundo de los tres.

\textbf{Las capas superiores son las caras de construir, no las de operar.} Las
filas con más entradas son L5, la entrega, con seis, y L1 y L6 con cinco. Un
proyecto que empiece por la interfaz de autoservicio está aceptando cinco
condiciones previas simultáneas; empezar por ahí es la causa habitual de que
un portal interno quede a medias.

\textbf{Hay dos elementos sin ninguna entrada}, L0 y M1, que son por tanto los
únicos puntos de arranque posibles de cualquier secuencia.

\subsec{04. Particionado: la secuencia de implantación}

Reordenar filas y columnas con la misma permutación no cambia la información
de la matriz, sólo la presentación. El particionado busca la permutación que
lleva el mayor número posible de marcas por debajo de la diagonal, porque una
matriz triangular inferior describe un orden en el que cada elemento se
construye cuando todo lo que necesita ya existe. La
Figura~\ref{fig:dsmarqord} recoge el resultado.

\begin{figure}[ht]
  \centering
  \begin{tikzpicture}[x=6.4mm, y=6.4mm, font=\scriptsize]
    \def\n{13}
    % Orden: M1 L0 M3 M2 L2 L1 M4 M5 L3 L4 L5 M6 L6
    % Posición nueva: 1=M1 2=L0 3=M3 4=M2 5=L2 6=L1 7=M4 8=M5 9=L3 10=L4
    %                 11=L5 12=M6 13=L6
    \def\marcas{3/2, 4/2,4/3, 5/1,5/2,5/3, 6/1,6/2,6/3,6/4, 7/1,7/6,
                8/1,8/6, 9/1,9/3,9/5,9/6, 10/1,10/3,10/5,10/6,
                11/1,11/3,11/6,11/7,11/8,11/10,
                12/2,12/4,12/6,12/11, 13/1,13/6,13/7,13/11,13/12}
    \fill[gray!4] (0,0) rectangle (\n,-\n);
    \foreach \i in {1,...,13}
      \fill[gray!35] (\i-1,-\i+1) rectangle (\i,-\i);
    % ---- marca de realimentación (por encima de la diagonal)
    \fill[acento] (6.14,-5.14) rectangle (6.86,-5.86);
    \foreach \r/\c in \marcas
      \fill[dsmarq] (\c-1+0.14,-\r+1-0.14) rectangle (\c-0.14,-\r+0.14);
    \draw[gray!55, thin, step=1] (0,0) grid (\n,-\n);
    \draw[thick] (0,0) rectangle (\n,-\n);
    % ---- escalonado de las olas de implantación, sobre la rejilla
    \foreach \a/\b in {0/2, 2/3, 3/5, 5/7, 7/10, 10/11, 11/12, 12/13}
      \draw[gray!75!black, line width=1pt] (\a,-\a) -- (\b,-\a) -- (\b,-\b);
    % ---- bloque acoplado L1-M4
    \draw[acento, very thick] (5,-5) rectangle (7,-7);
    \foreach \i/\lab in {1/M1,2/L0,3/M3,4/M2,5/L2,6/L1,7/M4,
                         8/M5,9/L3,10/L4,11/L5,12/M6,13/L6} {
      \node[anchor=east] at (-0.2,-\i+0.5) {\lab};
      \node[anchor=west, rotate=90] at (\i-0.5,0.2) {\lab};
    }
    \node[anchor=west, font=\scriptsize\itshape, text=acento]
      at (7.4,-6) {bloque acoplado};
    \draw[acento, thick] (7.05,-5.5) -- (7.35,-6);
  \end{tikzpicture}
  \caption{DSM de arquitectura particionada. En el orden M1, L0, M3, M2, L2,
  L1, M4, M5, L3, L4, L5, M6, L6 todas las dependencias caen por debajo de la
  diagonal salvo una, la que L1 tiene de M4, marcada en rojo. El escalonado
  gris separa las olas: los elementos de una misma ola no dependen entre sí y
  pueden abordarse en paralelo. Elaboración propia.}
  \label{fig:dsmarqord}
\end{figure}

El resultado es una matriz triangular inferior con una sola excepción, y esa
excepción es informativa. El escalonado define ocho olas, y los elementos que
comparten ola pueden abordarse en paralelo:

\begin{enumerate}
  \item \textbf{M1 y L0}: el repositorio y la cuenta del proveedor. Sin
    entradas, son el arranque obligado.
  \item \textbf{M3}: la gestión de secretos. Aparece la segunda, antes que
    cualquier capa, porque seis componentes la necesitan y ninguno puede
    autenticarse sin ella. Ponerla al final, que es lo habitual cuando se
    trata como un endurecimiento posterior, obliga a rehacer todo lo anterior.
  \item \textbf{M2 y L2}: la custodia del registro de estado y la construcción
    de imágenes, que no dependen la una de la otra.
  \item \textbf{L1 y M4}: el bloque acoplado, del que se habla enseguida.
  \item \textbf{M5, L3 y L4}: las pruebas, la configuración y la orquestación,
    los tres apoyados en el motor ya operativo.
  \item \textbf{L5}: la entrega continua, que necesita el motor, la
    orquestación, las políticas y las pruebas a la vez.
  \item \textbf{M6}: la observabilidad, que sólo puede medir lo que ya
    funciona.
  \item \textbf{L6}: la interfaz de autoservicio, la última porque envuelve a
    todo lo demás.
\end{enumerate}

\textbf{El bloque acoplado y cómo se rompe.} L1 necesita de M4 el veredicto de
la política antes de aplicar un cambio, y M4 necesita de L1 el plan sobre el
que pronunciarse. En una DSM eso es un ciclo, y un ciclo obliga a iterar o a
descomponer. La descomposición correcta consiste en partir L1 en dos
subprocesos, uno que produce el plan y otro que lo aplica, con la evaluación
de política entre ambos. Es exactamente lo que hace el diagrama SD1 del modelo
OPM de la sección~\ref{sec:opm}: al separar \emph{Planificación},
\emph{Validación} y \emph{Aplicación} en subprocesos distintos, el ciclo
desaparece y la secuencia queda lineal. La DSM detecta el acoplamiento y el
modelo OPM propone la partición que lo elimina, lo que ilustra por qué los
elementos 2 y 3 de la hoja de ruta se sostienen mutuamente.

\subsec{05. La DSM de la cartera de proyectos}

La segunda matriz cambia de objeto. Sus elementos ya no son componentes del
sistema sino los ocho proyectos con que la empresa pasa de la plataforma
actual a la plataforma de 2029. Tres de ellos se corresponden con los vectores
tecnológicos del Cuadro~\ref{tab:vectores}, que la sección~\ref{sec:comp}
utiliza para construir el gráfico vectorial; los otros cinco son los
habilitadores sin los cuales aquéllos no se sostienen. El
Cuadro~\ref{tab:dsmproy} los enumera.

\begin{table}[htbp]
  \centering
  \footnotesize
  \caption{Los ocho proyectos de la cartera, su correspondencia con los
  vectores tecnológicos del Cuadro~\ref{tab:vectores} y los componentes de la
  DSM de arquitectura sobre los que actúan. Las filas con la columna
  \emph{vector} en blanco son proyectos habilitadores, que no desplazan por sí
  solos ninguna coordenada del gráfico vectorial.}
  \label{tab:dsmproy}
  \begin{tabularx}{\textwidth}{@{}l>{\raggedright\arraybackslash}p{0.245\textwidth}
    >{\raggedright\arraybackslash}X l l@{}}
    \toprule
    \textbf{Cód.} & \textbf{Proyecto} & \textbf{Alcance} & \textbf{Vector} &
      \textbf{Comp.} \\
    \midrule
    R1 & Cierre de cobertura de codificación & Inventariar los recursos
         creados a mano e importarlos al motor hasta llevar la cobertura del
         55 al 86 por ciento. & & L1 \\
    R2 & Custodia del registro de estado & Estado remoto, cifrado, versionado
         y con bloqueo de concurrencia. & & M2 \\
    R3 & Catálogo de módulos certificados & Módulos versionados, firmados y
         publicados en un registro interno, con interfaz estrecha. & P1 &
         L1, L6 \\
    R4 & Política como código & Reglas de la organización evaluadas sobre el
         plan antes de aplicarlo. & P1 & M4 \\
    R5 & Malla de pruebas y entornos efímeros & Análisis estático, validación
         del plan y despliegue efímero verificado. & & M5 \\
    R6 & Sustitución del motor & Migración a la bifurcación abierta con
         estado gestionado propio. & P2 & L1 \\
    R7 & Interfaz de autoservicio & Portal que expone el catálogo al equipo
         de producto. & P3 & L6 \\
    R8 & Observabilidad de la plataforma & Deriva, coste, tiempo de ciclo y
         cumplimiento como series temporales. & & M6 \\
    \bottomrule
  \end{tabularx}
\end{table}

Las dependencias entre proyectos no son las mismas que entre componentes,
porque un proyecto puede necesitar de otro un resultado parcial y no el
componente entero. La Figura~\ref{fig:dsmproy} distingue por eso dos
intensidades: dependencia \emph{necesaria}, cuando el proyecto no puede
empezar sin el resultado del otro, y dependencia \emph{conveniente}, cuando
puede empezar pero a un coste o un riesgo mayores.

\begin{figure}[ht]
  \centering
  \begin{tikzpicture}[x=8.4mm, y=8.4mm, font=\scriptsize]
    \def\n{8}
    \def\fuerte{1/2, 3/1, 4/3, 5/3, 6/2, 6/3, 6/5, 7/3, 7/4, 8/1}
    \def\debil{3/2, 4/1, 5/2, 7/6, 8/2}
    \fill[gray!4] (0,0) rectangle (\n,-\n);
    \foreach \i in {1,...,8}
      \fill[gray!35] (\i-1,-\i+1) rectangle (\i,-\i);
    \foreach \r/\c in \fuerte
      \fill[dsmpro] (\c-1+0.16,-\r+1-0.16) rectangle (\c-0.16,-\r+0.16);
    \foreach \r/\c in \debil
      \draw[dsmpro, thick, fill=dsmpro!20]
        (\c-1+0.16,-\r+1-0.16) rectangle (\c-0.16,-\r+0.16);
    \draw[gray!55, thin, step=1] (0,0) grid (\n,-\n);
    \draw[thick] (0,0) rectangle (\n,-\n);
    \foreach \i/\lab in {1/R1,2/R2,3/R3,4/R4,5/R5,6/R6,7/R7,8/R8} {
      \node[anchor=east] at (-0.2,-\i+0.5) {\lab};
      \node[anchor=south] at (\i-0.5,0.1) {\lab};
    }
  \end{tikzpicture}
  \caption{DSM de la cartera de proyectos en el orden de codificación R1 a R8.
  Marca llena: el proyecto de la fila no puede empezar sin el resultado del
  proyecto de la columna. Marca hueca: puede empezar, pero a mayor coste o
  riesgo. Elaboración propia.}
  \label{fig:dsmproy}
\end{figure}

\subsec{06. Secuenciación de la cartera}

Aplicado el mismo particionado, la matriz de la cartera resulta ser
completamente triangular inferior: no hay ningún ciclo, de modo que existe una
secuencia que no obliga a rehacer nada. La Figura~\ref{fig:dsmproyord} la
recoge y la traduce en seis olas.

\begin{figure}[ht]
  \centering
  \begin{tikzpicture}[x=8.4mm, y=8.4mm, font=\scriptsize]
    \def\n{8}
    % Orden: R2 R1 R3 R8 R4 R5 R6 R7
    % Posición: 1=R2 2=R1 3=R3 4=R8 5=R4 6=R5 7=R6 8=R7
    \def\fuerte{2/1, 3/2, 4/2, 5/3, 6/3, 7/1, 7/3, 7/6, 8/3, 8/5}
    \def\debil{3/1, 4/1, 5/2, 6/1, 8/7}
    \fill[gray!4] (0,0) rectangle (\n,-\n);
    \foreach \i in {1,...,8}
      \fill[gray!35] (\i-1,-\i+1) rectangle (\i,-\i);
    \foreach \r/\c in \fuerte
      \fill[dsmpro] (\c-1+0.16,-\r+1-0.16) rectangle (\c-0.16,-\r+0.16);
    \foreach \r/\c in \debil
      \draw[dsmpro, thick, fill=dsmpro!20]
        (\c-1+0.16,-\r+1-0.16) rectangle (\c-0.16,-\r+0.16);
    \draw[gray!55, thin, step=1] (0,0) grid (\n,-\n);
    \draw[thick] (0,0) rectangle (\n,-\n);
    % ---- escalonado de las olas, sobre la rejilla
    \foreach \a/\b in {0/1, 1/2, 2/4, 4/6, 6/7, 7/8}
      \draw[gray!75!black, line width=1pt] (\a,-\a) -- (\b,-\a) -- (\b,-\b);
    \foreach \i/\lab in {1/R2,2/R1,3/R3,4/R8,5/R4,6/R5,7/R6,8/R7} {
      \node[anchor=east] at (-0.2,-\i+0.5) {\lab};
      \node[anchor=south] at (\i-0.5,0.1) {\lab};
    }
    \foreach \y/\ola in {0.5/{ola 1}, 1.5/{ola 2}, 3/{ola 3}, 5/{ola 4},
                         6.5/{ola 5}, 7.5/{ola 6}}
      \node[anchor=west, font=\scriptsize\itshape, text=gray!55!black]
        at (\n+0.35,-\y) {\ola};
  \end{tikzpicture}
  \caption{DSM de la cartera particionada. Todas las marcas quedan por debajo
  de la diagonal: la cartera admite una secuencia sin iteración. R3 y R8 caen
  en la misma ola porque no dependen entre sí, y lo mismo ocurre con R4 y R5.
  Elaboración propia.}
  \label{fig:dsmproyord}
\end{figure}

La secuencia resultante es R2; R1; R3 y R8 en paralelo; R4 y R5 en paralelo;
R6; y por último R7. Tres consecuencias merecen comentario.

\textbf{La custodia del estado va primero, y sola.} R2 no depende de nada y
casi todo depende de él. Es el único proyecto de la cartera que no admite ser
pospuesto, porque importar recursos a un registro de estado sin bloqueo ni
cifrado, que es lo que haría R1 sin él, multiplica el trabajo por dos: primero
se importa y después se migra.

\textbf{La observabilidad no es el remate, es la ola 3.} R8 aparece pronto
porque sin él las figuras de mérito del Cuadro~\ref{tab:fom} no son medibles,
y sin medida las posiciones del gráfico vectorial de la
sección~\ref{sec:comp} siguen siendo estimaciones. Situarlo al final, que es
la tentación natural, deja la hoja de ruta sin instrumento de verificación
justo durante los años en que se ejecuta.

\textbf{La matriz confirma el orden discutido del catálogo y el motor.} El
apartado 05 de la sección~\ref{sec:comp} defiende en prosa que el catálogo de
módulos (R3) debe construirse antes que la sustitución del motor (R6), y
reconoce que es la decisión más discutible de aquella sección. La DSM lo
convierte en una constatación estructural: R6 tiene marca necesaria en las
columnas R2, R3 y R5, de modo que en cualquier orden admisible R6 queda por
detrás de los tres. Invertir el orden no es una preferencia distinta, es una
marca por encima de la diagonal, esto es, una iteración forzada.

\subsec{07. Riesgos que la matriz hace visibles}

Cruzar las dos matrices identifica dónde conviene poner las reservas de
esfuerzo, que es el uso que la gestión del portafolio hace de la DSM.

\begin{table}[!ht]
  \centering
  \footnotesize
  \caption{Riesgos derivados de la estructura de dependencias y respuesta
  prevista.}
  \label{tab:dsmriesgo}
  \begin{tabularx}{\textwidth}{@{}>{\raggedright\arraybackslash}p{0.235\textwidth}
    >{\raggedright\arraybackslash}X
    >{\raggedright\arraybackslash}p{0.30\textwidth}@{}}
    \toprule
    \textbf{Origen en la matriz} & \textbf{Riesgo} & \textbf{Respuesta} \\
    \midrule
    M1 con ocho salidas & Una indisponibilidad del repositorio detiene ocho
      componentes; su compromiso los compromete a todos. &
      Réplica del repositorio y firma obligatoria de confirmaciones. \\
    M3 con seis salidas y ola temprana & Si la gestión de secretos llega tarde,
      lo construido antes queda con credenciales incrustadas. &
      Adelantarla a la ola 2 de la arquitectura, antes que cualquier capa
      intermedia. \\
    Bloque acoplado L1--M4 & La política evaluada después de aplicar convierte
      el incumplimiento en un hecho consumado. &
      Partir el motor en planificación y aplicación, con la evaluación
      intercalada, según el diagrama SD1 de la sección~\ref{sec:opm}. \\
    R2 con cinco entradas en su columna & Un retraso en la custodia del estado
      desplaza toda la cartera. &
      Ola 1 en solitario, con criterio de terminación explícito antes de
      abrir R1. \\
    R3 con cuatro entradas en su columna & El catálogo es el cuello de botella
      de la cartera: R4, R5, R6 y R7 lo esperan. &
      Publicar el catálogo por incrementos y abrir R4 y R5 contra los primeros
      módulos, sin esperar al catálogo completo. \\
    L3 y L6 con cero salidas & Ningún componente los espera, luego son los
      candidatos naturales a recortar si el calendario se comprime. &
      Registrar el recorte como decisión consciente, no como omisión. \\
    \bottomrule
  \end{tabularx}
\end{table}

\subsec{08. Límites de las dos matrices}

Tres reservas. La primera es que ambas matrices son \textbf{binarias en el
tipo}: registran que hay dependencia, no de qué naturaleza es. Una DSM más
elaborada distinguiría dependencias de información, de material, de energía y
espaciales, y la de este sistema sería casi enteramente de información, lo que
resta valor a la distinción.

La segunda es que la DSM de arquitectura describe una \textbf{implantación
completa desde cero}. Una organización que ya opera parte del sistema, que es
el caso de la empresa de referencia, no recorre las olas en orden sino que
entra en la secuencia por el punto en que se encuentra; el orden sigue siendo
válido para los componentes que faltan.

La tercera es que la matriz de cartera \textbf{no incorpora el calendario ni
el coste}. Dice qué debe ir antes que qué, no cuánto dura cada cosa ni cuántos
equipos caben en paralelo. Convertir estas olas en fechas exige el modelo
financiero, el elemento 8 del índice, que este documento declara fuera de
alcance.

% ==========================================================================
\section{Modelo de objetos y procesos (OPM)}
\label{sec:opm}
% ==========================================================================

\subsec{01. Por qué OPM y cómo se ha construido el modelo}

Esta sección corresponde al elemento 3 del índice de una hoja de ruta
tecnológica. La metodología de objetos y procesos, normalizada como ISO/PAS
19450:2015 y después como ISO 19450:2024
\parencite{iso19450:2015,iso19450:2024,dori2002,dori2016}, presenta dos rasgos
que la hacen adecuada aquí. El primero es que usa un único tipo de diagrama
para representar a la vez estructura y comportamiento, lo que evita el
problema habitual de mantener sincronizados un diagrama de bloques y un
diagrama de actividad. El segundo es su carácter bimodal: cada diagrama tiene
una traducción a lenguaje natural controlado, el OPL, de modo que el modelo se
verifica leyéndolo, y un enlace mal puesto se manifiesta como una frase falsa.

Las herramientas autorizadas para este elemento son OPCloud
\parencite{opcloud} y OPCAT \parencite{opcat}, y el modelo está enunciado en
el vocabulario cerrado que ambas admiten: cada entidad, cada enlace y cada
frase del OPL de los apartados siguientes se corresponden con un elemento de
la norma y con ningún otro. Las figuras de esta sección son la composición
tipográfica de ese mismo modelo, dibujada con la notación de la norma para que
quepa en la caja del documento y pueda referirse desde el texto.

\subsec{02. Vocabulario y notación empleados}

OPM tiene un vocabulario cerrado, y la mayor parte de los errores de modelado
consisten en usar un símbolo por otro. El Cuadro~\ref{tab:opmnotacion} fija el
que se usa en las tres figuras siguientes.

\begin{table}[htbp]
  \centering
  \footnotesize
  \caption{Vocabulario de OPM empleado en las Figuras \ref{fig:opmsd},
  \ref{fig:opmsd2} y \ref{fig:opmsd1}, con la frase OPL que genera cada
  símbolo.}
  \label{tab:opmnotacion}
  \begin{tabularx}{\textwidth}{@{}>{\raggedright\arraybackslash}p{0.215\textwidth}
    >{\raggedright\arraybackslash}p{0.305\textwidth}
    >{\raggedright\arraybackslash}X@{}}
    \toprule
    \textbf{Elemento} & \textbf{Símbolo} & \textbf{Frase OPL} \\
    \midrule
    \multicolumn{3}{@{}l}{\textit{Entidades}} \\
    Objeto & Rectángulo de contorno verde & \\
    Proceso & Elipse de contorno azul & \\
    Estado & Rectángulo redondeado dentro del objeto &
      \emph{O} puede estar \emph{s1}, \emph{s2} o \emph{s3}. \\
    Cosa física & Con sombra proyectada & \emph{C} es física. \\
    Cosa ambiental & Contorno discontinuo & \emph{C} es ambiental. \\
    \midrule
    \multicolumn{3}{@{}l}{\textit{Enlaces procedimentales}} \\
    Agente & Línea acabada en círculo relleno &
      \emph{O} maneja \emph{P}. \\
    Instrumento & Línea acabada en círculo hueco &
      \emph{P} requiere \emph{O}. \\
    Efecto & Flecha de doble punta &
      \emph{P} afecta a \emph{O}. \\
    Consumo & Flecha del objeto al proceso &
      \emph{P} consume \emph{O}. \\
    Resultado & Flecha del proceso al objeto &
      \emph{P} produce \emph{O}. \\
    Entrada y salida & Flechas entre estados concretos y el proceso &
      \emph{P} cambia \emph{O} de \emph{s1} a \emph{s2}. \\
    \midrule
    \multicolumn{3}{@{}l}{\textit{Enlaces estructurales}} \\
    Agregación & Triángulo relleno &
      \emph{A} consta de \emph{B}. \\
    Generalización & Triángulo hueco &
      \emph{B} es un \emph{A}. \\
    Exhibición & Triángulo dentro de triángulo &
      \emph{A} exhibe \emph{B}. \\
    \midrule
    \multicolumn{3}{@{}l}{\textit{Mecanismos de refinamiento}} \\
    Ampliación en detalle & El proceso ampliado contiene a sus subprocesos,
      ordenados de arriba abajo por precedencia temporal &
      \emph{P} se amplía en \emph{P1}, \emph{P2} y \emph{P3}. \\
    Despliegue & Diagrama aparte con la jerarquía estructural de un objeto &
      (las de sus enlaces estructurales) \\
    \bottomrule
  \end{tabularx}
\end{table}

Dos precisiones sobre el uso de estos símbolos, porque son las dos que con más
frecuencia se incumplen. La primera: la precedencia temporal entre subprocesos
de una ampliación se expresa \textbf{únicamente} por la posición vertical, y
no dibujando flechas entre ellos; una flecha entre dos procesos sería un
enlace de invocación, que significa otra cosa. La segunda: cuando el diagrama
muestra los estados de un objeto y el proceso transforma uno concreto en otro,
los enlaces deben partir y llegar \textbf{al estado} y no al objeto entero; el
enlace de efecto al objeto entero sólo es correcto cuando el cambio de estado
no se especifica.

\subsec{03. Diagrama de sistema (SD)}

El diagrama de sistema recoge el proceso principal, quién lo maneja, qué
instrumentos necesita y qué objetos transforma o produce. Los estados de la
\emph{Infraestructura Desplegada} se declaran aquí, pero las transiciones
concretas entre ellos no se especifican todavía: a este nivel el proceso la
afecta, sin más, y las transiciones aparecen en la ampliación en detalle. El
Cuadro~\ref{tab:entidades} enumera las entidades.

\begin{table}[htbp]
  \centering
  \scriptsize
  \caption{Entidades del modelo OPM y diagrama en el que aparecen.}
  \label{tab:entidades}
  \begin{tabularx}{\textwidth}{@{}>{\raggedright\arraybackslash}p{0.245\textwidth}
    >{\raggedright\arraybackslash}p{0.155\textwidth}
    >{\raggedright\arraybackslash}X@{}}
    \toprule
    \textbf{Entidad} & \textbf{Tipo} & \textbf{Papel en el modelo} \\
    \midrule
    Aprovisionamiento de Infraestructura & Proceso & Proceso principal; se
      amplía en detalle en SD1. \\
    Equipo de Plataforma & Objeto físico y ambiental & Agente humano que
      maneja el proceso. \\
    Especificación de Infraestructura & Objeto informático & El código;
      instrumento del proceso. Se despliega en SD2. Estados: borrador,
      validada. \\
    Repositorio Versionado & Objeto informático & Todo del que la
      especificación es parte (SD2). \\
    Módulo Reutilizable, Configuración de Entorno & Objetos informáticos &
      Especializaciones de la especificación (SD2). \\
    Versión & Objeto informático & Atributo exhibido por la especificación
      (SD2). \\
    Motor de IaC & Objeto informático & Instrumento que ejecuta la
      comparación y la conciliación. \\
    Registro de Estado & Objeto informático & Objeto afectado: el proceso lo
      lee y lo reescribe en cada aplicación. \\
    Proveedor de Nube & Objeto físico y ambiental & Instrumento: expone la
      interfaz de programación. \\
    Credenciales & Objeto informático & Instrumento con el que el motor se
      autentica. \\
    Conjunto de Políticas & Objeto informático & Instrumento de la
      validación. \\
    Infraestructura Desplegada & Objeto físico & Operando del sistema.
      Estados: inexistente, conforme, desviada, retirada. \\
    Plan de Cambios & Objeto informático & Resultado de la planificación,
      instrumento de la aplicación. Estados: pendiente, aprobado. \\
    Registro de Auditoría & Objeto informático & Resultado del proceso. \\
    \bottomrule
  \end{tabularx}
\end{table}

\begin{figure}[H]
  \centering
  \begin{tikzpicture}[scale=0.86, transform shape,
                      every node/.style={transform shape}]
    % ---- proceso principal
    \node[opmproc, minimum width=44mm, minimum height=17mm] (P) at (0,0)
      {Aprovisionamiento\\de Infraestructura};

    % ---- agente humano: objeto físico y ambiental
    \node[opmobj, opmamb, opmfis, text width=24mm] (eq) at (0,3.05)
      {Equipo de\\Plataforma};
    \draw[agente] (eq) -- (P);

    % ---- instrumentos
    \node[opmobj] (spec)  at (-5.8, 1.85) {Especificación de\\Infraestructura};
    \node[opmobj] (motor) at (-5.8, 0.35) {Motor de IaC};
    \node[opmobj] (pol)   at (-5.8,-1.15) {Conjunto de\\Políticas};
    \node[opmobj, opmamb, opmfis] (prov) at (5.8, 1.85)
      {Proveedor\\de Nube};
    \node[opmobj] (cred)  at ( 5.8, 0.35) {Credenciales};
    \foreach \nd in {spec,motor,pol,prov,cred} \draw[instr] (\nd) -- (P);

    % ---- objeto afectado
    \node[opmobj] (edo) at (-5.8,-2.75) {Registro\\de Estado};
    \draw[efecto] (edo) -- (P);

    % ---- resultados
    \node[opmobj] (plan) at (5.8,-1.15) {Plan de Cambios};
    \node[opmobj] (aud)  at (5.8,-2.75) {Registro de\\Auditoría};
    \draw[flecha] (P) -- (plan);
    \draw[flecha] (P) -- (aud);

    % ---- operando: objeto físico con sus cuatro estados
    \node[draw=opmobjc, very thick, fill=opmobjc!7, opmfis,
      minimum width=84mm, minimum height=18mm] (infra) at (0,-4.35) {};
    \node[anchor=north, font=\footnotesize\bfseries]
      at ([yshift=-1.6mm]infra.north) {Infraestructura Desplegada};
    \foreach \dx/\st in {-30/inexistente, -10/conforme, 10/desviada,
                          30/retirada}
      \node[opmstate] at ([xshift=\dx mm, yshift=-4.6mm]infra.center) {\st};
    \draw[efecto] (P) -- (infra);
  \end{tikzpicture}
  \caption{Diagrama de sistema (SD). El proceso principal, su agente, sus
  cinco instrumentos, el objeto que afecta, los dos que produce y el operando
  con sus cuatro estados. El \emph{Equipo de Plataforma} y el \emph{Proveedor
  de Nube} llevan contorno discontinuo por ser ambientales y sombra por ser
  físicos; la \emph{Infraestructura Desplegada} lleva sombra por ser física.
  Elaboración propia según \textcite{dori2016} y \textcite{iso19450:2024}.}
  \label{fig:opmsd}
\end{figure}

\subsec{04. El OPL del diagrama de sistema}

Las palabras reservadas del lenguaje van en versalitas los nombres de las
entidades y en redonda las palabras estructurales. Cada frase se corresponde
con un símbolo de la Figura~\ref{fig:opmsd} y con ningún otro.

\begin{quote}\small
\textsc{Equipo de Plataforma} es físico y ambiental.\\
\textsc{Proveedor de Nube} es físico y ambiental.\\
\textsc{Infraestructura Desplegada} es física.\\
\textsc{Infraestructura Desplegada} puede estar inexistente, conforme,
desviada o retirada.\\
\textsc{Equipo de Plataforma} maneja \textsc{Aprovisionamiento de
Infraestructura}.\\
\textsc{Aprovisionamiento de Infraestructura} requiere \textsc{Especificación
de Infraestructura}, \textsc{Motor de IaC}, \textsc{Conjunto de Políticas},
\textsc{Credenciales} y \textsc{Proveedor de Nube}.\\
\textsc{Aprovisionamiento de Infraestructura} afecta a \textsc{Registro de
Estado}.\\
\textsc{Aprovisionamiento de Infraestructura} afecta a \textsc{Infraestructura
Desplegada}.\\
\textsc{Aprovisionamiento de Infraestructura} produce \textsc{Plan de Cambios}
y \textsc{Registro de Auditoría}.
\end{quote}

Leído en voz alta, el conjunto es una descripción funcional completa del nivel
superior del sistema.

\subsec{05. Despliegue de la especificación (SD2)}

La \emph{Especificación de Infraestructura} tiene una estructura interna que
el diagrama de sistema no muestra, porque mezclarla allí saturaría el nivel
superior. OPM prevé para eso el despliegue: un diagrama aparte que expone la
jerarquía estructural de un objeto sin tocar su comportamiento. La
Figura~\ref{fig:opmsd2} lo recoge con las tres relaciones estructurales de la
norma, que son las tres que este objeto necesita.

\begin{figure}[H]
  \centering
  \begin{tikzpicture}[scale=0.92, transform shape,
                      every node/.style={transform shape}]
    \node[opmobj, text width=32mm] (spec) at (0,0)
      {Especificación de\\Infraestructura};

    % ---- agregación: el repositorio consta de la especificación
    \node[opmobj, text width=30mm] (repo) at (0,2.3)
      {Repositorio\\Versionado};
    \draw[very thick] (repo.south) -- (spec.north);
    \coordinate (ca) at ($(repo.south)!0.45!(spec.north)$);
    \opmagg{0}{ca}
    \node[etiq, anchor=west] at ([xshift=3mm]ca)
      {agregación: \emph{consta de}};

    % ---- exhibición: la especificación exhibe el atributo Versión
    \node[opmobj, text width=17mm] (ver) at (5.9,0) {Versión};
    \draw[very thick] (spec.east) -- (ver.west);
    \coordinate (ce) at ($(spec.east)!0.40!(ver.west)$);
    \opmexh{90}{ce}
    \node[etiq, anchor=south] at ([yshift=2.6mm]$(spec.east)!0.62!(ver.west)$)
      {exhibición: \emph{exhibe}};

    % ---- generalización: dos especializaciones
    \node[opmobj, text width=28mm] (mod) at (-3.0,-2.8)
      {Módulo\\Reutilizable};
    \node[opmobj, text width=28mm] (cfg) at ( 3.0,-2.8)
      {Configuración\\de Entorno};
    \coordinate (cg) at (0,-1.15);
    \draw[very thick] (spec.south) -- (0,-1.75);
    \draw[very thick] (mod.north) |- (0,-1.75) -| (cfg.north);
    \opmgen{0}{cg}
    \node[etiq, anchor=west] at ([xshift=3mm]cg)
      {generalización: \emph{es un}};
  \end{tikzpicture}
  \caption{SD2: despliegue del objeto \emph{Especificación de
  Infraestructura}. Agregación con el repositorio que la contiene,
  generalización en sus dos especializaciones y exhibición del atributo que la
  versiona. La base de cada triángulo se apoya en el todo, el general o el
  exhibidor, y el vértice apunta a la parte, la especialización o el atributo.
  Elaboración propia.}
  \label{fig:opmsd2}
\end{figure}

\begin{quote}\small
\textsc{Repositorio Versionado} consta de \textsc{Especificación de
Infraestructura}.\\
\textsc{Especificación de Infraestructura} exhibe \textsc{Versión}.\\
\textsc{Módulo Reutilizable} es una \textsc{Especificación de
Infraestructura}.\\
\textsc{Configuración de Entorno} es una \textsc{Especificación de
Infraestructura}.
\end{quote}

Conviene subrayar la diferencia entre la tercera frase y una que se le parece
mucho pero significa otra cosa. \emph{Es un} declara una especialización, esto
es, una clase de objeto distinta que hereda las propiedades de la general.
\emph{Puede estar} declara un estado, esto es, una situación transitoria del
mismo objeto. Un módulo reutilizable no es un estado de la especificación, es
un tipo de especificación; borrador y validada, en cambio, sí son estados.

\subsec{06. Ampliación en detalle del proceso principal (SD1)}

El diagrama SD1 amplía el proceso principal en siete subprocesos. La
convención de la norma es que la posición vertical, y sólo ella, denota
precedencia temporal: no hay flechas entre subprocesos. Los objetos que ya
aparecen en el diagrama de sistema permanecen fuera de la elipse ampliada, y
los enlaces cruzan su contorno hasta el subproceso concreto que los usa, que
es lo que permite ver en qué punto exacto interviene cada uno. Los estados de
la \emph{Especificación de Infraestructura}, del \emph{Plan de Cambios} y de
la \emph{Infraestructura Desplegada} se hacen explícitos aquí, y con ellos las
transiciones.

\begin{figure}[H]
  \centering
  \begin{tikzpicture}[scale=0.79, transform shape, x=1mm, y=1mm,
                      every node/.style={transform shape}]
    \tikzset{sp/.style={opmproc, minimum width=46mm, minimum height=10mm,
               text width=, font=\footnotesize\bfseries},
             ob/.style={opmobj, text width=23mm, minimum height=8mm,
               font=\scriptsize\bfseries},
             obe/.style={draw=opmobjc, very thick, fill=opmobjc!7,
               minimum width=34mm},
             est/.style={opmstate, font=\tiny, minimum width=20mm}}

    % ---- elipse del proceso ampliado
    \node[draw=opmprocc, very thick, ellipse, fill=opmprocc!5,
      minimum width=100mm, minimum height=124mm] (marco) at (0,0) {};
    \node[font=\scriptsize\bfseries, align=center] at (0,50)
      {Aprovisionamiento\\de Infraestructura};

    % ---- los siete subprocesos: la posición vertical fija la precedencia
    \node[sp] (c)  at (0, 36) {Codificación};
    \node[sp] (v)  at (0, 24) {Validación};
    \node[sp] (pl) at (0, 12) {Planificación};
    \node[sp] (au) at (0,  0) {Autorización};
    \node[sp] (ap) at (0,-12) {Aplicación};
    \node[sp] (co) at (0,-24) {Conciliación};
    \node[sp] (re) at (0,-36) {Retirada};

    % ---- eje temporal, fuera de todo
    \draw[flecha, gray!65] (-106,42) -- (-106,-42)
      node[midway, rotate=90, above=0.6mm, font=\scriptsize,
           text=gray!55!black] {precedencia temporal};

    % ---- agente, fuera del proceso ampliado
    \node[ob, opmamb, opmfis, text width=21mm] (eq) at (-72,48)
      {Equipo de\\Plataforma};

    % ---- objetos de la izquierda
    \node[obe, minimum height=28mm] (spec) at (-72,28) {};
    \node[anchor=north, font=\scriptsize\bfseries, align=center]
      at ([yshift=-1.1mm]spec.north) {Especificación de\\Infraestructura};
    \node[est] (sb) at (-72,30) {borrador};
    \node[est] (sv) at (-72,22) {validada};

    \node[ob] (edo)   at (-72, 8) {Registro\\de Estado};
    \node[ob] (motor) at (-72,-6) {Motor de IaC};

    \node[obe, minimum height=40mm, opmfis] (infra) at (-72,-33) {};
    \node[anchor=north, font=\scriptsize\bfseries, align=center]
      at ([yshift=-1.1mm]infra.north) {Infraestructura\\Desplegada};
    \node[est] (si) at (-72,-24) {inexistente};
    \node[est] (sc) at (-72,-32) {conforme};
    \node[est] (sd) at (-72,-40) {desviada};
    \node[est] (sr) at (-72,-48) {retirada};

    % ---- objetos de la derecha
    \node[ob] (pol) at (72,26) {Conjunto de\\Políticas};
    \node[obe, minimum height=20mm] (plan) at (72,6) {};
    \node[anchor=north, font=\scriptsize\bfseries] at ([yshift=-1.1mm]plan.north)
      {Plan de Cambios};
    \node[est] (pp) at (72,7.5) {pendiente};
    \node[est] (pa) at (72,0.5) {aprobado};
    \node[ob] (cred) at (72,-16) {Credenciales};
    \node[ob, opmamb, opmfis] (prov) at (72,-30) {Proveedor\\de Nube};
    \node[ob] (aud)  at (72,-46) {Registro de\\Auditoría};

    % ---- enlaces del agente
    \draw[agente] (eq.south east) to[out=-40,in=155] (c.north west);
    \draw[agente, rounded corners=3mm]
      (eq.west) -- (-96,48) -- (-96,0) -- (au.west);

    % ---- codificación, validación y planificación
    \draw[flecha] (c.west) to[out=180,in=10] ([yshift=1mm]sb.east);
    \draw[flecha] ([yshift=-1mm]sb.east) to[out=-10,in=170] (v.north west);
    \draw[flecha] (v.west) to[out=180,in=10] ([yshift=1mm]sv.east);
    \draw[instr]  ([yshift=-1mm]sv.east) to[out=-10,in=165] (pl.north west);
    \draw[instr]  (edo.east) -- (pl.west);
    \draw[instr]  (motor.north east) to[out=25,in=195] (pl.south west);

    % ---- autorización y aplicación
    \draw[flecha] (pl.east) to[out=0,in=170] (pp.north west);
    \draw[instr]  (pp.south west) to[out=-165,in=15] (au.east);
    \draw[flecha] (au.north east) to[out=15,in=190] (pa.west);
    \draw[instr]  (pa.south west) to[out=-160,in=20] (ap.east);
    \draw[instr]  (motor.east) to[out=0,in=170] (ap.north west);
    \draw[instr]  (cred.west) -- (ap.east);
    \draw[instr]  (prov.north west) to[out=155,in=-15] (ap.south east);
    \draw[efecto] (ap.west) to[out=170,in=-25] (edo.south east);
    \draw[flecha] (si.east) to[out=0,in=200] (ap.south west);
    \draw[flecha] (ap.west) to[out=195,in=15] ([yshift=1.4mm]sc.east);

    % ---- conciliación y retirada
    \draw[instr]  (prov.west) to[out=180,in=-15] (co.east);
    \draw[flecha] (sd.east) to[out=0,in=195] (co.west);
    \draw[flecha] (co.north west) to[out=160,in=-5] ([yshift=-1.4mm]sc.east);
    \draw[flecha] (co.south east) to[out=-40,in=165] (aud.west);
    \draw[instr]  (prov.south west) to[out=205,in=-10] (re.east);
    \draw[flecha, rounded corners=2mm]
      ([yshift=-3.2mm]sc.east) -- (-49,-35.2) -- (-32,-35.2) -- (re.west);
    \draw[flecha] (re.south west) to[out=210,in=-10] (sr.south east);
  \end{tikzpicture}
  \caption{SD1: ampliación en detalle de \emph{Aprovisionamiento de
  Infraestructura} en sus siete subprocesos. La precedencia temporal la fija
  la posición vertical, sin flechas entre procesos. Los objetos que ya
  aparecen en el diagrama de sistema quedan fuera de la elipse y se enlazan
  con el subproceso que los usa. \emph{Aplicación}, \emph{Conciliación} y
  \emph{Retirada} actúan sobre el mismo objeto desde estados de partida
  distintos, que es la formulación precisa de la diferencia entre desplegar,
  corregir la deriva y desmantelar. Elaboración propia.}
  \label{fig:opmsd1}
\end{figure}

\subsec{07. El OPL de la ampliación}

\begin{quote}\small
\textsc{Aprovisionamiento de Infraestructura} se amplía en
\textsc{Codificación}, \textsc{Validación}, \textsc{Planificación},
\textsc{Autorización}, \textsc{Aplicación}, \textsc{Conciliación} y
\textsc{Retirada}, en ese orden temporal.\\
\textsc{Especificación de Infraestructura} puede estar borrador o validada.\\
\textsc{Plan de Cambios} puede estar pendiente o aprobado.\\
\textsc{Equipo de Plataforma} maneja \textsc{Codificación} y
\textsc{Autorización}.\\
\textsc{Codificación} produce \textsc{Especificación de Infraestructura}
borrador.\\
\textsc{Validación} requiere \textsc{Conjunto de Políticas}.\\
\textsc{Validación} cambia \textsc{Especificación de Infraestructura} de
borrador a validada.\\
\textsc{Planificación} requiere \textsc{Motor de IaC}, \textsc{Registro de
Estado} y \textsc{Especificación de Infraestructura} validada.\\
\textsc{Planificación} produce \textsc{Plan de Cambios} pendiente.\\
\textsc{Autorización} requiere \textsc{Plan de Cambios} pendiente.\\
\textsc{Autorización} cambia \textsc{Plan de Cambios} de pendiente a
aprobado.\\
\textsc{Aplicación} requiere \textsc{Motor de IaC}, \textsc{Plan de Cambios}
aprobado, \textsc{Credenciales} y \textsc{Proveedor de Nube}.\\
\textsc{Aplicación} cambia \textsc{Infraestructura Desplegada} de inexistente
a conforme.\\
\textsc{Aplicación} afecta a \textsc{Registro de Estado}.\\
\textsc{Conciliación} requiere \textsc{Proveedor de Nube}.\\
\textsc{Conciliación} cambia \textsc{Infraestructura Desplegada} de desviada a
conforme.\\
\textsc{Conciliación} produce \textsc{Registro de Auditoría}.\\
\textsc{Retirada} requiere \textsc{Proveedor de Nube}.\\
\textsc{Retirada} cambia \textsc{Infraestructura Desplegada} de conforme a
retirada.
\end{quote}

\subsec{08. Qué hace explícito el modelo}

Merece la pena señalar cuatro cosas que el modelo obliga a decir y que las
descripciones en prosa dejan implícitas.

\begin{enumerate}
  \item \textsc{Registro de Estado} es el único objeto que aparece a la vez
    como instrumento de la planificación y como objeto afectado por la
    aplicación. Esa doble condición es exactamente la razón por la que
    necesita bloqueo de concurrencia, y el modelo la hace visible sin
    argumentarla.
  \item \textsc{Aplicación}, \textsc{Conciliación} y \textsc{Retirada} cambian
    el mismo objeto desde estados de partida distintos. Con enlaces de efecto
    al objeto entero las tres serían indistinguibles; con enlaces de entrada y
    salida entre estados quedan formalmente separadas, que es la diferencia
    entre desplegar, corregir la deriva y desmantelar.
  \item El equipo humano interviene sólo en dos de los siete subprocesos, la
    codificación y la autorización. Eso delimita el alcance real de la
    automatización y, de paso, señala dónde está el cuello de botella de
    personas.
  \item La \textsc{Validación} se sitúa entre la \textsc{Codificación} y la
    \textsc{Planificación}, y la \textsc{Autorización} entre la
    \textsc{Planificación} y la \textsc{Aplicación}. Esa partición es la que
    rompe el bloque acoplado L1--M4 que la
    Figura~\ref{fig:dsmarqord} detecta en la matriz de dependencias: el
    veredicto de la política deja de necesitar un motor que a su vez lo
    necesita a él, porque cada uno actúa sobre un subproceso distinto.
\end{enumerate}

% ==========================================================================
\section{Figuras de mérito del sistema}
\label{sec:fom}
% ==========================================================================

Esta sección y la siguiente componen el elemento 4 del índice de una hoja de
ruta tecnológica. Ésta define las figuras de mérito y fija sus anclas; la
sección~\ref{sec:evol} las aplica para observar cómo ha evolucionado la
tecnología.

\subsec{01. Las ocho figuras}

Un sistema de IaC es en sí mismo una tecnología, y como tal admite figuras de
mérito que permiten situarlo y seguir su evolución. El
Cuadro~\ref{tab:fom} propone ocho, con su unidad y su procedimiento de
medida. Las cuatro primeras coinciden con las métricas de rendimiento de
entrega ampliamente utilizadas \parencite{forsgren2018}; las cuatro últimas
son específicas de la infraestructura.

\begin{table}[htbp]
  \centering
  \footnotesize
  \caption{Figuras de mérito propuestas para un sistema de IaC.}
  \label{tab:fom}
  \begin{tabularx}{\textwidth}{@{}>{\raggedright\arraybackslash}p{0.27\textwidth}
    >{\raggedright\arraybackslash}p{0.17\textwidth}
    >{\raggedright\arraybackslash}X@{}}
    \toprule
    \textbf{Figura de mérito} & \textbf{Unidad} & \textbf{Cómo se mide} \\
    \midrule
    Tiempo de ciclo del cambio & horas & Del registro del cambio en el
      repositorio a su aplicación efectiva en producción. \\
    Frecuencia de aplicación & cambios/semana & Recuento de aplicaciones
      correctas por entorno. \\
    Tasa de fallo del cambio & \% & Aplicaciones que exigen reversión o
      corrección inmediata sobre el total. \\
    Tiempo de restablecimiento & minutos & De la detección del fallo a la
      recuperación del servicio. \\
    Cobertura de codificación & \% & Recursos gestionados por el motor sobre
      el total de recursos existentes en el proveedor. \\
    Tasa de deriva & recursos desviados / 1000 · semana & Diferencias
      detectadas en las ejecuciones de sólo lectura. \\
    Tiempo de reconstrucción & horas & De un entorno vacío a un entorno
      equivalente al de producción, partiendo sólo del repositorio. \\
    Cumplimiento de política & \% & Planes que superan la evaluación de
      políticas sin excepción manual. \\
    \bottomrule
  \end{tabularx}
\end{table}

El tiempo de reconstrucción es la figura más reveladora de todas, porque
verifica de un solo golpe las cuatro propiedades exigidas en la
sección~\ref{sec:alcance}. Una
organización que no puede reconstruir un entorno completo desde el repositorio
no tiene un sistema de IaC completo, con independencia de cuántos ficheros
declarativos posea.


\subsec{02. Por qué estas ocho y no otras}

Una figura de mérito no es cualquier métrica que pueda recogerse. Para servir
en una hoja de ruta tiene que cumplir cuatro condiciones, y las ocho del
Cuadro~\ref{tab:fom} se han seleccionado exigiéndolas todas
\parencite{deweck2022}.

\begin{enumerate}
  \item \textbf{Es observable con un procedimiento único.} La columna derecha
    del Cuadro~\ref{tab:fom} no es un adorno: si dos observadores pueden
    obtener valores distintos, la figura no sirve para comparar posiciones ni
    para seguir una evolución.
  \item \textbf{Tiene un sentido de mejora inequívoco.} Cada figura mejora
    monótonamente en una sola dirección. El coste, por ejemplo, queda excluido
    del cuadro por este motivo y se trata aparte, como índice, en la
    sección~\ref{sec:comp}.
  \item \textbf{Es sensible a las decisiones de la hoja de ruta.} Una figura
    que no se mueva con ninguno de los ocho proyectos de la cartera es
    irrelevante, por interesante que resulte.
  \item \textbf{Tiene un límite conocido.} Sin límite superior o asíntota no
    se puede normalizar a utilidad ni ajustar una curva S, que es lo que hacen
    la sección~\ref{sec:comp} y la sección~\ref{sec:evol}.
\end{enumerate}

Las ocho se agrupan además en dos familias que conviene no mezclar, porque
sirven a interlocutores distintos. Las cuatro primeras son \emph{figuras de
entrega}: miden lo que el equipo de producto experimenta, existe literatura
que las mide de forma comparable entre organizaciones
\parencite{forsgren2018,dora2024} y son las que alimentan el eje de valor del
gráfico vectorial. Las cuatro últimas son \emph{figuras de infraestructura}:
miden propiedades del sistema que el equipo de producto no percibe
directamente pero que determinan su riesgo, y ninguna fuente pública las
publica de forma sistemática.

\subsec{03. Anclas, valor actual y objetivo}

Para que las ocho figuras sean operativas hace falta un tercer dato además de
la unidad y el procedimiento: las anclas que delimitan su recorrido útil. El
Cuadro~\ref{tab:anclas} las fija, junto con el valor observado en la
plataforma de 2026 y el objetivo comprometido para 2029, que la
sección~\ref{sec:estrategia} convierte en la declaración de estrategia.

\begin{table}[htbp]
  \centering
  \footnotesize
  \setlength{\tabcolsep}{4pt}
  \caption{Anclas de normalización, valor actual y objetivo de cada figura de
  mérito. El ancla \emph{peor} es el valor por debajo del cual el sistema se
  considera no apto; el ancla \emph{mejor} es la frontera de la tecnología
  estimada en la sección~\ref{sec:evol}. La columna \emph{familia} distingue
  las figuras de entrega (E) de las de infraestructura (I).}
  \label{tab:anclas}
  \begin{tabular}{@{}llrrrrr@{}}
    \toprule
    \textbf{Figura de mérito} & \textbf{Fam.} & \textbf{Peor} &
      \textbf{Mejor} & \textbf{2026} & \textbf{2029} &
      \textbf{Sentido} \\
    \midrule
    Tiempo de ciclo (h)            & E &  48 &   4 &  26 &   8 & menos mejor \\
    Frecuencia (cambios/sem.)      & E &   2 &  90 &  12 &  40 & más mejor \\
    Tasa de fallo (\pct)           & E &  20 &   3 &  14 &   6 & menos mejor \\
    Restablecimiento (min)         & E & 480 &  30 & 180 &  45 & menos mejor \\
    Cobertura (\pct)               & I &  40 &  95 &  55 &  86 & más mejor \\
    Tasa de deriva (rec./1000 sem.)& I &  60 &   4 &  38 &   8 & menos mejor \\
    Reconstrucción (h)             & I &  72 &   6 &  40 &  11 & menos mejor \\
    Cumplimiento (\pct)            & I &  50 &  98 &  62 &  94 & más mejor \\
    \bottomrule
  \end{tabular}
\end{table}

Cuatro de las ocho figuras, las de la familia de entrega marcadas en el
Cuadro~\ref{tab:anclas}, son las que la sección~\ref{sec:comp} normaliza a
utilidad para construir el eje de valor del gráfico vectorial. Las otras
cuatro no entran en ese eje porque el equipo de producto no las percibe, pero
sí gobiernan el modelo técnico de la sección~\ref{sec:tecnico}: son las
variables independientes de las que las primeras dependen.

La distancia entre las columnas de 2026 y de 2029 no es un deseo. Cada salto
está asignado a proyectos concretos de la cartera en el
Cuadro~\ref{tab:trazafom}, y su verosimilitud se comprueba en el modelo
técnico de la sección~\ref{sec:tecnico}, que deriva las figuras de entrega de
las de infraestructura en lugar de estimarlas por separado.

% ==========================================================================
\section{Evolución de las figuras de mérito}
\label{sec:evol}
% ==========================================================================

La sección anterior define las figuras de mérito; ésta las aplica para
observar cómo ha evolucionado la tecnología. Ambas componen el elemento 4 del
índice, que no se agota en proponer métricas: una figura de mérito sin serie
temporal no permite decidir nada, porque lo que orienta una hoja de ruta no es
el valor de una métrica sino su derivada \parencite{deweck2022}.

Se emplean los tres modelos disponibles, y no uno solo, porque cada uno
responde a una pregunta distinta y las tres respuestas resultan ser
divergentes, lo que constituye el hallazgo principal de la sección. La curva S
sitúa a la tecnología en su ciclo de vida \parencite{foster1986}; el ajuste
log-lineal al modo de la ley de Moore cuantifica el ritmo de mejora
\parencite{moore1965}; y el frente de Pareto identifica qué posiciones son
eficientes y cuánto cuesta avanzar sobre él.

\subsec{01. Qué se mide y de quién}

Antes de trazar una sola curva hay que distinguir dos objetos que las series
de figuras de mérito confunden con facilidad, y cuya confusión invalida
cualquier conclusión.

\begin{enumerate}
  \item \textbf{La frontera de la tecnología} es el mejor valor de cada figura
    de mérito alcanzable en un año dado con las herramientas y las prácticas
    disponibles en ese año, por una organización que las adopte
    correctamente. Es una propiedad de la tecnología, no de nadie en
    particular, y es lo que se modela en esta sección.
  \item \textbf{La posición de la empresa} es el valor que la organización
    concreta obtiene hoy. Es lo que recoge la fila de referencia del
    Cuadro~\ref{tab:fomcomp}, y está por detrás de la frontera: el retraso
    entre ambas, y no el valor absoluto, es lo que la hoja de ruta se propone
    cerrar.
\end{enumerate}

La distinción tiene una consecuencia práctica inmediata. Si la frontera de una
figura de mérito está saturada, invertir en acercarse a ella tiene un retorno
acotado y conocido; si sigue avanzando con rapidez, cualquier posición
alcanzada se degrada por sí sola en términos relativos, y la inversión hay que
repetirla. Saber en qué régimen está cada figura decide, por tanto, cuáles
merecen un proyecto y cuáles no.

Los valores de frontera empleados a continuación proceden de tres fuentes,
declaradas para que cada cifra sea rastreable: los hitos documentados de
capacidad de las herramientas, recogidos en el Cuadro~\ref{tab:hitos}; los
percentiles de los grupos de mayor rendimiento publicados en la serie anual de
informes sobre el estado de la práctica \parencite{forsgren2018,dora2024},
para las cuatro figuras de entrega; y la estimación razonada de ingeniería
para las cuatro figuras específicas de infraestructura, que ninguna fuente
pública mide de forma sistemática. La tercera fuente es la más débil y se
señala como tal en el apartado 09.

\subsec{02. La serie histórica de la tecnología}

El Cuadro~\ref{tab:hitos} fija los hitos que anclan las series. Se han
seleccionado por un criterio único: cada uno amplía lo que puede declararse y
reconciliarse como código, o bien cierra una de las cuatro propiedades
exigidas en la sección~\ref{sec:alcance}. Los hitos de mercado sin efecto
sobre la capacidad quedan fuera.

\begin{table}[htbp]
  \centering
  \footnotesize
  \setlength{\tabcolsep}{4pt}
  \caption{Hitos documentados que anclan la serie histórica de las figuras de
  mérito. La columna de la derecha indica qué propiedad de las exigidas en la
  sección~\ref{sec:alcance} queda cubierta por primera vez, o qué extensión
  del alcance declarable aporta el hito.}
  \label{tab:hitos}
  \begin{tabularx}{\textwidth}{@{}l>{\raggedright\arraybackslash}p{0.235\textwidth}
    >{\raggedright\arraybackslash}X@{}}
    \toprule
    \textbf{Año} & \textbf{Hito} & \textbf{Qué aporta} \\
    \midrule
    1993 & CFEngine \parencite{burgess1995} & Especificación externalizada y
      convergencia sobre la configuración del sistema operativo. No alcanza a
      la máquina ni a la red. \\
    1998 & ISconf y la máquina virtual de empresa
      \parencite{traugott1998} & Trata la infraestructura entera como un solo
      sistema distribuido, no como un conjunto de máquinas. \\
    2002 & Equivalencia de Turing en la administración de sistemas
      \parencite{traugott2002} & Demuestra que el orden de aplicación importa:
      funda la reproducibilidad como propiedad exigible. \\
    2006 & Interfaz de programación de cómputo elástico & Por primera vez la
      máquina misma, y no sólo su contenido, es un recurso declarable. \\
    2009 & Chef, y la infraestructura como código como término & Extiende la
      convergencia a la cadena de herramientas de despliegue. \\
    2011 & CloudFormation \parencite{cloudformationlaunch} & Primer motor
      declarativo de plantilla completa: 48 tipos de recurso, 13 de los 15
      servicios entonces disponibles. \\
    2013 & Pruebas de idempotencia \parencite{hummer2013} & Convierte la
      idempotencia en propiedad verificable y no en promesa. \\
    2014 & Terraform y Kubernetes \parencite{terraformdocs,burns2016} &
      Motor independiente del proveedor y bucle de conciliación continua: se
      cierra la cuarta propiedad. \\
    2017 & Registro público de módulos & La unidad de reutilización deja de ser
      el fichero y pasa a ser el módulo versionado. \\
    2018 & Pulumi \parencite{pulumilaunch} y Crossplane
      \parencite{crossplane} & Lenguajes de propósito general y plano de
      control como alternativas al lenguaje específico. \\
    2019 & Registro extensible de recursos \parencite{cfnregistry} y estudios
      de campo \parencite{rahman2019,guerriero2019} & El catálogo de recursos
      declarables se abre a terceros y la práctica pasa a estar medida. \\
    2020 & Bicep \parencite{bicepdocs} y política como código madura
      \parencite{opa} & La validación previa a la aplicación se generaliza. \\
    2023 & 3.000 proveedores y 12.000 módulos en el registro
      \parencite{tfecosystem}; bifurcación abierta del motor
      \parencite{hashicorpbusl,opentofuga} & La superficie declarable deja de
      ser el factor limitante; el gobierno de la licencia pasa a serlo. \\
    2026 & Estado de la práctica al inicio de esta hoja de ruta & Ninguna de
      las cuatro propiedades carece de solución disponible. \\
    \bottomrule
  \end{tabularx}
\end{table}

Una observación sobre la forma de la serie. Los cinco primeros hitos amplían
el alcance de lo declarable; los cuatro últimos no lo amplían apenas y
mejoran, en cambio, el gobierno de lo ya declarable. Ese cambio de naturaleza
en 2014, y no una cifra concreta, es el primer indicio de que la tecnología ha
pasado el codo de su curva.

\subsec{03. Primer modelo: la curva S de la cobertura de codificación}

Se toma como primera figura de mérito la cobertura de codificación, definida
en el Cuadro~\ref{tab:fom}, medida sobre la frontera: la fracción de la
infraestructura de un servicio de producción representativo que la mejor
herramienta de cada año permite declarar \emph{y} reconciliar. Es la figura
más apropiada para una curva S porque tiene un límite físico evidente, el
100 \%, que es justamente lo que un ajuste logístico necesita para estar
determinado.

Se ajusta por mínimos cuadrados la función logística

\[
y(t) = \frac{L}{1 + e^{-k\,(t - t_0)}},
\]

con $L$ acotado a 100 \% por construcción, sobre las doce anclas del
Cuadro~\ref{tab:hitos}. El resultado es
$L = 100{,}0\ \text{\pct}$, $k = 0{,}171\ \text{año}^{-1}$ y
$t_0 = 2012{,}0$, con un error cuadrático medio de 2,9 puntos porcentuales y
un coeficiente de determinación de 0,991. La Figura~\ref{fig:curvas}
superpone el ajuste a las anclas.

\begin{figure}[ht]
  \centering
  \begin{tikzpicture}
    \begin{axis}[
      width=15.2cm, height=7.4cm,
      xmin=1992, xmax=2034, ymin=0, ymax=104,
      xlabel={\footnotesize Año},
      ylabel={\footnotesize Cobertura de codificación alcanzable (\pct)},
      xlabel style={font=\footnotesize}, ylabel style={font=\footnotesize},
      xtick={1995,2000,2005,2010,2015,2020,2025,2030},
      ytick={0,20,40,60,80,100},
      xticklabel style={font=\scriptsize}, yticklabel style={font=\scriptsize},
      ymajorgrids, xmajorgrids, grid style={gray!18},
      axis lines=left,
      legend style={font=\scriptsize, at={(0.97,0.05)}, anchor=south east,
                    draw=gray!45, fill=white},
    ]
      % límite superior
      \addplot[gray!55, dashed, thick, forget plot]
        coordinates {(1992,100) (2034,100)};
      % ajuste logístico
      \addplot[capa, very thick, smooth] coordinates {
        (1993,3.8) (1994,4.5) (1995,5.2) (1996,6.2) (1997,7.2) (1998,8.4)
        (1999,9.9) (2000,11.5) (2001,13.3) (2002,15.4) (2003,17.8)
        (2004,20.4) (2005,23.3) (2006,26.5) (2007,30.0) (2008,33.7)
        (2009,37.6) (2010,41.7) (2011,45.9) (2012,50.1) (2013,54.4)
        (2014,58.6) (2015,62.7) (2016,66.6) (2017,70.2) (2018,73.7)
        (2019,76.9) (2020,79.8) (2021,82.4) (2022,84.7) (2023,86.8)
        (2024,88.6) (2025,90.2) (2026,91.6) (2027,92.9) (2028,93.9)
        (2029,94.8) (2030,95.6) (2031,96.3) (2032,96.8) (2033,97.3)
        (2034,97.7)};
      \addlegendentry{Ajuste logístico ($L=100$, $k=0{,}171$, $t_0=2012{,}0$)}
      % anclas observadas
      \addplot[only marks, mark=*, mark size=1.9pt, acento] coordinates {
        (1993,8) (1998,14) (2002,18) (2006,25) (2009,34) (2011,45)
        (2014,58) (2017,68) (2019,78) (2021,85) (2023,90) (2026,93)};
      \addlegendentry{Anclas del Cuadro \ref{tab:hitos}}
      % inflexión
      \draw[trans, thick, densely dotted] (axis cs:2012,0) -- (axis cs:2012,50.1);
      \node[trans, circle, fill, inner sep=1.5pt] at (axis cs:2012,50.1) {};
      \draw[trans, thin] (axis cs:2012,50.1) -- (axis cs:2008.6,66);
      \node[font=\scriptsize, text=trans, anchor=south east, align=right]
        at (axis cs:2008.8,65) {Punto de inflexión\\$t_0 = 2012$};
      % posición actual
      \draw[acento, thick, densely dotted] (axis cs:2026,0) -- (axis cs:2026,93);
      \node[font=\scriptsize, text=acento, anchor=east, align=right]
        at (axis cs:2025.4,64) {Hoy\\92 \pct{} del límite};
      \node[font=\scriptsize, text=gray!45!black, anchor=north east]
        at (axis cs:2033.6,99.2) {$L = 100$ \pct};
    \end{axis}
  \end{tikzpicture}
  \caption{Curva S de la cobertura de codificación alcanzable. El punto de
  inflexión queda en 2012, catorce años antes del inicio de esta hoja de
  ruta, y el 90 \pct{} del límite se alcanzó en 2025. Elaboración propia
  sobre las anclas del Cuadro~\ref{tab:hitos}.}
  \label{fig:curvas}
\end{figure}

Tres lecturas, en orden de importancia decreciente.

\textbf{El punto de inflexión ya pasó, y hace tiempo.} El ajuste lo sitúa en
2012, esto es, catorce años antes del arranque de esta hoja de ruta. La
pendiente máxima fue de 4,3 puntos porcentuales por año y hoy es de 1,3, la
tercera parte. La tecnología no está en la fase de crecimiento explosivo que
justificaría apostar por una plataforma emergente: está en el tercio superior
de su curva.

\textbf{Queda poco margen y se sabe cuánto.} El ajuste da 91,6 \% para 2026,
94,8 \% para 2029 y 96,8 \% para 2032. Todo el avance de la frontera en los
seis años del plan es de tres puntos porcentuales. La empresa, en cambio,
está en el 55 \% (Cuadro~\ref{tab:fomcomp}): su retraso respecto a la
frontera es de 37 puntos, doce veces mayor que el avance esperable de la
frontera misma.

\textbf{De ahí se sigue el criterio de inversión de toda la hoja de ruta.}
Con la frontera casi detenida y un retraso de 37 puntos, el riesgo relevante
no es apostar por la tecnología equivocada sino no cerrar la distancia con la
tecnología ya disponible. Esto explica por qué la sección~\ref{sec:comp}
concluye que la elección de motor aporta 1,9 de los 46,9 puntos de mejora:
el motor es precisamente la parte de la tecnología que está saturada.

\subsec{04. La misma curva sobre el tiempo de reconstrucción}

La cobertura tiene un techo aritmético, de modo que su saturación podría ser
un artefacto de la definición y no una propiedad de la tecnología. Conviene
por eso repetir el ejercicio sobre una figura de mérito sin techo evidente: el
tiempo de reconstrucción de un entorno completo desde el repositorio, que la
sección~\ref{sec:fom} considera la más reveladora de las ocho. Se ajusta la
logística decreciente

\[
y(t) = L_{\min} + \frac{Y_0 - L_{\min}}{1 + e^{k\,(t - t_0)}},
\]

sin acotar ninguno de sus cuatro parámetros. El ajuste devuelve una asíntota
inferior $L_{\min} = 3{,}6$ horas, $Y_0 = 1141$ horas,
$k = 0{,}211\ \text{año}^{-1}$ y $t_0 = 2000{,}5$, con un error cuadrático
medio de 2,7 horas sobre once anclas y un coeficiente de determinación de
0,9998. La Figura~\ref{fig:recons} lo representa en escala logarítmica, donde
la logística aparece como una recta que se dobla al final.

\begin{figure}[ht]
  \centering
  \begin{tikzpicture}
    \begin{semilogyaxis}[
      width=15.2cm, height=6.8cm,
      xmin=1996, xmax=2034, ymin=2.6, ymax=1400,
      xlabel={\footnotesize Año},
      ylabel={\footnotesize Tiempo de reconstrucción (horas)},
      xlabel style={font=\footnotesize}, ylabel style={font=\footnotesize},
      xtick={2000,2005,2010,2015,2020,2025,2030},
      ytick={4,10,40,100,400,1000},
      yticklabels={4,10,40,100,400,1000},
      xticklabel style={font=\scriptsize}, yticklabel style={font=\scriptsize},
      ymajorgrids, xmajorgrids, grid style={gray!18},
      axis lines=left, log ticks with fixed point,
      legend style={font=\scriptsize, at={(0.98,0.96)}, anchor=north east,
                    draw=gray!45, fill=white},
    ]
      \addplot[gray!55, dashed, thick, forget plot]
        coordinates {(1996,3.55) (2034,3.55)};
      \addplot[capa, very thick, smooth] coordinates {
        (1996,823.8) (1998,719.2) (2000,602.6) (2002,483.3) (2004,371.6)
        (2006,275.1) (2008,197.6) (2010,138.7) (2012,95.9) (2014,65.9)
        (2016,45.2) (2018,31.2) (2020,21.8) (2022,15.6) (2024,11.5)
        (2026,8.8) (2028,7.0) (2030,5.8) (2032,5.0) (2034,4.5)};
      \addlegendentry{Ajuste logístico decreciente ($L_{\min}=3{,}6$ h)}
      \addplot[only marks, mark=*, mark size=1.9pt, acento] coordinates {
        (1998,720) (2002,480) (2006,280) (2009,168) (2011,110) (2014,64)
        (2017,40) (2019,26) (2021,18) (2023,13) (2026,10)};
      \addlegendentry{Anclas de frontera}
      \node[font=\scriptsize, text=gray!45!black, anchor=south east]
        at (axis cs:2033.6,3.7) {Asíntota $L_{\min} = 3{,}6$ h};
      \draw[acento, thick, densely dotted]
        (axis cs:2026,2.6) -- (axis cs:2026,10);
      \node[font=\scriptsize, text=acento, anchor=south west, align=left]
        at (axis cs:2021.2,120) {Hoy: 10 h,\\2,8 veces la asíntota};
    \end{semilogyaxis}
  \end{tikzpicture}
  \caption{Tiempo de reconstrucción de un entorno completo, en escala
  logarítmica. El aplanamiento de la recta a partir de 2020 es la firma del
  techo de la curva S: el tiempo de división por dos se alarga de 4,3 a 5,1
  años. Elaboración propia.}
  \label{fig:recons}
\end{figure}

El diagnóstico coincide con el de la cobertura, y por una vía independiente.
El tiempo que la frontera tarda en dividir por dos esta figura era de 4,3 años
en el intervalo 2011--2026 y es de 5,1 años en 2019--2026: se \emph{alarga},
que es exactamente la firma del techo de una curva S
\parencite{christensen1992}. Y la asíntota estimada, 3,6 horas, deja claro
cuál es el premio máximo restante: desde las 10 horas de hoy no quedan órdenes
de magnitud por ganar, sino un factor de 2,8.

\subsec{05. Segundo modelo: la tasa de mejora al modo de la ley de Moore}

Los dos ajustes anteriores describen razones y tiempos, magnitudes acotadas
por su propia naturaleza. Un tercer indicador no lo está: el número absoluto
de tipos de recurso que la tecnología permite declarar, esto es, la superficie
declarable. Es la magnitud análoga al recuento de componentes por circuito
sobre el que se formuló la ley de Moore \parencite{moore1965}, y la que la
literatura sobre tendencias cuantitativas de rendimiento técnico encuentra
más frecuentemente en régimen exponencial \parencite{koh2006,magee2016}.

Se dispone de tres mediciones publicadas por sus responsables, que son las
únicas cifras auditables de este apartado:

\begin{enumerate}
  \item Febrero de 2011: el primer motor declarativo de plantilla completa
    cubre 48 tipos de recurso, correspondientes a 13 de los 15 servicios
    entonces existentes \parencite{cloudformationlaunch}.
  \item Marzo de 2023: el registro público del motor independiente del
    proveedor supera los 3.000 proveedores de recursos y los 12.000 módulos,
    con más de 250 socios \parencite{tfecosystem}.
  \item 2025: el mismo registro publica 7.207 proveedores y 24.384 módulos.
\end{enumerate}

Sobre los dos últimos, que son homogéneos entre sí, el ajuste log-lineal da un
crecimiento del 47,4 \% anual en proveedores y del 36,9 \% anual en módulos,
esto es, tiempos de duplicación de 1,79 y 2,21 años respectivamente. El
Cuadro~\ref{tab:tasas} recoge el resultado junto con las tasas de las ocho
figuras de mérito, y la Figura~\ref{fig:moore} lo representa.

\begin{figure}[ht]
  \centering
  \begin{tikzpicture}
    \begin{semilogyaxis}[
      width=15.2cm, height=6.4cm,
      xmin=2010, xmax=2032, ymin=30, ymax=90000,
      xlabel={\footnotesize Año},
      ylabel={\footnotesize Superficie declarable (recuento)},
      xlabel style={font=\footnotesize}, ylabel style={font=\footnotesize},
      xtick={2011,2014,2017,2020,2023,2026,2029,2032},
      ytick={100,1000,10000,100000},
      yticklabels={$10^2$,$10^3$,$10^4$,$10^5$},
      xticklabel style={font=\scriptsize}, yticklabel style={font=\scriptsize},
      ymajorgrids, xmajorgrids, grid style={gray!18},
      axis lines=left,
      legend style={font=\scriptsize, at={(0.02,0.96)}, anchor=north west,
                    draw=gray!45, fill=white},
    ]
      % recta de duplicación cada 2 años desde 48 en 2011
      \addplot[gray!60, dashed, thick, domain=2011:2032, samples=40]
        {48*2^((x-2011.15)/2)};
      \addlegendentry{Referencia: duplicación cada 2 años}
      % proveedores: 3000 (2023,24) -> 7207 (2025,5), Td = 1,79
      \addplot[jugA, very thick, domain=2023.24:2025.5, samples=12]
        {3000*2^((x-2023.24)/1.79)};
      \addlegendentry{Proveedores de recursos ($T_2 = 1{,}79$ años)}
      % módulos: 12000 -> 24384, Td = 2,21
      \addplot[jugB, very thick, domain=2023.24:2025.5, samples=12]
        {12000*2^((x-2023.24)/2.21)};
      \addlegendentry{Módulos publicados ($T_2 = 2{,}21$ años)}
      % prolongación al horizonte del plan: extrapolación, no medida
      \addplot[jugA, thick, densely dotted, domain=2025.5:2032, samples=24,
        forget plot] {3000*2^((x-2023.24)/1.79)};
      \addplot[jugB, thick, densely dotted, domain=2025.5:2032, samples=24]
        {12000*2^((x-2023.24)/2.21)};
      \addlegendentry{Prolongación al horizonte del plan (extrapolación)}
      \addplot[only marks, mark=*, mark size=2.1pt, acento]
        coordinates {(2011.15,48) (2023.24,3000) (2025.5,7207)
                     (2023.24,12000) (2025.5,24384)};
      \addlegendentry{Cifras publicadas}
      \node[font=\scriptsize, anchor=west, text=acento]
        at (axis cs:2011.5,48) {48 tipos (2011)};
    \end{semilogyaxis}
  \end{tikzpicture}
  \caption{Superficie declarable en escala logarítmica. Las dos series
  homogéneas se duplican cada 1,8 y 2,2 años, un régimen indistinguible del
  de la ley de Moore, que la recta discontinua gris representa. El trazo
  continuo cubre el único intervalo con dos cifras publicadas y el punteado es
  extrapolación. Elaboración propia sobre \textcite{cloudformationlaunch} y
  \textcite{tfecosystem}.}
  \label{fig:moore}
\end{figure}

El resultado es notable por su coincidencia: un tiempo de duplicación de 1,8 a
2,2 años es indistinguible de los dos años de la formulación original de la
ley de Moore. Conviene sin embargo no leer más de lo que el dato sostiene. Un
proveedor de recursos añadido al registro no equivale a un transistor: la
calidad es heterogénea, hay duplicidades y el recuento incluye integraciones
marginales. Lo que la cifra sostiene con solidez es la dirección y el orden de
magnitud del ritmo, no su valor exacto.

\subsec{06. Tasa de mejora de cada figura de mérito}

El Cuadro~\ref{tab:tasas} aplica el mismo procedimiento a las ocho figuras de
mérito del Cuadro~\ref{tab:fom}, comparando el tiempo de duplicación (o de
división por dos, según el sentido de cada figura) en dos intervalos
consecutivos. La comparación entre ambos, y no su valor absoluto, es lo que
diagnostica la fase: si el segundo intervalo es más lento que el primero, la
figura desacelera.

\begin{table}[htbp]
  \centering
  \footnotesize
  \setlength{\tabcolsep}{4pt}
  \caption{Valores de frontera y tasa de mejora de las ocho figuras de mérito.
  $T_2$ es el tiempo, en años, que la frontera tarda en duplicar la figura o
  en dividirla por dos, según su sentido. Un $T_2$ mayor en el segundo
  intervalo indica desaceleración.}
  \label{tab:tasas}
  \begin{tabular}{@{}lrrrrrrl@{}}
    \toprule
    \textbf{Figura de mérito} & \textbf{2011} & \textbf{2019} &
      \textbf{2026} & \multicolumn{2}{c}{\textbf{$T_2$ (años)}} &
      \textbf{Mejora} & \textbf{Fase} \\
    \cmidrule(lr){5-6}
    & & & & \textbf{11--19} & \textbf{19--26} &
      \textbf{anual} & \\
    \midrule
    Tiempo de ciclo (h)          & 168 &  24 &   6 &  2,9 &  3,5 & 21,9 \pct & Desacelera \\
    Frecuencia (cambios/sem.)    &   2 &  25 &  90 &  2,2 &  3,8 & 20,1 \pct & Desacelera \\
    Tasa de fallo (\pct)         &  30 &  12 &   5 &  6,1 &  5,5 & 13,3 \pct & Estable \\
    Restablecimiento (min)       & 240 &  60 &  30 &  4,0 &  7,0 & 10,4 \pct & Desacelera \\
    Cobertura (\pct)             &  45 &  78 &  93 & 10,1 & 27,6 &  2,5 \pct & Saturada \\
    Tasa de deriva (rec./1000)   &  40 &  12 &   4 &  4,6 &  4,4 & 17,0 \pct & Estable \\
    Reconstrucción (h)           & 110 &  26 &  10 &  3,8 &  5,1 & 14,6 \pct & Desacelera \\
    Cumplimiento (\pct)          &  20 &  60 &  88 &  5,1 & 12,7 &  5,6 \pct & Saturada \\
    \midrule
    \multicolumn{8}{@{}l}{\textit{Para contraste, la magnitud no acotada del
      apartado 05:}} \\
    Superficie declarable (rec.) & 48 & -- & 7207 & -- & 1,8 & 47,4 \pct &
      Exponencial \\
    \bottomrule
  \end{tabular}
\end{table}

El cuadro ordena las ocho figuras en tres grupos, y el criterio que los separa
no es tecnológico sino aritmético.

\textbf{Las dos figuras expresadas en porcentaje están saturadas}, con tiempos
de duplicación de 27,6 y 12,7 años. No podía ser de otro modo: una magnitud
acotada por el 100 \% tiene que saturar, y su saturación no informa sobre el
estado de la tecnología sino sobre la definición de la métrica. Es la razón
por la que este apartado no se detiene en ellas.

\textbf{Las cuatro figuras de tiempo desaceleran, pero siguen mejorando
rápido}, entre el 10 y el 22 \% anual. Su tiempo de división por dos se alarga
en todos los casos, de 2,9 a 3,5 años en el tiempo de ciclo y de 4,0 a 7,0 en
el de restablecimiento. Aquí sí hay un diagnóstico tecnológico, y es el mismo
que dio la Figura~\ref{fig:recons}.

\textbf{Sólo la magnitud no acotada permanece en régimen exponencial estable},
duplicándose cada 1,8 años. La conclusión conjunta es la del apartado 08.

\subsec{07. Tercer modelo: el frente de Pareto}

Los dos modelos anteriores describen la evolución en el tiempo. El tercero
describe la eficiencia en un instante: dadas las posiciones alcanzables en
2029, cuáles no están dominadas por ninguna otra. Se aplica sobre las dos
coordenadas ya construidas en el apartado 03 de la sección~\ref{sec:comp},
$\Delta V$ y $\Delta C$, con el criterio de dominancia natural del problema,
maximizar valor y minimizar coste: una posición está dominada si existe otra
con coste igual o menor y valor igual o mayor, y estrictamente mejor en al
menos uno de los dos.

Aplicado a las seis posiciones de los Cuadros~\ref{tab:fomcomp}
y~\ref{tab:coste}, el resultado es el del Cuadro~\ref{tab:pareto} y la
Figura~\ref{fig:pareto}.

\begin{table}[htbp]
  \centering
  \footnotesize
  \setlength{\tabcolsep}{5pt}
  \caption{Frente de Pareto de las seis posiciones, ordenadas por coste. La
  tasa de canje es el valor que se gana por cada punto de coste añadido al
  pasar al punto siguiente del frente.}
  \label{tab:pareto}
  \begin{tabular}{@{}lrrlr@{}}
    \toprule
    \textbf{Posición} & \textbf{$\Delta C$} & \textbf{$\Delta V$} &
      \textbf{Situación} & \textbf{Canje al siguiente} \\
    \midrule
    D. Contención de coste  & $-26$ & $-16{,}6$ & En el frente & 25,00 \\
    B. Bicep sobre Azure    & $-24$ & $+33{,}4$ & En el frente & 0,68 \\
    \textbf{Ruta propia}    & $-4$  & $+46{,}9$ & \textbf{En el frente} &
      \textbf{0,07} \\
    Producto de referencia  & $0$   & $0{,}0$   &
      Dominado por la ruta propia & -- \\
    A. Terraform gestionado & $+32$ & $+45{,}5$ &
      Dominado por la ruta propia & -- \\
    C. Pulumi               & $+44$ & $+50{,}3$ & En el frente & -- \\
    \bottomrule
  \end{tabular}
\end{table}

\begin{figure}[ht]
  \centering
  \begin{tikzpicture}
    \begin{axis}[
      width=14.6cm, height=8.2cm,
      xmin=-34, xmax=54, ymin=-26, ymax=60,
      xlabel={\footnotesize $\Delta C$: variación del coste anual de propiedad
              (puntos de índice)},
      ylabel={\footnotesize $\Delta V$: variación de valor para el cliente
              (puntos de utilidad)},
      xlabel style={font=\footnotesize}, ylabel style={font=\footnotesize},
      xtick={-30,-20,-10,0,10,20,30,40,50},
      ytick={-20,-10,0,10,20,30,40,50},
      xticklabel style={font=\scriptsize}, yticklabel style={font=\scriptsize},
      ymajorgrids, xmajorgrids, grid style={gray!18},
      axis lines=middle, axis on top,
      every axis x label/.style={at={(axis description cs:0.5,-0.14)}},
      every axis y label/.style={at={(axis description cs:-0.07,0.5)},rotate=90},
    ]
      % región dominada por la ruta propia
      \addplot[fill=capa!7, draw=none, forget plot]
        coordinates {(-4,46.9) (54,46.9) (54,-26) (-4,-26)} \closedcycle;
      % frente de Pareto
      \addplot[gray!55!black, very thick, dashed, forget plot]
        coordinates {(-26,-16.6) (-24,33.4) (-4,46.9) (44,50.3)};
      % puntos del frente
      \addplot[only marks, mark=*, mark size=2.9pt, descart, forget plot]
        coordinates {(-26,-16.6)};
      \addplot[only marks, mark=*, mark size=2.9pt, jugB, forget plot]
        coordinates {(-24,33.4)};
      \addplot[only marks, mark=*, mark size=3.6pt, propia, forget plot]
        coordinates {(-4,46.9)};
      \addplot[only marks, mark=*, mark size=2.9pt, jugC, forget plot]
        coordinates {(44,50.3)};
      % puntos dominados
      \addplot[only marks, mark=o, mark size=2.9pt, jugA, thick, forget plot]
        coordinates {(32,45.5)};
      \addplot[only marks, mark=o, mark size=2.9pt, gray!55!black, thick,
        forget plot] coordinates {(0,0)};
      % etiquetas
      \node[font=\scriptsize, anchor=west, text=descart, align=left]
        at (axis cs:-24.4,-19.5) {D. Contención};
      \node[font=\scriptsize, anchor=east, text=jugB, align=right]
        at (axis cs:-25.4,33.4) {B. Bicep};
      \node[font=\scriptsize\bfseries, anchor=north west, text=propia,
        align=left] at (axis cs:2.6,44.2) {Ruta propia};
      \node[font=\scriptsize, anchor=south east, text=jugC, align=right]
        at (axis cs:42.6,51.6) {C. Pulumi};
      \node[font=\scriptsize, anchor=west, text=jugA, align=left]
        at (axis cs:33.6,41.2) {A. Terraform\\\itshape dominado};
      \node[font=\scriptsize, anchor=north west, text=gray!45!black,
        align=left] at (axis cs:1.4,-1.4) {Referencia 2026\\\itshape dominada};
      \node[font=\scriptsize, text=capa!75!black, anchor=south east, align=right]
        at (axis cs:52,-24.4) {Región dominada\\por la ruta propia};
      % tasas de canje entre puntos consecutivos del frente
      \node[font=\scriptsize\itshape, text=gray!40!black, anchor=north]
        at (axis cs:-15,38.6) {0,68};
      \node[font=\scriptsize\itshape, text=gray!40!black, anchor=south]
        at (axis cs:16,50.6) {0,07 puntos de valor por punto de coste};
      \node[font=\scriptsize\itshape, text=gray!40!black, anchor=west]
        at (axis cs:-23.2,4) {25,00};
    \end{axis}
  \end{tikzpicture}
  \caption{Frente de Pareto de las posiciones alcanzables en 2029. Los puntos
  llenos están en el frente; los huecos están dominados. El sombreado marca la
  región dominada por la ruta propia, que contiene a Terraform gestionado y a
  la plataforma actual. Elaboración propia.}
  \label{fig:pareto}
\end{figure}

El frente aporta dos resultados que el gráfico vectorial de la
sección~\ref{sec:comp} no hacía explícitos.

\textbf{Dos de las seis posiciones están dominadas, y una de ellas es un
competidor.} Terraform gestionado ofrece 45,5 puntos de valor por 32 puntos
de coste añadido; la ruta propia ofrece 46,9 puntos por 4 puntos \emph{menos}
de coste. Es dominancia estricta en ambos ejes: no existe ninguna
ponderación de los dos objetivos, ninguna preferencia sobre el canje entre
valor y coste, que haga preferible el jugador A a la ruta propia. Es una
conclusión más fuerte que la del gráfico vectorial, que sólo mostraba que la
ruta propia quedaba mejor situada. La plataforma actual está dominada por la
misma razón, lo que equivale a decir que no actuar no es una opción eficiente.

\textbf{La ruta propia está en el codo del frente.} Las tasas de canje del
Cuadro~\ref{tab:pareto} caen en dos órdenes de magnitud a lo largo del
frente: pasar de la contención de coste a Bicep compra 25 puntos de valor por
punto de coste; pasar de Bicep a la ruta propia, 0,68; y pasar de la ruta
propia a Pulumi, 0,07. Los últimos 3,4 puntos de valor disponibles cuestan 48
puntos de coste, catorce veces más caros que los anteriores. La ruta propia es
el último punto del frente antes de que el rendimiento marginal se desplome.

Conviene rectificar aquí una afirmación del apartado 08 de la
sección~\ref{sec:comp}, que el frente permite precisar. Bajo el escenario
multinube, el jugador B se desplaza a $(-4; +10{,}9)$ y la ruta propia a
$(-1; +45{,}0)$: B no queda \emph{dominado}, porque conserva tres puntos de
ventaja en coste. Lo que ocurre es más elocuente que la dominancia. B sigue
formalmente en el frente, pero la tasa de canje entre B y la ruta propia salta
de 0,68 a 11,37, es decir, esos tres puntos de coste ahorrado cuestan 34,1
puntos de valor. B pasa de ser el punto barato y razonable del frente a ser
un extremo degenerado del que ninguna función de preferencia plausible
escogería.

\subsec{08. Dónde está la tecnología en su curva}

Los tres modelos, aplicados a las mismas figuras de mérito, dan un diagnóstico
único y en dos regímenes simultáneos, que es la conclusión de esta sección.

\textbf{Lo que la tecnología entrega, medido como razón o como tiempo, está
saturando.} La cobertura alcanzable pasó su inflexión en 2012 y ha recorrido
el 92 \% de su límite. El tiempo de reconstrucción se acerca a una asíntota de
3,6 horas y su tiempo de división por dos se alarga de 4,3 a 5,1 años. Las
cuatro figuras de entrega desaceleran todas. Ninguna de las ocho figuras de
mérito promete otro orden de magnitud.

\textbf{Lo que la tecnología abarca, medido como recuento, sigue en régimen
exponencial.} La superficie declarable se duplica cada 1,8 a 2,2 años, sin
signo de desaceleración.

Las dos afirmaciones no se contradicen, y entender por qué no lo hacen es
entender el estado de la tecnología. La superficie declarable crece
exponencialmente, pero la superficie \emph{a declarar} crece al mismo ritmo,
porque son la misma cosa vista desde los dos lados: cada servicio nuevo del
proveedor aparece casi a la vez como recurso declarable y como recurso que hay
que declarar. La razón entre ambas, que es la cobertura, se mantiene por
tanto casi constante y satura. La tecnología no se está quedando corta:
avanza deprisa en términos absolutos y en el sitio, en términos relativos.

De aquí se siguen tres consecuencias para la hoja de ruta, y las tres
reaparecen en la declaración de estrategia de la sección~\ref{sec:estrategia}.

\begin{enumerate}
  \item \textbf{No hay ventaja competitiva en la elección de herramienta.}
    Una tecnología en el tercio superior de su curva S no discrimina entre
    quienes la adoptan. Es coherente con el hallazgo de la
    sección~\ref{sec:comp}: el vector que sustituye el motor aporta 1,9 de
    los 46,9 puntos de mejora.
  \item \textbf{La ventaja está en cerrar el retraso propio.} 37 puntos de
    cobertura separan a la empresa de la frontera, doce veces el avance
    esperable de la frontera en el mismo periodo. El competidor relevante es
    la propia plataforma de 2026, no Terraform ni Pulumi.
  \item \textbf{El régimen exponencial de la superficie declarable convierte
    la cobertura en un objetivo móvil.} Si el parque crece y la superficie
    declarable se duplica cada dos años, mantener la cobertura al 86 \% exige
    trabajo recurrente y no un proyecto que termina. Esto es lo que justifica
    que R1 tenga continuidad operativa después de 2029 en la
    sección~\ref{sec:portfolio}, en lugar de cerrarse.
\end{enumerate}

\subsec{09. Límites de los tres ajustes}

Un ajuste con un coeficiente de determinación de 0,99 sobre anclas escogidas
por el propio autor transmite más certidumbre de la que merece. Cuatro
reservas, en orden de gravedad.

\begin{enumerate}
  \item \textbf{Las anclas de las cuatro figuras específicas de
    infraestructura son estimaciones, no medidas.} La cobertura, la deriva, la
    reconstrucción y el cumplimiento no las mide ninguna fuente pública de
    forma sistemática, y sus valores de frontera proceden de juicio de
    ingeniería sobre las capacidades documentadas de cada hito. El ajuste
    logístico de la Figura~\ref{fig:curvas} tiene, por eso, el valor de una
    hipótesis ordenada, no de una regresión sobre datos. Las cuatro figuras de
    entrega están en mejor situación, porque los informes anuales citados
    publican percentiles comparables \parencite{forsgren2018,dora2024}.
  \item \textbf{Un ajuste logístico siempre encuentra una saturación.} La
    forma funcional la impone por construcción, de modo que el ajuste no puede
    usarse como prueba de que la tecnología satura. Lo que sí es informativo,
    y no depende de la forma funcional elegida, es el alargamiento del tiempo
    de división por dos entre dos intervalos consecutivos, que el
    Cuadro~\ref{tab:tasas} calcula sin ajustar ninguna curva. Ése es el
    verdadero soporte del diagnóstico.
  \item \textbf{La tercera cifra del apartado 05 no está publicada por su
    responsable.} Los recuentos de 2011 y de marzo de 2023 proceden de
    comunicaciones oficiales; el de 2025 procede de una fuente secundaria. El
    tiempo de duplicación descansa sobre dos puntos, uno de ellos de calidad
    inferior, y bastaría un error del 20 \% en el segundo para moverlo de 1,8
    a 2,2 años. La conclusión sobrevive a ese error, porque el intervalo
    resultante sigue siendo el de la ley de Moore, pero la cifra puntual no
    debe citarse aislada.
  \item \textbf{El frente de Pareto hereda los supuestos del índice de
    valor.} Sus coordenadas son las del Cuadro~\ref{tab:fomcomp}, construidas
    con los pesos declarados en la sección~\ref{sec:comp}. La dominancia de
    Terraform gestionado por la ruta propia es, en cambio, robusta a los
    pesos: se da en ambos ejes a la vez, de modo que ninguna reponderación de
    las cuatro figuras la revierte mientras se conserven las estimaciones de
    partida.
\end{enumerate}


% ==========================================================================
\section{Correspondencia con las motivaciones estratégicas}
\label{sec:motiv}
% ==========================================================================

Este es el elemento 5 del índice, y es el que decide si la hoja de ruta es un
plan de empresa o un ejercicio de ingeniería. Su función es doble: comprobar
que cada figura de mérito responde a una motivación declarada de la
organización, y comprobar en sentido inverso que ninguna motivación declarada
se queda sin figura de mérito que la vigile. Las dos comprobaciones fallan de
manera distinta y ambas son frecuentes: la primera produce métricas que nadie
usa para decidir, y la segunda, objetivos de empresa que el plan técnico no
puede demostrar que atiende.

\subsec{01. Las seis motivaciones estratégicas}

El Cuadro~\ref{tab:motiv} enuncia las seis motivaciones que la dirección
sostiene para la plataforma. Están redactadas de forma deliberadamente
incómoda: cada una incluye la condición que la haría fracasar, porque una
motivación que no puede fracasar no orienta ninguna decisión.

\begin{table}[htbp]
  \centering
  \footnotesize
  \caption{Las seis motivaciones estratégicas de la empresa para la
  plataforma de infraestructura como código, y el fracaso que cada una define.}
  \label{tab:motiv}
  \begin{tabularx}{\textwidth}{@{}l>{\raggedright\arraybackslash}p{0.235\textwidth}
    >{\raggedright\arraybackslash}X@{}}
    \toprule
    \textbf{Cód.} & \textbf{Motivación} & \textbf{Qué constituiría un fracaso} \\
    \midrule
    E1 & Plazo de salida al mercado & Que el aprovisionamiento de
      infraestructura siga siendo el camino crítico de un lanzamiento de
      producto. \\
    E2 & Escalabilidad del coste de operación & Que el equipo de plataforma
      tenga que crecer en la misma proporción que el parque de recursos que
      gobierna. \\
    E3 & Cumplimiento normativo demostrable & Que cada auditoría exija un
      esfuerzo extraordinario de recopilación de pruebas en lugar de
      resolverse con la evidencia que el sistema genera por sí solo. \\
    E4 & Continuidad ante desastre & Que la recuperación de una región
      completa dependa de conocimiento no escrito o de personas concretas. \\
    E5 & Capacidad de negociación con el proveedor & Que un cambio de
      condiciones comerciales, de licencia o de precio no admita respuesta
      técnica en un plazo razonable. \\
    E6 & Atracción y retención de ingeniería & Que la operación manual
      recurrente consuma la parte del trabajo que hace atractivos los puestos
      de plataforma. \\
    \bottomrule
  \end{tabularx}
\end{table}

\subsec{02. Correspondencia con las figuras de mérito}

El Cuadro~\ref{tab:motivfom} cruza las seis motivaciones con las ocho figuras
de mérito del Cuadro~\ref{tab:fom}. La marca llena indica que la figura es un
indicador \emph{principal} de la motivación, esto es, que un cambio apreciable
en la motivación se manifiesta necesariamente en esa figura; la marca hueca
indica que la figura es un indicador secundario, sensible pero no suficiente.

\begin{table}[htbp]
  \centering
  \footnotesize
  \setlength{\tabcolsep}{4pt}
  \caption{Correspondencia entre las motivaciones estratégicas y las figuras
  de mérito. $\bullet$ indicador principal, $\circ$ indicador secundario. La
  fila inferior cuenta los indicadores principales de cada figura y la columna
  derecha, los de cada motivación.}
  \label{tab:motivfom}
  \begin{tabular}{@{}l*{8}{c}r@{}}
    \toprule
    & \multicolumn{8}{c}{\textbf{Figuras de mérito}} & \\
    \cmidrule(lr){2-9}
    \textbf{Motivación} & \rotatebox{60}{Tiempo de ciclo} &
      \rotatebox{60}{Frecuencia} & \rotatebox{60}{Tasa de fallo} &
      \rotatebox{60}{Restablecimiento} & \rotatebox{60}{Cobertura} &
      \rotatebox{60}{Tasa de deriva} & \rotatebox{60}{Reconstrucción} &
      \rotatebox{60}{Cumplimiento} & \textbf{Ppal.} \\
    \midrule
    E1. Plazo de salida     & $\bullet$ & $\bullet$ & $\circ$   &           & $\bullet$ &           &           &           & 3 \\
    E2. Coste escalable     & $\circ$   & $\circ$   & $\bullet$ &           & $\bullet$ & $\bullet$ &           & $\circ$   & 3 \\
    E3. Cumplimiento        &           &           & $\circ$   &           & $\bullet$ & $\bullet$ &           & $\bullet$ & 3 \\
    E4. Continuidad         &           &           & $\circ$   & $\bullet$ & $\bullet$ & $\circ$   & $\bullet$ &           & 3 \\
    E5. Negociación         &           &           &           &           & $\bullet$ &           & $\bullet$ &           & 2 \\
    E6. Retención           & $\circ$   & $\circ$   &           & $\bullet$ &           & $\bullet$ & $\circ$   &           & 2 \\
    \midrule
    \textbf{Principales}    & 1 & 1 & 1 & 2 & \textbf{5} & 3 & 3 & 1 & \\
    \bottomrule
  \end{tabular}
\end{table}

Las dos comprobaciones que motivan la sección se resuelven leyendo el cuadro
por sus dos márgenes, y las dos salen bien, pero con una asimetría que merece
comentario.

\textbf{Ninguna fila está vacía}, de modo que las seis motivaciones tienen al
menos dos indicadores principales. Ninguna columna está vacía tampoco: las
ocho figuras de mérito sirven a alguna motivación declarada y ninguna es
métrica por sí misma.

\textbf{La cobertura de codificación es indicador principal de cinco de las
seis motivaciones}, y ninguna otra figura pasa de tres. Es un resultado que no
se buscaba y que atraviesa todo el documento: la misma figura que la
sección~\ref{sec:evol} identifica como la que más satura en la frontera y la
que más separa a la empresa de ella (37 puntos de retraso), y la que el
análisis de sensibilidad del modelo financiero de la
sección~\ref{sec:financiero} señala como la variable que más mueve el valor
actual neto, es la que conecta con más motivaciones de la dirección. La
convergencia de cuatro análisis independientes sobre la misma figura es el
argumento más sólido que la hoja de ruta puede ofrecer para su prioridad.

\textbf{La única motivación con sólo dos indicadores principales y ninguno
propio es E5, la capacidad de negociación.} Comparte cobertura y
reconstrucción con otras motivaciones y no tiene ninguna figura que la vigile
en exclusiva. Es una debilidad reconocida del cuadro de mando: la evaluación
de madurez de la sección~\ref{sec:madurez} la recoge como tal, y su
consecuencia práctica se discute en el apartado 04.

\subsec{03. Trazabilidad completa hasta los proyectos}

La correspondencia con las figuras de mérito no basta. Una motivación
estratégica sólo está atendida si existe un proyecto de la cartera cuyo
resultado mueve la figura que la vigila. El Cuadro~\ref{tab:trazafom} cierra
la cadena en sus cuatro eslabones: motivación, figura de mérito, salto
comprometido para 2029 y proyecto responsable.

\begin{table}[htbp]
  \centering
  \footnotesize
  \setlength{\tabcolsep}{4pt}
  \caption{Cadena de trazabilidad de la hoja de ruta: de la motivación
  estratégica al proyecto que la ejecuta. Los códigos de proyecto son los del
  Cuadro~\ref{tab:dsmproy} y los vectores, los del Cuadro~\ref{tab:vectores}.
  El proyecto en negrita es el responsable principal del salto.}
  \label{tab:trazafom}
  \begin{tabular}{@{}llrrll@{}}
    \toprule
    \textbf{Motiv.} & \textbf{Figura de mérito} & \textbf{2026} &
      \textbf{2029} & \textbf{Proyectos} & \textbf{Vector} \\
    \midrule
    E1 & Tiempo de ciclo (h)             &  26 &   8 & \textbf{R7}, R3, R5 & P3, P1 \\
    E1 & Frecuencia (cambios/sem.)       &  12 &  40 & \textbf{R7}, R3     & P3, P1 \\
    E1 & Cobertura (\pct)                &  55 &  86 & \textbf{R1}, R3     & P1 \\
    \midrule
    E2 & Tasa de fallo (\pct)            &  14 &   6 & \textbf{R5}, R4     & P1 \\
    E2 & Cobertura (\pct)                &  55 &  86 & \textbf{R1}         & P1 \\
    E2 & Tasa de deriva (rec./1000 sem.) &  38 &   8 & \textbf{R8}, R1     & -- \\
    \midrule
    E3 & Cobertura (\pct)                &  55 &  86 & \textbf{R1}         & P1 \\
    E3 & Tasa de deriva (rec./1000 sem.) &  38 &   8 & \textbf{R8}         & -- \\
    E3 & Cumplimiento (\pct)             &  62 &  94 & \textbf{R4}, R5     & P1 \\
    \midrule
    E4 & Restablecimiento (min)          & 180 &  45 & \textbf{R8}, R2     & -- \\
    E4 & Cobertura (\pct)                &  55 &  86 & \textbf{R1}         & P1 \\
    E4 & Reconstrucción (h)              &  40 &  11 & \textbf{R1}, R2, R3 & P1 \\
    \midrule
    E5 & Cobertura (\pct)                &  55 &  86 & \textbf{R1}         & P1 \\
    E5 & Reconstrucción (h)              &  40 &  11 & \textbf{R3}, R6     & P1, P2 \\
    \midrule
    E6 & Restablecimiento (min)          & 180 &  45 & \textbf{R8}         & -- \\
    E6 & Tasa de deriva (rec./1000 sem.) &  38 &   8 & \textbf{R8}, R1     & -- \\
    \bottomrule
  \end{tabular}
\end{table}

Dos lecturas de la cadena completa.

\textbf{R1 y R8 aparecen como responsables principales en once de las
dieciséis filas, y ninguno de los dos desplaza por sí solo el gráfico
vectorial.} Son los dos proyectos que el Cuadro~\ref{tab:dsmproy} clasifica
como habilitadores, con la columna de vector en blanco. La cadena de
trazabilidad revela así una distorsión del elemento 6 que conviene tener
presente: el gráfico vectorial mide valor percibido por el equipo de producto,
y por eso los proyectos que sirven al cumplimiento, a la continuidad y a la
retención le resultan invisibles, aunque atiendan cuatro de las seis
motivaciones de la dirección.

\textbf{Ningún proyecto de la cartera queda sin motivación.} Los ocho
aparecen en el cuadro, lo que descarta el fallo más común de una cartera
técnica: el proyecto que se sostiene por su interés propio. El caso más
justificado de sospecha era R6, la sustitución del motor, que aporta 1,9 de
los 46,9 puntos de valor del gráfico vectorial; el cuadro muestra que su
justificación no está en el eje de valor sino en E5, la capacidad de
negociación, que ninguna otra figura de mérito recoge bien.

\subsec{04. Lo que esta correspondencia deja sin resolver}

Tres reservas, y las tres se trasladan a la evaluación de madurez de la
sección~\ref{sec:madurez}.

\begin{enumerate}
  \item \textbf{E5 no tiene un indicador propio.} La capacidad de negociación
    con el proveedor se manifiesta en el coste de una sustitución, no en
    ninguna de las ocho figuras de mérito. Una novena figura, el tiempo
    estimado de sustitución del motor manteniendo el servicio, la vigilaría
    mejor, pero no se ha incluido en el Cuadro~\ref{tab:fom} porque no cumple
    la primera condición del apartado 02 de la sección~\ref{sec:fom}: no es
    observable con un procedimiento único, sino sólo estimable. Se deja
    consignada la carencia en lugar de resolverla con una métrica que no
    resistiría el criterio.
  \item \textbf{Las motivaciones compiten entre sí y el cuadro no lo muestra.}
    E1 empuja a bajar la barrera de despliegue y E3 a subirla. La política
    como código (R4) es precisamente el mecanismo que permite atender a las
    dos a la vez, al trasladar el control del momento de la revisión humana al
    momento de la evaluación automática del plan, pero el
    Cuadro~\ref{tab:motivfom} no expresa esa tensión: sólo registra que ambas
    motivaciones tienen indicadores.
  \item \textbf{Las seis motivaciones son internas a la organización de
    tecnología.} No hay ninguna motivación de mercado, ni de ingreso, ni de
    diferenciación de producto, porque la infraestructura como código es una
    capacidad interna y no un producto de catálogo. Es coherente con el
    planteamiento del apartado 01 de la sección~\ref{sec:comp}, pero conviene
    decirlo: esta hoja de ruta justifica una inversión en eficiencia y en
    riesgo, no en crecimiento.
\end{enumerate}


% ==========================================================================
\section{Posicionamiento frente a la competencia}
\label{sec:comp}
% ==========================================================================

Esta sección corresponde al elemento 6 del índice de una hoja de ruta
tecnológica: el posicionamiento de la empresa con respecto a la competencia
expresado sobre las figuras de mérito de la tecnología
\parencite{deweck2022}. El instrumento empleado es el gráfico vectorial, que
sitúa a cada jugador en un plano cuyos ejes miden la variación de valor
percibida por el cliente y la variación de coste soportada por el productor, y
que descompone el desplazamiento de cada uno en los vectores tecnológicos
concretos que lo producen. Todo el análisis se apoya en las ocho figuras de
mérito definidas en el Cuadro~\ref{tab:fom}.

\subsec{01. Qué se compara y desde qué punto de vista}

Un gráfico vectorial obliga a fijar tres cosas antes de dibujar la primera
flecha, y ninguna de ellas es evidente cuando la tecnología analizada es una
capacidad interna y no un producto de catálogo.

\begin{enumerate}
  \item \textbf{El productor} es la organización de plataforma: el equipo que
    construye y opera el sistema de infraestructura como código descrito en
    las secciones anteriores. El coste que soporta es el coste anual de
    propiedad de ese sistema, no el coste de la infraestructura desplegada
    con él.
  \item \textbf{El cliente} es el equipo de producto que consume la
    plataforma. El valor que percibe no es abstracto: son las cuatro figuras
    de mérito del Cuadro~\ref{tab:fom} que se manifiestan en su trabajo
    diario, a saber, el tiempo de ciclo del cambio, la cobertura de
    codificación, la tasa de fallo del cambio y el tiempo de reconstrucción.
    Las otras cuatro son internas del productor y no entran en el eje de
    valor.
  \item \textbf{El producto de referencia} es la plataforma tal como está
    hoy, en 2026: motor de aprovisionamiento comunitario, estado remoto,
    canalización de integración continua, sin política como código, sin
    catálogo certificado y sin interfaz de autoservicio. Ocupa el origen de
    coordenadas por construcción, y todo el gráfico mide desviaciones
    respecto a él.
\end{enumerate}

Los jugadores comparados no son las empresas proveedoras entre sí, sino las
posiciones que alcanza una organización adoptante al construir su plataforma
sobre cada una de las tres familias tecnológicas al horizonte de 2029, es
decir, hoy más tres años. Esta distinción importa: la pregunta que responde el
gráfico no es cuál herramienta es mejor en abstracto, sino dónde queda situada
la empresa según cuál adopte.

\subsec{02. Los tres competidores considerados}

El Cuadro~\ref{tab:competidores} caracteriza a los tres jugadores. Se han
elegido porque cubren las tres estrategias de producto que hoy se disputan la
capa L1 de la Figura~\ref{fig:pila}: el estándar de hecho multinube
(Terraform), la integración vertical de un proveedor único (Bicep) y la
apuesta por los lenguajes de propósito general (Pulumi).

\begin{table}[htbp]
  \centering
  \footnotesize
  \setlength{\tabcolsep}{4pt}
  \caption{Caracterización de los tres jugadores en la capa de
  aprovisionamiento.}
  \label{tab:competidores}
  \begin{tabularx}{\textwidth}{@{}>{\raggedright\arraybackslash}p{0.15\textwidth}
    >{\raggedright\arraybackslash}X
    >{\raggedright\arraybackslash}X
    >{\raggedright\arraybackslash}X@{}}
    \toprule
    & \textbf{Jugador A: Terraform} & \textbf{Jugador B: Bicep} &
      \textbf{Jugador C: Pulumi} \\
    \midrule
    Titular & HashiCorp, integrada en IBM desde 2025 & Microsoft &
      Pulumi Corporation \\
    Lenguaje & HCL, lenguaje declarativo propio
      \parencite{terraformdocs} & DSL declarativo que se
      traduce a plantillas ARM & lenguajes de propósito general: TypeScript,
      Python, Go, C\# \\
    Ámbito & multinube, con un registro de proveedores muy extenso &
      exclusivamente Azure & multinube, reutiliza los proveedores de
      Terraform \\
    Estado & fichero de estado explícito, remoto, cifrado y con bloqueo
      \parencite{terraformstate} & sin fichero de estado: lo mantiene el
      gestor de recursos de Azure \parencite{deploymentstacks} & fichero de
      estado, en servicio gestionado o autoalojado \\
    Licencia & BUSL desde 2023 \parencite{hashicorpbusl}; servicio gestionado
      de pago & gratuito, incluido en la suscripción de Azure
      \parencite{bicepdocs} & núcleo abierto; servicio gestionado con tarifa
      por recurso \parencite{pulumidocs} \\
    Política como código & Sentinel u OPA \parencite{sentinel,opa} &
      Azure Policy, integrado en el plano de control &
      CrossGuard \\
    Reutilización & registro de módulos, público y privado & módulos
      verificados publicados por el propio proveedor \parencite{avm} &
      componentes que son clases del lenguaje anfitrión \\
    Barrera de adopción & expresividad limitada de HCL y coste de licencia al
      crecer & dependencia de un proveedor único & exige competencia real de
      programación en el equipo de plataforma \\
    \bottomrule
  \end{tabularx}
\end{table}

Junto a los tres jugadores se sitúa la ruta propuesta para la empresa, que no
coincide con ninguno de ellos porque combina un motor compatible de código
abierto \parencite{opentofu} con dos inversiones propias, el catálogo de
módulos certificados y la interfaz de autoservicio. El gráfico incluye además
una quinta posición, la contención de coste, que se evaluó y se descartó, y
cuya única función aquí es mostrar que el cuadrante inferior izquierdo del
plano existe y es alcanzable.

\subsec{03. Construcción de las dos coordenadas}

Los dos ejes del gráfico no son magnitudes observables sino índices
construidos, y conviene declarar cómo se calculan antes de leer ninguna
posición.

\textbf{Eje de ordenadas: variación de valor para el cliente.} Cada una de las
cuatro figuras de mérito orientadas al cliente se convierte en una utilidad
acotada entre 0 y 100 mediante interpolación lineal entre dos anclas, el peor
valor aceptable y el mejor valor alcanzable con la tecnología disponible, al
modo de la exploración de espacios de diseño multiatributo
\parencite{ross2004}. Para las figuras en las que menos es mejor,

\[
u_i = 100 \cdot \frac{x_i^{\text{peor}} - x_i}
                     {x_i^{\text{peor}} - x_i^{\text{mejor}}},
\]

y con el numerador invertido para la cobertura de codificación, que es la
única en la que más es mejor. Las anclas adoptadas son: tiempo de ciclo, de
48 a 4 horas; cobertura, del 40 al 95 \%; tasa de fallo, del 20 al 3 \%;
tiempo de reconstrucción, de 72 a 6 horas. El valor agregado es la media
ponderada

\[
V = 0{,}30\,u_{\text{ciclo}} + 0{,}25\,u_{\text{cobertura}}
  + 0{,}25\,u_{\text{fallo}} + 0{,}20\,u_{\text{reconstrucción}},
\]

y la coordenada del gráfico es $\Delta V = V - V_{\text{ref}}$, con
$V_{\text{ref}} = 40{,}3$ puntos para el producto de referencia. Los pesos
reflejan una decisión explícita: se prima la velocidad percibida sobre la
capacidad de reconstrucción, porque la primera se experimenta cada día y la
segunda sólo ante un incidente grave.

\textbf{Eje de abscisas: variación de coste para el productor.} Es un índice
de coste anual de propiedad por cada cien recursos gestionados, con la
plataforma actual normalizada a 100 puntos. Se descompone en tres sumandos:
licencia y servicio gestionado, ingeniería de mantenimiento del propio sistema
de plataforma, y formación e incorporación de personas. La coordenada del
gráfico es $\Delta C = C - 100$.

Una advertencia sobre el eje de coste: mide el régimen permanente en 2029, no
la inversión de transición. El coste de una sola vez de cada migración
pertenece al modelo financiero, que es el elemento 8 del índice y se desarrolla
en la sección~\ref{sec:financiero}. Un jugador puede parecer barato en el
gráfico y exigir un desembolso inicial considerable para llegar a esa
posición: la ruta propia, por ejemplo, aparece aquí con un $\Delta C$ de
$-4$ puntos y exige una exposición máxima de tesorería de 335,6 miles de
euros para alcanzarlo.

\subsec{04. Estimación de las figuras de mérito al horizonte 2029}

El Cuadro~\ref{tab:fomcomp} recoge las estimaciones y las utilidades
resultantes. Son estimaciones razonadas de ingeniería, no medidas: se apoyan
en el comportamiento observado de la plataforma actual y en las propiedades
declaradas de cada tecnología, y su función es ordenar posiciones, no predecir
cifras.

\begin{table}[htbp]
  \centering
  \footnotesize
  \setlength{\tabcolsep}{4.5pt}
  \caption{Figuras de mérito estimadas al horizonte 2029 y valor agregado
  resultante. Las cuatro figuras son las orientadas al cliente del
  Cuadro~\ref{tab:fom}.}
  \label{tab:fomcomp}
  \begin{tabular}{@{}lrrrrrr@{}}
    \toprule
    \textbf{Posición} & \textbf{Ciclo} & \textbf{Cobert.} &
      \textbf{Fallo} & \textbf{Reconstr.} & \textbf{$V$} &
      \textbf{$\Delta V$} \\
    & (h) & (\pct) & (\pct) & (h) & & \\
    \midrule
    Producto de referencia (2026) & 26 & 55 & 14 & 40 & 40,3 & 0,0 \\
    \midrule
    A. Terraform gestionado       &  9 & 85 &  6 & 12 & 85,8 & $+45{,}5$ \\
    B. Bicep sobre Azure          & 12 & 72 &  8 & 16 & 73,7 & $+33{,}4$ \\
    C. Pulumi                     &  7 & 88 &  5 & 10 & 90,6 & $+50{,}3$ \\
    \textbf{Ruta propia}          &  8 & 86 &  6 & 11 & 87,3 & $+46{,}9$ \\
    \midrule
    D. Contención de coste (descartada) & 32 & 44 & 15 & 60 & 23,7 &
      $-16{,}6$ \\
    \bottomrule
  \end{tabular}
\end{table}

El Cuadro~\ref{tab:coste} descompone el otro eje. La diferencia decisiva no
está en la licencia sino en la ingeniería de mantenimiento: un catálogo de
módulos certificados y una política como código reducen el trabajo recurrente
de revisión y corrección lo bastante para compensar el esfuerzo de
construirlos y sostenerlos.

\begin{table}[htbp]
  \centering
  \footnotesize
  \setlength{\tabcolsep}{5pt}
  \caption{Descomposición del índice de coste anual de propiedad por cada cien
  recursos gestionados. La plataforma actual se normaliza a 100 puntos.}
  \label{tab:coste}
  \begin{tabular}{@{}lrrrrr@{}}
    \toprule
    \textbf{Posición} & \textbf{Licencia y} & \textbf{Ingeniería de} &
      \textbf{Formación} & \textbf{Total} & \textbf{$\Delta C$} \\
    & \textbf{servicio} & \textbf{mantenimiento} & & & \\
    \midrule
    Producto de referencia (2026) &  8 & 76 & 16 & 100 & 0 \\
    \midrule
    A. Terraform gestionado       & 52 & 60 & 20 & 132 & $+32$ \\
    B. Bicep sobre Azure          &  0 & 54 & 22 &  76 & $-24$ \\
    C. Pulumi                     & 46 & 62 & 36 & 144 & $+44$ \\
    \textbf{Ruta propia}          & 12 & 66 & 18 &  96 & $-4$ \\
    \midrule
    D. Contención de coste (descartada) & 6 & 56 & 12 & 74 & $-26$ \\
    \bottomrule
  \end{tabular}
\end{table}

La Figura~\ref{fig:perfilfom} compara los perfiles de utilidad figura por
figura, antes de agregarlos. Muestra algo que la coordenada agregada esconde:
las cuatro posiciones avanzadas son casi indistinguibles en velocidad y en
fiabilidad, y toda su diferencia se concentra en la cobertura de codificación.
Bicep no pierde por ser lento, sino por dejar recursos fuera del alcance del
motor.

\begin{figure}[ht]
  \centering
  \begin{tikzpicture}
    \begin{axis}[
      width=15.2cm, height=6.4cm,
      ybar=1pt, bar width=7pt,
      ymin=0, ymax=104,
      ylabel={\footnotesize Utilidad (0 a 100)},
      ylabel style={font=\footnotesize},
      symbolic x coords={f1,f2,f3,f4},
      xtick={f1,f2,f3,f4},
      xticklabels={Tiempo de ciclo, Cobertura de codificación,
                   Tasa de fallo, Tiempo de reconstrucción},
      xticklabel style={font=\scriptsize},
      yticklabel style={font=\scriptsize},
      ytick={0,20,40,60,80,100},
      ymajorgrids, grid style={gray!22},
      axis lines=left,
      enlarge x limits=0.16,
      legend style={font=\scriptsize, at={(0.5,1.04)}, anchor=south,
                    legend columns=5, draw=none, column sep=1.1ex},
    ]
      \addplot[fill=gray!40,  draw=gray!70!black]
        coordinates {(f1,50.0) (f2,27.3) (f3,35.3) (f4,48.5)};
      \addplot[fill=jugA!35,  draw=jugA]
        coordinates {(f1,88.6) (f2,81.8) (f3,82.4) (f4,90.9)};
      \addplot[fill=jugB!35,  draw=jugB]
        coordinates {(f1,81.8) (f2,58.2) (f3,70.6) (f4,84.8)};
      \addplot[fill=jugC!30,  draw=jugC]
        coordinates {(f1,93.2) (f2,87.3) (f3,88.2) (f4,93.9)};
      \addplot[fill=propia!35, draw=propia]
        coordinates {(f1,90.9) (f2,83.6) (f3,82.4) (f4,92.4)};
      \legend{Referencia 2026, Terraform, Bicep, Pulumi, Ruta propia}
    \end{axis}
  \end{tikzpicture}
  \caption{Perfil de utilidad de las cuatro figuras de mérito orientadas al
  cliente, antes de la agregación ponderada. La única diferencia relevante
  entre las cuatro posiciones avanzadas está en la cobertura de codificación.
  Elaboración propia.}
  \label{fig:perfilfom}
\end{figure}

\subsec{05. Descomposición en vectores tecnológicos}

El gráfico vectorial no dibuja las posiciones finales como puntos aislados
sino como la suma de los saltos tecnológicos que conducen a ellas. El
Cuadro~\ref{tab:vectores} identifica cada vector, la inversión que representa
y su contribución a las dos coordenadas.

\begin{table}[htbp]
  \centering
  \footnotesize
  \setlength{\tabcolsep}{5pt}
  \caption{Descomposición de cada trayectoria en vectores tecnológicos. Las
  contribuciones son aditivas y suman la posición final de cada jugador.}
  \label{tab:vectores}
  \begin{tabularx}{\textwidth}{@{}l>{\raggedright\arraybackslash}X rr@{}}
    \toprule
    \textbf{Vector} & \textbf{Inversión tecnológica que representa} &
      \textbf{$\Delta C$} & \textbf{$\Delta V$} \\
    \midrule
    \multicolumn{4}{@{}l}{\textit{Jugador A: Terraform}} \\
    A1 & Ejecución y estado gestionados en el servicio del proveedor, con
         control de acceso por espacio de trabajo & $+24$ & $+30{,}0$ \\
    A2 & Despliegue multientorno por pilas y política como código con Sentinel
       & $+8$ & $+15{,}5$ \\
    \midrule
    \multicolumn{4}{@{}l}{\textit{Jugador B: Bicep}} \\
    B1 & Sustitución de las plantillas ARM por el DSL nativo, que suprime el
         fichero de estado y su gobierno & $-20$ & $+14{,}0$ \\
    B2 & Pilas de despliegue, módulos verificados del proveedor y Azure Policy
         en el plano de control & $-4$ & $+19{,}4$ \\
    \midrule
    \multicolumn{4}{@{}l}{\textit{Jugador C: Pulumi}} \\
    C1 & Descripción de la infraestructura en un lenguaje de propósito general
         con pruebas unitarias reales & $+32$ & $+26{,}0$ \\
    C2 & Gestión de entornos y secretos y interfaz programable de automatización
         \parencite{pulumiesc} & $+12$ & $+24{,}3$ \\
    \midrule
    \multicolumn{4}{@{}l}{\textit{Ruta propia}} \\
    P1 & Catálogo de módulos certificados y política como código con OPA
         \parencite{opa} & $+6$ & $+21{,}0$ \\
    P2 & Sustitución del motor por la bifurcación abierta y estado gestionado
         propio \parencite{opentofu} & $-18$ & $+1{,}9$ \\
    P3 & Interfaz de autoservicio construida sobre el catálogo
         \parencite{backstage} & $+8$ & $+24{,}0$ \\
    \midrule
    \multicolumn{4}{@{}l}{\textit{Opción descartada}} \\
    D1 & Contención de coste: reducir el alcance codificado al núcleo crítico
         y devolver el resto a operación manual & $-26$ & $-16{,}6$ \\
    \bottomrule
  \end{tabularx}
\end{table}

El orden de los vectores de la ruta propia no es arbitrario, y es la decisión
más discutible de toda la sección. El catálogo de módulos se construye
\emph{antes} de sustituir el motor, y no después, porque el catálogo es
precisamente lo que abarata la sustitución: cuando los equipos de producto
consumen módulos y no recursos, el motor queda detrás de una interfaz estable
y cambiarlo deja de ser un proyecto de migración para convertirse en un
cambio interno del catálogo. Invertir el orden convertiría P2 en el vector más
caro y arriesgado de los tres.

\subsec{06. El gráfico vectorial}

La Figura~\ref{fig:vectorial} compone los vectores del
Cuadro~\ref{tab:vectores} sobre el plano de las dos coordenadas. Se lee así:
todas las trayectorias arrancan del origen, que es la plataforma actual; cada
flecha es una inversión tecnológica concreta y se encadena a la anterior, de
modo que la posición de un jugador es la suma vectorial de sus inversiones y
no un punto postulado; el punto grueso del final de cada cadena es la posición
en 2029; y la flecha discontinua es la opción que se evaluó y se descartó.
Cuanto más arriba, más valor recibe el equipo de producto; cuanto más a la
izquierda, menos cuesta sostener la plataforma. El cuadrante superior
izquierdo es, por tanto, el favorable.

\begin{figure}[ht]
  \centering
  \begin{tikzpicture}[x=0.125cm, y=0.113cm,
      vec/.style  = {line width=1.05pt, -{Stealth[length=5.5pt]}},
      pto/.style  = {circle, fill, inner sep=2.1pt},
      cod/.style  = {font=\scriptsize\bfseries, inner sep=1.6pt},
      jugl/.style = {font=\footnotesize\bfseries, align=center, inner sep=2pt}]

    % rejilla de diez puntos de índice
    \foreach \x in {-40,-30,-20,-10,10,20,30,40,50}
      \draw[gray!20, thin] (\x,-24) -- (\x,56);
    \foreach \y in {-20,-10,10,20,30,40,50}
      \draw[gray!20, thin] (-42,\y) -- (58,\y);

    % ejes
    \draw[thick,-{Stealth[length=6pt]}] (-44,0) -- (62,0);
    \draw[thick,-{Stealth[length=6pt]}] (0,-26) -- (0,60);
    \node[font=\footnotesize\bfseries, align=center, anchor=south]
      at (2,61) {$\Delta$ Valor\\para el cliente};
    \node[font=\footnotesize\bfseries, align=center, anchor=north]
      at (52,-5) {$\Delta$ Coste para\\el productor};

    % marcas de escala
    \foreach \x in {-40,-30,-20,20,30,40,50}
      \node[font=\tiny, text=gray!55!black, below=1pt] at (\x,0) {\x};
    \foreach \y in {-10,10,20,30,40}
      \node[font=\tiny, text=gray!55!black, left=1pt] at (0,\y) {\y};

    \node[font=\scriptsize\itshape, text=gray!60!black, anchor=north west]
      at (-43,60) {Referencia: 2026 + 3 años};

    % ------------------------------------------------ jugador A: Terraform
    \draw[vec, jugA] (0,0) -- (24,30)
      node[cod, text=jugA, pos=0.55, below right=-1pt] {A1};
    \draw[vec, jugA] (24,30) -- (32,45.5)
      node[cod, text=jugA, midway, right=0pt] {A2};
    \node[pto, jugA] at (32,45.5) {};
    \node[jugl, text=jugA, anchor=east] at (29,45.5) {Jugador A\\Terraform};

    % ------------------------------------------------ jugador C: Pulumi
    \draw[vec, jugC] (0,0) -- (32,26)
      node[cod, text=jugC, pos=0.6, below right=-1pt] {C1};
    \draw[vec, jugC] (32,26) -- (44,50.3)
      node[cod, text=jugC, pos=0.32, right=0pt] {C2};
    \node[pto, jugC] at (44,50.3) {};
    \node[jugl, text=jugC, anchor=south] at (44,53) {Jugador C\\Pulumi};

    % ------------------------------------------------ jugador B: Bicep
    \draw[vec, jugB] (0,0) -- (-20,14)
      node[cod, text=jugB, pos=0.55, above right=-2pt] {B1};
    \draw[vec, jugB] (-20,14) -- (-24,33.4)
      node[cod, text=jugB, midway, right=0pt] {B2};
    \node[pto, jugB] at (-24,33.4) {};
    \node[jugl, text=jugB, anchor=east] at (-27,33.4) {Jugador B\\Bicep};

    % ------------------------------------------------ ruta propia
    \draw[vec, propia] (0,0) -- (6,21)
      node[cod, text=propia, pos=0.6, right=0pt] {P1};
    \draw[vec, propia] (6,21) -- (-12,22.9)
      node[cod, text=propia, pos=0.78, below=1pt] {P2};
    \draw[vec, propia] (-12,22.9) -- (-4,46.9)
      node[cod, text=propia, midway, left=0pt] {P3};
    \node[pto, propia] at (-4,46.9) {};
    \node[jugl, text=propia, anchor=east] at (-7,50)
      {Nuestra empresa\\ruta propuesta};

    % ------------------------------------------------ opción descartada
    \draw[vec, descart, dashed] (0,0) -- (-26,-16.6)
      node[cod, text=descart, pos=0.5, below=1pt] {D1};
    \node[pto, descart] at (-26,-16.6) {};
    \node[font=\scriptsize\itshape, text=descart, align=center, anchor=north]
      at (-26,-19.5) {Contención de coste\\(descartada)};

    % ------------------------------------------------ origen
    \node[circle, draw, very thick, fill=white, inner sep=1.7pt] at (0,0) {};
    \node[font=\scriptsize, align=left, anchor=north west] at (7,-7)
      {Producto de referencia:\\plataforma actual, 2026};
  \end{tikzpicture}
  \caption{Gráfico vectorial de posicionamiento competitivo al horizonte 2029.
  El origen es la plataforma actual. Cada flecha es un vector tecnológico del
  Cuadro~\ref{tab:vectores} y cada punto, la posición que alcanza una
  organización que adopta esa trayectoria. La rejilla es de diez puntos de
  índice. Elaboración propia.}
  \label{fig:vectorial}
\end{figure}

\subsec{07. Lectura del gráfico}

El gráfico dice cuatro cosas, y ninguna de ellas es que exista una herramienta
mejor que las demás.

\textbf{Primera: la frontera de eficiencia la forman tres posiciones, no una.}
Entre las cuatro trayectorias de inversión, Pulumi, la ruta propia y Bicep no
se dominan entre sí. Pulumi ofrece el mayor
valor, $+50{,}3$ puntos, y cuesta $+44$; Bicep es el más barato, $-24$, y
entrega $+33{,}4$; la ruta propia se sitúa en $-4$ y $+46{,}9$. Terraform
gestionado, en cambio, sí está dominado: la ruta propia entrega más valor
($+46{,}9$ frente a $+45{,}5$) por 36 puntos menos de coste. Ésa es la
conclusión operativa más limpia de todo el análisis, y explica por qué el
vector P2 existe.

\textbf{Segunda: el valor no lo aporta el motor, lo aportan las capas de
encima.} En la ruta propia, el vector que sustituye el motor, P2, aporta
$+1{,}9$ puntos de valor sobre 46,9, es decir, un 4 \%. Los otros 45 puntos
vienen del catálogo y del autoservicio, que son inversiones propias y no
dependen del proveedor elegido. La consecuencia práctica es que la decisión de
herramienta, que suele consumir la mayor parte del debate técnico, es la menos
determinante de las tres.

\textbf{Tercera: la Figura~\ref{fig:perfilfom} explica de dónde sale la
distancia.} Las cuatro posiciones avanzadas están entre 81 y 93 puntos de
utilidad en velocidad y entre 70 y 88 en fiabilidad, márgenes estrechos. La
cobertura de codificación, en cambio, va de 58 a 87. Lo que separa a los
jugadores no es lo bien que hacen aquello que hacen, sino cuánto de la
infraestructura queda dentro de su alcance.

\textbf{Cuarta: el cuadrante inferior izquierdo es una trampa reconocible.} La
opción descartada D1 ahorra 26 puntos de coste, dos puntos más que Bicep, y
destruye 16,6 puntos de valor. Conviene ser exacto aquí: al ser la más barata
de todas, D1 tampoco está dominada, y una regla de eficiencia de Pareto
aplicada sin más la conservaría en la frontera. Se descarta por un criterio
que el plano no contiene, un suelo de valor aceptable, y esto ilustra el
límite del instrumento: el gráfico ordena posiciones, pero no decide cuál es
admisible. Reducir el alcance codificado no elimina el
trabajo de infraestructura, lo devuelve a la forma manual, donde no se mide y
no aparece en el índice de coste de la plataforma porque se ha desplazado a
los equipos de producto. El vector es barato en el eje que el gráfico mide y caro en uno
que no mide.

\subsec{08. Sensibilidad al supuesto de cartera}

Toda la comparación descansa en un supuesto que conviene poner a prueba: que
la cartera de infraestructura seguirá concentrada en un solo proveedor de
nube. El Cuadro~\ref{tab:sensib} recalcula las cuatro posiciones bajo el
escenario contrario, con un 40 \% de la carga fuera de Azure, que obliga al
jugador B a sostener una segunda cadena de herramientas.

\begin{table}[htbp]
  \centering
  \footnotesize
  \setlength{\tabcolsep}{6pt}
  \caption{Efecto del escenario multinube sobre las posiciones. Sólo el
  jugador B se desplaza de forma apreciable.}
  \label{tab:sensib}
  \begin{tabular}{@{}lrrrrrr@{}}
    \toprule
    & \multicolumn{2}{c}{\textbf{Proveedor único}} &
      \multicolumn{2}{c}{\textbf{Cartera multinube}} &
      \multicolumn{2}{c}{\textbf{Desplazamiento}} \\
    \cmidrule(lr){2-3}\cmidrule(lr){4-5}\cmidrule(lr){6-7}
    \textbf{Posición} & $\Delta C$ & $\Delta V$ & $\Delta C$ & $\Delta V$ &
      $\Delta C$ & $\Delta V$ \\
    \midrule
    A. Terraform gestionado & $+32$ & $+45{,}5$ & $+35$ & $+43{,}6$ &
      $+3$ & $-1{,}9$ \\
    B. Bicep sobre Azure    & $-24$ & $+33{,}4$ & $-4$  & $+10{,}9$ &
      $+20$ & $-22{,}5$ \\
    C. Pulumi               & $+44$ & $+50{,}3$ & $+47$ & $+48{,}4$ &
      $+3$ & $-1{,}9$ \\
    \textbf{Ruta propia}    & $-4$  & $+46{,}9$ & $-1$  & $+45{,}0$ &
      $+3$ & $-1{,}9$ \\
    \bottomrule
  \end{tabular}
\end{table}

El resultado es asimétrico y por eso es útil. Las tres posiciones multinube se
mueven tres puntos en coste y dos en valor, un desplazamiento despreciable. El
jugador B pierde 22,5 puntos de valor y sube 20 de coste, hasta dejar de ser
la opción barata. Conviene precisar en qué sentido, porque el frente de
Pareto del apartado 07 de la sección~\ref{sec:evol} muestra que B no queda
\emph{dominado}: con $\Delta C = -4$ conserva tres puntos de ventaja en coste
sobre la ruta propia, que queda en $-1$, y por tanto sigue formalmente en la
frontera de eficiencia. Lo que ocurre es más elocuente que la dominancia. La
tasa de canje entre B y la ruta propia pasa de 0,68 a 11,37 puntos de valor
por punto de coste, es decir, esos tres puntos de coste ahorrado cuestan 34,1
puntos de valor. B deja de ser el punto barato y razonable de la frontera para
convertirse en un extremo del que ninguna función de preferencia plausible
escogería. Su ventaja competitiva no es una propiedad de la tecnología sino
una propiedad del supuesto de cartera, y desaparece con él.

Esto convierte la elección en una apuesta explícita y no en una preferencia
técnica: adoptar Bicep es afirmar que la empresa no abrirá una segunda nube en
el horizonte del plan. Si esa afirmación no puede sostenerse por escrito ante
la dirección, el jugador B queda fuera con independencia de lo que diga el
gráfico principal.

\subsec{09. Límites del análisis}

Tres reservas que conviene dejar escritas junto al gráfico, porque un plano
con flechas transmite más certidumbre de la que el método admite.

\begin{enumerate}
  \item \textbf{Las coordenadas son estimaciones, no medidas.} Las figuras de
    mérito del Cuadro~\ref{tab:fomcomp} proceden del comportamiento observado
    de la plataforma actual y de las propiedades declaradas de cada
    tecnología. El propio Cuadro~\ref{tab:fom} define cómo medirlas: la
    manera correcta de validar este análisis es instrumentar la referencia y
    volver a dibujar el gráfico con datos de un trimestre real.
  \item \textbf{Los pesos del índice de valor son una decisión de gobierno.}
    Al ponderar el tiempo de ciclo con 0,30 y el de reconstrucción con 0,20 se
    está escogiendo una postura ante el riesgo. Una organización sujeta a un
    requisito de recuperación ante desastres invertiría esos pesos, y con
    ellos el orden de los dos primeros clasificados.
  \item \textbf{El eje de coste ignora la transición.} Como se advirtió en el
    apartado 03, mide el régimen permanente de 2029. La ruta propia exige tres
    inversiones sucesivas cuyo desembolso inicial no aparece en ninguna
    coordenada, y esa cifra es la que decide si el plan es financiable. Su
    lugar es el elemento 8 del índice, y la sección~\ref{sec:financiero} lo
    resuelve: 820,2 miles de euros de inversión, 335,6 de exposición máxima y
    un valor actual neto de 743,3 al 9 \pct.
\end{enumerate}

% ==========================================================================
\section{Modelo técnico}
\label{sec:tecnico}
% ==========================================================================

El elemento 7 del índice pide el modelo técnico: el conjunto de relaciones
cuantitativas que explican \emph{por qué} las figuras de mérito toman los
valores que toman. Es el elemento que distingue una hoja de ruta de una lista
de objetivos. Sin él, las columnas de 2029 del Cuadro~\ref{tab:anclas} son
aspiraciones; con él, son consecuencias de parámetros que la cartera de
proyectos modifica y que pueden medirse por separado.

El modelo se construye sobre una hipótesis de estructura: las cuatro figuras
de entrega no son independientes, sino funciones de las cuatro figuras de
infraestructura y de tres parámetros de adopción que los proyectos de la
cartera desplazan. Esos tres parámetros son:

\begin{description}[leftmargin=1.6em, style=nextline]
  \item[$\sigma$, adopción del autoservicio] Fracción de las peticiones de
    infraestructura que el equipo de producto resuelve por sí mismo, sin
    entrar en la cola del equipo de plataforma. Lo desplaza R7.
  \item[$\gamma$, cumplimiento automatizado] Fracción de los planes que la
    evaluación de políticas resuelve sin excepción manual. Coincide con la
    figura de mérito de cumplimiento del Cuadro~\ref{tab:fom}. Lo desplaza R4.
  \item[$\kappa$, penetración del catálogo] Fracción del contenido de un
    cambio típico construida con módulos certificados en lugar de recursos
    declarados a mano. Lo desplaza R3.
\end{description}

\subsec{01. Ecuación del tiempo de ciclo}

El tiempo de ciclo se descompone en los cuatro tramos del ciclo de operación
de la sección~\ref{sec:ciclo}, y cada tramo depende de un parámetro distinto:

\[
T_c = \underbrace{(1-\sigma)\,q}_{\text{espera}}
    + \underbrace{(1-\gamma)\,h_h + \gamma\,h_a}_{\text{revisión}}
    + \underbrace{t_v}_{\text{validación}}
    + \underbrace{(1-\kappa)\,a_m + \kappa\,a_c}_{\text{aplicación}} .
\]

Los coeficientes son propiedades medibles del proceso y no del plan:
$q = 16$ h es la espera media en la cola de plataforma, $h_h = 14$ h la
demora de una revisión humana y $h_a = 0{,}2$ h la de una evaluación
automática, $t_v = 1{,}8$ h la validación, y $a_m = 3{,}2$ h frente a
$a_c = 0{,}7$ h el trabajo de aplicación según se declare a mano o por
módulo.

Sustituidos los valores de 2026 ($\sigma = 0$, $\gamma = 0{,}62$,
$\kappa = 0{,}15$) la ecuación devuelve 26,1 horas, y sustituidos los de 2029
($\sigma = 0{,}75$, $\gamma = 0{,}94$, $\kappa = 0{,}80$) devuelve 8,0 horas.
Coinciden con las dos columnas del Cuadro~\ref{tab:anclas}, que no se habían
obtenido de este modelo sino estimadas por separado en la
sección~\ref{sec:comp}: la coincidencia es la primera comprobación de
consistencia del documento.

Más útil que el valor es la atribución. Activando un parámetro a la vez, el
salto de 26 a 8 horas se reparte así: la interfaz de autoservicio aporta 12,0
de las 18 horas ganadas (66,5 \%), la política como código 4,4 horas (24,5 \%)
y el catálogo de módulos 1,6 horas (9,0 \%). El resultado desmiente la
intuición de que el catálogo es lo que acelera el ciclo: el catálogo acelera
la \emph{aplicación}, que es el tramo más corto de los cuatro. Lo que domina
el tiempo de ciclo es la espera en la cola, y eso sólo lo elimina el
autoservicio.

Esto no degrada a R3 sino que precisa su papel: el catálogo importa por la
tasa de fallo (apartado 03) y como requisito de R7, no por el tiempo de ciclo.
La DSM de la sección~\ref{sec:dsm} ya lo había establecido por otra vía, al
situar R3 como dependencia necesaria de R7.

\subsec{02. Ecuación de la deriva}

La deriva es la única figura de mérito que se comporta como un estado y no
como un suceso, y por eso admite un modelo dinámico. Si $D$ es el número de
recursos desviados, $N$ el parque, $C$ la cobertura, $\lambda$ la tasa de
intervención manual por recurso no codificado y $\mu$ la tasa de conciliación,

\[
\frac{dD}{dt} = \lambda\,(1-C)\,N - \mu\,D
\qquad\Longrightarrow\qquad
D^{*} = \frac{\lambda\,(1-C)\,N}{\mu},
\]

de modo que la deriva por cada mil recursos en régimen permanente es
$d^{*} = 1000\,\lambda\,(1-C)/\mu$. La ecuación dice algo que la figura de
mérito aislada no dice: la deriva no depende del tamaño del parque, sólo de la
cobertura y de la frecuencia con que se concilia.

Calibrado sobre 2026 ($C = 0{,}55$, $\mu = 0{,}5$ semanas$^{-1}$, esto es una
conciliación efectiva cada catorce días, y $d^{*} = 38$) resulta
$\lambda = 0{,}0422$ intervenciones por recurso y semana. Con ese valor fijo,
el objetivo de $d^{*} = 8$ para 2029 exige $\mu = 0{,}739$, es decir, acortar
el intervalo de detección de catorce a nueve días y medio.

La conclusión es la que justifica el par de proyectos, y ninguno de los dos
por separado:

\begin{itemize}
  \item Con la cobertura al 86 \% pero el intervalo de detección sin cambiar,
    $d^{*} = 11{,}8$: se incumple el objetivo.
  \item Con el intervalo acortado a nueve días y medio pero la cobertura al
    55 \%, $d^{*} = 25{,}7$: se incumple con holgura.
  \item Con ambos, $d^{*} = 8{,}0$.
\end{itemize}

R1 y R8 son por tanto complementarios y no alternativos, lo que explica por
qué el Cuadro~\ref{tab:trazafom} los asigna conjuntamente a la deriva y por
qué la sección~\ref{sec:dsm} sitúa R8 en la tercera ola y no al final.

\subsec{03. Ecuaciones de la reconstrucción y del fallo}

\textbf{Tiempo de reconstrucción.} Se descompone en la parte que el motor
reconstruye por sí mismo y el resto que hay que rehacer a mano:

\[
T_r = T_{\text{auto}} + (1-C)\,N\,\tau_m .
\]

Con $T_{\text{auto}} = 4{,}0$ h, $N = 1200$, $C = 0{,}55$ y
$\tau_m = 4$ minutos por recurso, resulta $T_r = 40{,}0$ h para 2026; con
$T_{\text{auto}} = 3{,}2$ h, $N = 1825$, $C = 0{,}86$ y
$\tau_m = 1{,}8$ minutos, resulta 11,0 h para 2029. Ambos coinciden de nuevo
con el Cuadro~\ref{tab:anclas}.

El reparto es el dato relevante: en 2026 el resto manual son 36 de las 40
horas, el 90 \% del total, y en 2029 son 7,8 de 11, el 71 \%. El tiempo de
reconstrucción es, en la práctica, una medida indirecta de la fracción no
codificada del parque. Esto respalda con una ecuación lo que la
sección~\ref{sec:fom} afirmaba de manera cualitativa, que es la figura que
mejor resume el estado de una implantación.

\textbf{Tasa de fallo del cambio.} Se modela como una tasa base atenuada
multiplicativamente por los dos controles, cada uno con su cobertura de modos
de fallo:

\[
f = f_0\,(1 - \pi\gamma)\,(1 - \theta\kappa),
\]

con $f_0 = 0{,}20$ la tasa sin ningún control. Ajustando $\pi$ y $\theta$
para que la ecuación reproduzca el 14 \% de 2026 y el 6 \% de 2029 se obtiene
$\pi = 0{,}354$ y $\theta = 0{,}688$.

Los dos coeficientes son el resultado más contraintuitivo del modelo
técnico. El catálogo certificado atrapa el 68,8 \% de los modos de fallo y la
política como código sólo el 35,4 \%, casi la mitad. Actuando por separado
desde la situación de 2026, el catálogo llevaría la tasa de fallo al 7,0 \% y
la política sólo al 12,0 \%. La razón es de naturaleza y no de calidad: la
política evalúa reglas declaradas de antemano y por tanto sólo detecta los
fallos que alguien anticipó, mientras que un módulo certificado impide de raíz
las combinaciones de parámetros que no ha validado nadie. La política atrapa
lo previsto; el catálogo elimina lo imprevisto.

\subsec{04. Sensibilidad del modelo}

Un modelo que reproduce sus dos puntos de calibración no está validado: hay
que saber cuál de sus parámetros lo domina. La Figura~\ref{fig:tornado}
recoge el análisis sobre el tiempo de ciclo en el punto de 2029, variando
cada parámetro de adopción en $\pm 0{,}15$ absolutos y cada coeficiente de
proceso en $\pm 25$ \%.

\begin{figure}[ht]
  \centering
  \begin{tikzpicture}
    \begin{axis}[
      width=13.6cm, height=6.8cm,
      xbar stacked, bar width=9pt,
      xmin=5.0, xmax=11.0,
      y=0.74cm,
      ytick={1,2,3,4,5,6,7,8},
      yticklabels={$a_c$ ($\pm 25$ \pct), $a_m$ ($\pm 25$ \pct),
                   $h_h$ ($\pm 25$ \pct), $\kappa$ ($\pm 0{,}15$),
                   $t_v$ ($\pm 25$ \pct), $q$ ($\pm 25$ \pct),
                   $\gamma$ ($\pm 0{,}15$), $\sigma$ ($\pm 0{,}15$)},
      yticklabel style={font=\scriptsize},
      xticklabel style={font=\scriptsize},
      xtick={5,5.5,6,6.5,7,7.5,8,8.5,9,9.5,10,10.5,11},
      xlabel={\footnotesize Tiempo de ciclo resultante (horas)},
      xlabel style={font=\footnotesize},
      xmajorgrids, grid style={gray!20},
      axis lines=left, axis on top,
      legend style={font=\scriptsize, at={(0.5,1.04)}, anchor=south,
                    legend columns=2, draw=none, column sep=1.6ex},
      area legend,
    ]
      % Tramo invisible: del origen al extremo bajo de cada parámetro.
      \addplot[draw=none, fill=none, forget plot] coordinates {
        (7.89,1) (7.87,2) (7.82,3) (7.65,4) (7.58,5) (7.03,6) (7.20,7)
        (5.63,8)};
      % Del extremo bajo al valor base (8,03 h).
      \addplot[fill=capa!30, draw=capa] coordinates {
        (0.14,1) (0.16,2) (0.21,3) (0.38,4) (0.45,5) (1.00,6) (0.83,7)
        (2.40,8)};
      \addlegendentry{Valor del parámetro en su extremo favorable}
      % Del valor base al extremo alto.
      \addplot[fill=acento!30, draw=acento] coordinates {
        (0.14,1) (0.16,2) (0.21,3) (0.37,4) (0.45,5) (1.00,6) (2.07,7)
        (2.40,8)};
      \addlegendentry{Valor del parámetro en su extremo adverso}
      \draw[gray!30!black, very thick]
        ({axis cs:8.03,0}|-{rel axis cs:0,0}) --
        ({axis cs:8.03,0}|-{rel axis cs:0,1});
      \node[font=\scriptsize, text=gray!30!black, anchor=south]
        at (axis cs:8.03,8.55) {Base: 8,0 h};
    \end{axis}
  \end{tikzpicture}
  \caption{Sensibilidad del tiempo de ciclo de 2029 a cada parámetro del
  modelo, ordenada por amplitud. Cada barra parte del valor base de 8,0 horas
  y llega al resultado que produce el parámetro en cada extremo de su
  intervalo. La adopción del autoservicio domina, con una amplitud de 4,8
  horas. Elaboración propia.}
  \label{fig:tornado}
\end{figure}

El resultado concentra el riesgo técnico de la hoja de ruta en un solo
parámetro. La adopción del autoservicio, $\sigma$, produce una amplitud de
4,8 horas sobre un valor base de 8,0: si la adopción se queda en el 60 \% en
lugar del 75 \%, el tiempo de ciclo es de 10,4 horas y el objetivo se
incumple, con independencia de que todo lo demás salga bien. Los dos
parámetros siguientes, $\gamma$ y $q$, producen amplitudes de 2,9 y 2,0
horas, y los cinco restantes juntos no llegan a 2,5.

Y $\sigma$ no es un parámetro técnico. Es una medida de cuántos equipos de
producto prefieren usar el portal antes que abrir una petición, esto es, un
problema de diseño de la interfaz y de confianza en el catálogo, no de
infraestructura. El modelo técnico termina por tanto señalando que el mayor
riesgo de una hoja de ruta de infraestructura no está en la infraestructura.
La consecuencia se recoge en el prototipo P4 de la
sección~\ref{sec:portfolio}, cuyo único objetivo es medir $\sigma$ con
usuarios reales antes de comprometer el presupuesto de R7.

\subsec{05. Qué no modela}

El modelo tiene tres huecos declarados. No modela la frecuencia de aplicación,
la cuarta figura de entrega, porque depende de la demanda de cambio del
negocio y no de la plataforma: la plataforma retira el límite, pero no genera
la demanda. No modela el tiempo de restablecimiento, que depende del
procedimiento de guardia, expresamente fuera de la frontera fijada en el
Cuadro~\ref{tab:frontera}. Y trata $\lambda$, $\pi$ y $\theta$ como
constantes, cuando es razonable esperar que las intervenciones manuales que
quedan tras codificar el 86 \% del parque sean precisamente las más difíciles,
y por tanto que $\lambda$ crezca al subir la cobertura. Ese efecto empujaría
la deriva por encima del objetivo, y no se ha incorporado por falta de base
para estimarlo.

% ==========================================================================
\section{Modelo financiero}
\label{sec:financiero}
% ==========================================================================

Las versiones anteriores de este documento declaraban el elemento 8 fuera de
alcance, y el apartado 09 de la sección~\ref{sec:comp} señalaba la carencia
como la más grave de las tres que enumeraba: el eje de coste del gráfico
vectorial mide el régimen permanente de 2029 y no la inversión de transición,
de modo que ninguna coordenada respondía a la pregunta de si el plan es
financiable. Esta sección la responde.

El modelo se construye con la misma disciplina que el resto del documento: los
beneficios no se postulan, se derivan de las figuras de mérito del
Cuadro~\ref{tab:anclas} mediante las ecuaciones de la
sección~\ref{sec:tecnico}. Un beneficio que no pueda escribirse como función
de una figura de mérito no entra en el modelo, y por eso queda fuera el
concepto más citado en las justificaciones de este tipo de inversión, la
productividad del desarrollador, que ninguna de las ocho figuras mide.

\subsec{01. La inversión: los ocho proyectos}

El Cuadro~\ref{tab:inversion} cuantifica la cartera. La unidad de esfuerzo es
el mes-persona, valorado a 7,92 miles de euros según un coste anual completo
de 95 mil euros por ingeniero. El calendario es el de la secuencia en seis
olas que la DSM produjo en la sección~\ref{sec:dsm}, y no una distribución
uniforme.

\begin{table}[htbp]
  \centering
  \footnotesize
  \setlength{\tabcolsep}{5pt}
  \caption{Inversión por proyecto, en meses-persona y en miles de euros. El
  año de ejecución procede del particionado de la DSM de la cartera
  (Figura~\ref{fig:dsmproyord}). La partida de otros conceptos recoge
  herramienta, servicio y formación.}
  \label{tab:inversion}
  \begin{tabular}{@{}llrrr@{}}
    \toprule
    \textbf{Cód.} & \textbf{Proyecto} & \textbf{Meses-p.} &
      \textbf{Miles €} & \textbf{Año} \\
    \midrule
    R2 & Custodia del registro de estado      &  6 &  47{,}5 & 2026 \\
    R1 & Cierre de cobertura de codificación  & 14 & 110{,}8 & 2026 \\
    R3 & Catálogo de módulos certificados     & 18 & 142{,}5 & 2027 \\
    R8 & Observabilidad de la plataforma      &  8 &  63{,}3 & 2027 \\
    R4 & Política como código                 & 10 &  79{,}2 & 2028 \\
    R5 & Malla de pruebas y entornos efímeros & 12 &  95{,}0 & 2028 \\
    R6 & Sustitución del motor                &  9 &  71{,}2 & 2028 \\
    R7 & Interfaz de autoservicio             & 16 & 126{,}7 & 2029 \\
    \midrule
    \multicolumn{2}{@{}l}{\textbf{Total de ingeniería}} & \textbf{93} &
      \textbf{736{,}2} & \\
    \multicolumn{2}{@{}l}{Otros conceptos}   & -- &  84{,}0 & \\
    \multicolumn{2}{@{}l}{\textbf{Inversión total}} & -- &
      \textbf{820{,}2} & \\
    \bottomrule
  \end{tabular}
\end{table}

Los 93 meses-persona equivalen a 2,6 ingenieros a tiempo completo durante los
tres años y medio del plan. Es una cifra que conviene retener, porque sitúa la
hoja de ruta en su escala: no es un programa de transformación sino el trabajo
sostenido de un equipo pequeño.

\subsec{02. Los beneficios, derivados de las figuras de mérito}

Se contabilizan cuatro corrientes de beneficio, cada una anclada a una figura
de mérito y a la ecuación que la gobierna. El parque de recursos parte de
1.200 en 2026 y se supone que crece un 15 \% anual; la cobertura sigue la
trayectoria de la cartera, y el caso de referencia contra el que se compara es
el de no hacer nada, en el que la cobertura permanece en el 55 \% y la tasa de
fallo en el 14 \%.

\begin{description}[leftmargin=1.6em, style=nextline]
  \item[B1. Operación manual evitada] Los recursos no codificados consumen
    trabajo recurrente, estimado en 400 euros por recurso y año. El beneficio
    es $N(t)\,[C(t) - 0{,}55]\times 0{,}4$ miles de euros, y procede
    directamente de la figura de cobertura.
  \item[B2. Retrabajo por fallo evitado] Cada cambio fallido cuesta 1.600
    euros entre corrección e impacto en el servicio. El beneficio es
    $n_c(t)\,[0{,}14 - f(t)]\times 1{,}6$, con $f(t)$ la tasa de fallo que
    predice la ecuación del apartado 03 de la sección~\ref{sec:tecnico}.
  \item[B3. Capacidad de plataforma liberada] Cada cambio entregado consume
    hoy 2,1 horas de acompañamiento del equipo de plataforma que el catálogo
    y el autoservicio suprimen, a 60 euros por hora.
  \item[B4. Menor coste unitario de propiedad] Los cuatro puntos de índice que
    la ruta propia ahorra en el Cuadro~\ref{tab:coste}, aplicados al parque
    creciente. Es la corriente más pequeña de las cuatro, y esa
    desproporción es en sí misma un resultado.
\end{description}

El Cuadro~\ref{tab:flujos} recoge el flujo de caja resultante.

\begin{table}[htbp]
  \centering
  \footnotesize
  \setlength{\tabcolsep}{4.5pt}
  \caption{Flujo de caja del plan, en miles de euros. El beneficio de 2026 es
  nulo porque las dos primeras olas no mueven ninguna figura de mérito hasta
  su cierre.}
  \label{tab:flujos}
  \begin{tabular}{@{}lrrrrrrrrr@{}}
    \toprule
    \textbf{Año} & \textbf{Parque} & \textbf{Cobert.} & \textbf{B1} &
      \textbf{B2} & \textbf{B3} & \textbf{B4} & \textbf{Benef.} &
      \textbf{Invers.} & \textbf{Flujo} \\
    \midrule
    2026 & 1.200 & 55 \pct &   0{,}0 &   0{,}0 &   0{,}0 &  0{,}0 &   0{,}0 & 178{,}3 & $-178{,}3$ \\
    2027 & 1.380 & 62 \pct &  38{,}6 &  23{,}0 &  20{,}5 &  3{,}1 &  85{,}3 & 229{,}8 & $-144{,}6$ \\
    2028 & 1.587 & 74 \pct & 120{,}6 &  64{,}8 &  62{,}6 &  9{,}7 & 257{,}7 & 270{,}4 & $-12{,}7$ \\
    2029 & 1.825 & 86 \pct & 226{,}3 & 115{,}2 & 113{,}4 & 18{,}2 & 473{,}1 & 141{,}7 & $+331{,}5$ \\
    2030 & 2.099 & 88 \pct & 277{,}1 & 128{,}0 & 134{,}1 & 22{,}3 & 561{,}5 &   0{,}0 & $+561{,}5$ \\
    2031 & 2.414 & 88 \pct & 318{,}6 & 140{,}8 & 147{,}5 & 25{,}7 & 632{,}7 &   0{,}0 & $+632{,}7$ \\
    \bottomrule
  \end{tabular}
\end{table}

Con un coste de capital del 9 \%, el valor actual neto del plan es de
\textbf{743,3 miles de euros}, la tasa interna de retorno del
\textbf{51,8 \%} y el retorno de la inversión se produce en el primer mes de
\textbf{2030}, cuatro años después del arranque. La
Figura~\ref{fig:flujoacum} representa el flujo acumulado.

\begin{figure}[ht]
  \centering
  \begin{tikzpicture}
    \begin{axis}[
      width=15.2cm, height=6.2cm,
      xmin=2025.4, xmax=2031.7, ymin=-600, ymax=1340,
      xlabel={\footnotesize Año}, xlabel style={font=\footnotesize},
      ylabel={\footnotesize Miles de euros},
      ylabel style={font=\footnotesize},
      xtick={2026,2027,2028,2029,2030,2031},
      ytick={-400,-200,0,200,400,600,800,1000,1200},
      xticklabel style={font=\scriptsize}, yticklabel style={font=\scriptsize},
      ymajorgrids, grid style={gray!18}, axis lines=left,
      legend style={font=\scriptsize, at={(0.02,0.97)}, anchor=north west,
                    draw=gray!45, fill=white},
      area legend,
    ]
      \addplot[ybar, bar width=13pt, fill=acento!22, draw=acento]
        coordinates {(2026,-178.3) (2027,-229.8) (2028,-270.4)
                     (2029,-141.7) (2030,0) (2031,0)};
      \addlegendentry{Inversión del año}
      \addplot[ybar, bar width=13pt, fill=capa!22, draw=capa]
        coordinates {(2026,0) (2027,85.3) (2028,257.7) (2029,473.1)
                     (2030,561.5) (2031,632.7)};
      \addlegendentry{Beneficio del año}
      \addplot[propia, very thick, mark=*, mark size=2.2pt, sharp plot]
        coordinates {(2026,-178.3) (2027,-322.9) (2028,-335.6)
                     (2029,-4.1) (2030,557.4) (2031,1190.1)};
      \addlegendentry{Flujo acumulado a fin de año}
      \draw[gray!45!black, thick] (axis cs:2025.4,0) -- (axis cs:2031.7,0);
      % retorno de la inversión: el acumulado cruza el cero al entrar en 2030
      \draw[propia, thin] (axis cs:2029.1,-40) -- (axis cs:2029.6,-200);
      \node[font=\scriptsize, text=propia, anchor=west, align=left]
        at (axis cs:2029.65,-215) {Retorno de la\\inversión: 2030};
      % exposición máxima de tesorería, al cierre de 2028
      \draw[acento, thin] (axis cs:2028,-352) -- (axis cs:2027.42,-450);
      \node[font=\scriptsize, text=acento, anchor=east, align=right]
        at (axis cs:2027.38,-455) {Máxima exposición\\de tesorería: 335,6};
    \end{axis}
  \end{tikzpicture}
  \caption{Flujo de caja anual y acumulado del plan. La exposición máxima de
  tesorería, 335,6 miles de euros, se alcanza al cierre de 2028, y es la cifra
  que la dirección tiene que autorizar. Elaboración propia.}
  \label{fig:flujoacum}
\end{figure}

La cifra que la dirección necesita no es el valor actual neto sino la
exposición máxima de tesorería: 335,6 miles de euros al cierre de 2028. Es lo
que hay que financiar antes de que el plan se sostenga solo.

\subsec{03. Sensibilidad y umbral de rentabilidad}

El Cuadro~\ref{tab:sensfin} recalcula el valor actual neto bajo siete
variaciones de los supuestos. Se ha escogido una por cada supuesto que un
comité de inversión cuestionaría.

\begin{table}[htbp]
  \centering
  \footnotesize
  \setlength{\tabcolsep}{6pt}
  \caption{Sensibilidad del valor actual neto a los supuestos del modelo, en
  miles de euros.}
  \label{tab:sensfin}
  \begin{tabular}{@{}lrr@{}}
    \toprule
    \textbf{Escenario} & \textbf{VAN} & \textbf{Variación} \\
    \midrule
    Caso base                                     & 743{,}3 & -- \\
    \midrule
    Cobertura final del 74 \pct{} y no del 86 \pct & 308{,}7 & $-434{,}6$ \\
    Sólo B1 y B2 (sin capacidad ni coste unitario) & 393{,}4 & $-349{,}9$ \\
    Parque plano, sin crecimiento                  & 443{,}5 & $-299{,}8$ \\
    Sobrecoste de la inversión del 40 \pct         & 452{,}8 & $-290{,}5$ \\
    Coste de la operación manual un 30 \pct{} menor & 528{,}8 & $-214{,}5$ \\
    Coste de capital del 15 \pct                   & 539{,}9 & $-203{,}4$ \\
    Coste de la operación manual un 30 \pct{} mayor & 957{,}8 & $+214{,}5$ \\
    \bottomrule
  \end{tabular}
\end{table}

Tres conclusiones.

\textbf{El plan es rentable en los siete escenarios.} El peor de ellos deja
un valor actual neto de 308,7 miles de euros, y ninguna combinación razonable
de dos supuestos adversos lo lleva a cero. La decisión de invertir no depende,
por tanto, de acertar en los parámetros.

\textbf{Pero una sola variable domina el resultado.} Quedarse en el 74 \% de
cobertura en lugar del 86 \% destruye 434,6 miles de euros, más del doble que
un sobrecoste del 40 \% en toda la cartera. Es la tercera vez que la cobertura
de codificación aparece como la variable crítica del documento, después del
frente de Pareto de la sección~\ref{sec:evol} y de la matriz de motivaciones
de la sección~\ref{sec:motiv}, y por caminos independientes. La conclusión
operativa es incómoda y clara: es preferible un sobrecoste del 40 \% que
alcance el 86 \% de cobertura a un plan en presupuesto que se detenga en el
74 \%.

\textbf{La corriente de beneficio más blanda no es la que sostiene el
plan.} Eliminando por completo B3 y B4, esto es, quedándose sólo con la
operación manual evitada y el retrabajo, el valor actual neto sigue siendo de
393,4 miles de euros. El caso de inversión no descansa en la partida
discutible.

\subsec{04. Límites del modelo financiero}

\begin{enumerate}
  \item \textbf{Los precios unitarios son supuestos, no cotizaciones.} Los 400
    euros por recurso no codificado y año, los 1.600 por cambio fallido y las
    2,1 horas de acompañamiento por cambio son estimaciones de ingeniería. El
    Cuadro~\ref{tab:sensfin} acota el efecto del primero, que es el mayor de
    los tres, y el proyecto R8 los convertiría en medidas durante 2027.
  \item \textbf{No hay coste de oportunidad del equipo.} Los 93 meses-persona
    se valoran a su coste y no al valor de lo que el equipo dejará de hacer
    para ejecutarlos. Si esa alternativa tuviera un retorno superior al 51,8
    \% de la tasa interna de retorno del plan, la comparación cambiaría.
  \item \textbf{El horizonte de seis años favorece al plan.} Dos de los seis
    años son de beneficio puro sin inversión. Un horizonte de cuatro años,
    que es el que muchos comités aplican, deja el valor actual neto en
    territorio marginal, y el argumento pasa a depender del valor terminal de
    la plataforma, que este modelo no calcula.
\end{enumerate}

% ==========================================================================
\section{Cartera de proyectos y prototipos}
\label{sec:portfolio}
% ==========================================================================

El elemento 9 del índice estaba atendido sólo en parte: el
Cuadro~\ref{tab:dsmproy} enumera los ocho proyectos y la
Figura~\ref{fig:dsmproyord} deriva su secuencia, pero una cartera no queda
descrita por su lista y su orden. Falta lo que convierte una lista en cartera:
la madurez de partida de cada proyecto, los prototipos que reducen su
incertidumbre antes de comprometer el presupuesto, y las puertas de decisión
en las que la cartera puede detenerse.

\subsec{01. Nivel de madurez de partida}

El Cuadro~\ref{tab:trl} asigna a cada proyecto su nivel de madurez
tecnológica de partida, en la escala de nueve niveles de uso corriente
\parencite{mankins1995}, y el nivel al que tiene que llegar para considerarse
cerrado.

\begin{table}[htbp]
  \centering
  \footnotesize
  \setlength{\tabcolsep}{4pt}
  \caption{Nivel de madurez tecnológica de partida y de llegada de cada
  proyecto, y naturaleza de la incertidumbre dominante. Ningún proyecto parte
  por debajo del nivel 6.}
  \label{tab:trl}
  \begin{tabularx}{\textwidth}{@{}llcc>{\raggedright\arraybackslash}X@{}}
    \toprule
    \textbf{Cód.} & \textbf{Proyecto} & \textbf{Part.} & \textbf{Lleg.} &
      \textbf{Incertidumbre dominante} \\
    \midrule
    R2 & Custodia del estado      & 9 & 9 & Ninguna tecnológica: es
      configuración de un servicio disponible. \\
    R1 & Cierre de cobertura      & 8 & 9 & Volumen: cuántos recursos
      resisten la importación automática. \\
    R3 & Catálogo de módulos      & 7 & 9 & Diseño: si una interfaz estrecha
      cubre los casos reales sin proliferar variantes. \\
    R4 & Política como código     & 8 & 9 & Cobertura de reglas: cuántos modos
      de fallo son expresables como política. \\
    R5 & Malla de pruebas         & 7 & 9 & Coste de ejecución de los entornos
      efímeros y tiempo de la malla. \\
    R6 & Sustitución del motor    & 8 & 9 & Compatibilidad del estado
      existente con la bifurcación abierta. \\
    R7 & Interfaz de autoservicio & 6 & 9 & \textbf{Adopción}: el parámetro
      $\sigma$, que domina el modelo técnico. \\
    R8 & Observabilidad           & 7 & 9 & Instrumentación: si las ocho
      figuras de mérito son medibles sin trabajo manual. \\
    \bottomrule
  \end{tabularx}
\end{table}

El cuadro tiene una lectura que va más allá de su contenido, y confirma por
tercera vía el diagnóstico de la sección~\ref{sec:evol}. Ningún proyecto de la
cartera parte por debajo del nivel 6, y seis de los ocho parten del 7 o
superior. No hay investigación en esta cartera: no hay ni un solo componente
cuyo funcionamiento esté por demostrar. Es exactamente lo que cabe esperar de
una tecnología cuya curva S pasó su inflexión hace catorce años, y es la razón
por la que la columna de la derecha no contiene ninguna incertidumbre de
física ni de algoritmo. Las incertidumbres son de volumen, de diseño de
interfaz, de coste y, en el caso decisivo, de adopción humana.

\subsec{02. Los cinco prototipos}

Un prototipo, en esta cartera, no sirve para demostrar que algo funciona: se
sabe que funciona. Sirve para \emph{medir un parámetro} del modelo técnico de
la sección~\ref{sec:tecnico} antes de que el presupuesto dependa de su valor
supuesto. El Cuadro~\ref{tab:proto} los define con ese criterio, y por eso
cada fila declara qué parámetro mide y qué decisión desbloquea.

\begin{table}[htbp]
  \centering
  \footnotesize
  \setlength{\tabcolsep}{4pt}
  \caption{Prototipos de la cartera. Cada uno mide un parámetro del modelo
  técnico y habilita una puerta de decisión. La columna de esfuerzo está
  incluida en la del proyecto al que pertenece en el
  Cuadro~\ref{tab:inversion}.}
  \label{tab:proto}
  \begin{tabularx}{\textwidth}{@{}ll>{\raggedright\arraybackslash}X ccl@{}}
    \toprule
    \textbf{Cód.} & \textbf{Proy.} & \textbf{Prototipo y pregunta que
      responde} & \textbf{Param.} & \textbf{M-p.} & \textbf{Trim.} \\
    \midrule
    P1 & R1 & Importador masivo sobre una cuenta real de tamaño medio.
      ¿Qué fracción de los recursos creados a mano admite importación
      automática y cuánto cuesta el resto? & $\tau_m$ & 2 & 2026 T3 \\
    P2 & R3 & Un módulo certificado de referencia, con contrato de interfaz y
      tres consumidores reales. ¿Una interfaz estrecha cubre los casos sin
      proliferar variantes? & $\kappa$, $\theta$ & 3 & 2027 T1 \\
    P3 & R4 & Evaluación de políticas sobre el plan en la canalización, en
      modo de sólo aviso. ¿Cuántos modos de fallo son expresables y cuál es la
      tasa de falso positivo? & $\pi$, $\gamma$ & 2 & 2027 T4 \\
    P4 & R7 & Portal con tres plantillas y diez equipos de producto
      voluntarios. \textbf{¿Qué fracción de las peticiones se resuelve sin
      entrar en la cola?} & $\sigma$ & 3 & 2028 T2 \\
    P5 & R1, R2, R3 & Reconstrucción completa de un entorno equivalente al de
      producción en una región secundaria, cronometrada. ¿Cuánto vale de
      verdad el tiempo de reconstrucción? & $T_{\text{auto}}$, $\tau_m$ & 2 &
      2028 T4 \\
    \bottomrule
  \end{tabularx}
\end{table}

P4 merece el resaltado que lleva. El análisis de sensibilidad de la
Figura~\ref{fig:tornado} identificó $\sigma$ como el parámetro que domina el
resultado del plan, con una amplitud de 4,8 horas sobre un tiempo de ciclo
objetivo de 8,0, y advirtió que no es un parámetro técnico sino una medida de
comportamiento de los equipos de producto. P4 existe únicamente para medirlo
con diez equipos reales y tres meses-persona, antes de comprometer los 126,7
miles de euros de R7. Es el prototipo con mejor razón entre lo que cuesta y lo
que decide de toda la cartera.

P5 tiene un papel distinto y también singular: no mide un parámetro para
decidir una inversión, sino que verifica de una vez las cuatro propiedades
exigidas en la sección~\ref{sec:alcance}, tal como argumenta la
sección~\ref{sec:fom}. Es la prueba de aceptación de la hoja de ruta completa,
adelantada un año para que quede margen de corrección.

\subsec{03. Puertas de decisión}

El Cuadro~\ref{tab:puertas} define los cuatro puntos en los que la cartera se
somete a decisión explícita. Cada puerta tiene un criterio numérico, no un
juicio, y una consecuencia declarada de antemano para el caso de no
superarse.

\begin{table}[htbp]
  \centering
  \footnotesize
  \setlength{\tabcolsep}{4pt}
  \caption{Puertas de decisión de la cartera. El criterio se evalúa sobre el
  resultado del prototipo indicado.}
  \label{tab:puertas}
  \begin{tabularx}{\textwidth}{@{}lll>{\raggedright\arraybackslash}X@{}}
    \toprule
    \textbf{Puerta} & \textbf{Momento} & \textbf{Prot.} &
      \textbf{Criterio y consecuencia de no superarlo} \\
    \midrule
    G1 & 2026 T4 & P1 & Al menos el 70 \pct{} de los recursos importa de forma
      automática. Si no: se replantea el objetivo de cobertura de 2029 a la
      baja y se recalcula el modelo financiero antes de seguir. \\
    G2 & 2027 T4 & P2, P3 & El módulo de referencia sirve a tres consumidores
      sin bifurcarse, y la política tiene menos del 15 \pct{} de falsos
      positivos. Si no: R4 se degrada a modo de aviso permanente y se
      renuncia al objetivo de cumplimiento del 94 \pct. \\
    G3 & 2028 T3 & P4 & $\sigma \geq 0{,}60$ con diez equipos. \textbf{Si no:
      R7 se cancela} y su presupuesto se reasigna a ampliar el catálogo, con
      un tiempo de ciclo objetivo de 14 h en lugar de 8. \\
    G4 & 2029 T1 & P5 & Reconstrucción completa por debajo de 16 h. Si no: no
      se declara alcanzada la motivación E4 y se abre un proyecto específico
      de continuidad, fuera de esta cartera. \\
    \bottomrule
  \end{tabularx}
\end{table}

G3 es la única puerta que puede cancelar un proyecto entero, y es deliberado.
El escenario de cancelación está cuantificado: sin R7, el tiempo de ciclo se
queda en 14 horas, la motivación E1 se atiende sólo en parte y el valor actual
neto del plan cae, pero sigue siendo positivo, porque las tres corrientes de
beneficio principales dependen de la cobertura y no del autoservicio. La
cartera admite perder su proyecto más visible sin dejar de ser rentable, y
saberlo por adelantado es lo que permite comprometerla.

\subsec{04. Calendario de la cartera}

La Figura~\ref{fig:gantt} sitúa los ocho proyectos, los cinco prototipos y las
cuatro puertas en el tiempo. El orden es el de las seis olas que produjo el
particionado de la DSM en la sección~\ref{sec:dsm}, y no una programación
independiente.

\begin{figure}[ht]
  \centering
  \begin{tikzpicture}[x=0.855cm, y=0.60cm, font=\scriptsize]
    % --- rejilla de trimestres y años
    \foreach \i in {0,...,14} \draw[gray!22] (\i,0.55) -- (\i,-9.9);
    \foreach \i in {0,4,8,12,14} \draw[gray!60, thick] (\i,0.55) -- (\i,-9.9);
    \draw[gray!60, thick] (0,-8.7) -- (14,-8.7);
    \foreach \i/\a in {2/2026, 6/2027, 10/2028, 13/2029}
      \node[anchor=south, font=\scriptsize\bfseries] at (\i,0.62) {\a};
    \foreach \i/\t in {0.5/T1, 1.5/T2, 2.5/T3, 3.5/T4, 4.5/T1, 5.5/T2,
                       6.5/T3, 7.5/T4, 8.5/T1, 9.5/T2, 10.5/T3, 11.5/T4,
                       12.5/T1, 13.5/T2}
      \node[anchor=north, text=gray!50!black, font=\tiny] at (\i,-9.9) {\t};
    % --- barras de proyecto: inicio / fin / fila / color
    \foreach \ini/\fin/\fila/\col in {
      0/2/1/capa, 1.6/4.4/2/capa, 4/7.6/3/trans, 4.4/6.4/4/trans,
      8/10/5/acento, 8.4/10.8/6/acento, 10.4/12/7/jugA, 11.6/14/8/propia}
      \fill[\col!24, draw=\col, thick, rounded corners=1.5pt]
        (\ini,-\fila+0.26) rectangle (\fin,-\fila-0.26);
    % --- columna de nombres, a la izquierda de la rejilla
    \foreach \cod/\nom/\fila in {
      R2/{Custodia del estado}/1, R1/{Cierre de cobertura}/2,
      R3/{Catálogo de módulos}/3, R8/{Observabilidad}/4,
      R4/{Política como código}/5, R5/{Malla de pruebas}/6,
      R6/{Sustitución del motor}/7, R7/{Interfaz de autoservicio}/8} {
      \node[anchor=east, font=\scriptsize\bfseries] at (-3.75,-\fila) {\cod};
      \node[anchor=west, font=\tiny, text=gray!25!black]
        at (-3.64,-\fila) {\nom};
    }
    % --- carril de prototipos, bajo la línea gruesa
    \node[anchor=east, font=\scriptsize\bfseries] at (-3.75,-9.3) {Prot.};
    \foreach \p/\x in {P1/2.5, P2/4.5, P4/9.5, P5/11.5} {
      \node[star, star points=5, star point ratio=2.2, fill=acento,
            draw=acento!60!black, inner sep=1.5pt] at (\x,-9.3) {};
      \node[anchor=west, font=\tiny\bfseries, text=acento!70!black]
        at (\x+0.22,-9.3) {\p};
    }
    % P3 rotula a la izquierda: a su derecha pasa la línea de la puerta G2
    \node[star, star points=5, star point ratio=2.2, fill=acento,
          draw=acento!60!black, inner sep=1.5pt] at (7.5,-9.3) {};
    \node[anchor=east, font=\tiny\bfseries, text=acento!70!black]
      at (7.28,-9.3) {P3};
    % líneas de conducción del prototipo a la fila del proyecto que lo ejecuta
    \foreach \x/\fila in {2.5/2, 4.5/3, 7.5/5, 9.5/8, 11.5/3}
      \draw[acento!55, thin, densely dotted] (\x,-9.05) -- (\x,-\fila-0.3);
    % --- puertas de decisión
    \foreach \g/\x in {G1/3.9, G2/7.9, G3/10.9, G4/12.4} {
      \draw[gray!25!black, very thick, densely dashed] (\x,1.32) -- (\x,-9.8);
      \node[fill=white, draw=gray!25!black, thick, rounded corners=1.5pt,
            inner sep=1.7pt, font=\tiny\bfseries] at (\x,1.5) {\g};
    }
  \end{tikzpicture}
  \caption{Calendario de la cartera. Barras: los ocho proyectos, coloreados
  por ola de la secuencia de la Figura~\ref{fig:dsmproyord}. Carril inferior:
  los cinco prototipos, con línea punteada hasta el proyecto que los ejecuta.
  Líneas discontinuas verticales: las cuatro puertas de decisión. Se aprecia
  el adelanto de P3 y P4 respecto de R4 y R7, que es lo que permite que G2 y
  G3 decidan con medida. Elaboración propia.}
  \label{fig:gantt}
\end{figure}

Dos rasgos del calendario merecen justificación, porque ninguno es la opción
natural.

\textbf{R1 no se cierra: continúa después de 2029.} El apartado 08 de la
sección~\ref{sec:evol} estableció por qué. Si la superficie declarable se
duplica cada 1,8 a 2,2 años y el parque crece, mantener la cobertura en el
86 \% exige trabajo recurrente. La cobertura no es un proyecto que termina
sino un nivel que se sostiene, y por eso el modelo financiero mantiene su
corriente de beneficio en 2030 y 2031 con la cobertura constante en el 88 \% y
no creciendo.

\textbf{Los prototipos se adelantan a su proyecto, no lo abren.} P2 se ejecuta
en el primer trimestre de 2027 mientras R3 arranca en el mismo periodo, pero
P3 y P4 se adelantan un año completo a R4 y a R7 respectivamente. Es lo que
permite que las puertas G2 y G3 decidan con medida y no con estimación, y es
la única razón por la que la cartera puede cancelar R7 en 2028 sin haber
gastado sus 126,7 miles de euros.

% ==========================================================================
\section{Publicaciones, presentaciones y patentes clave}
\label{sec:publica}
% ==========================================================================

El elemento 10 del índice pide identificar las publicaciones, presentaciones y
patentes que determinan el estado del arte de la tecnología. En una hoja de
ruta de tecnología emergente su función es detectar de dónde vendrá el
siguiente salto. Aquí la función es la contraria, y hay que decirlo antes de
las tablas: el diagnóstico de la sección~\ref{sec:evol} sitúa la tecnología en
el tercio superior de su curva S, de modo que lo esperable es un panorama
documental maduro, sin frentes abiertos, y lo que este elemento tiene que
comprobar es precisamente eso. Un hallazgo de investigación básica activa
sería, en este caso, una señal de que el diagnóstico está equivocado.

\subsec{01. Publicaciones fundacionales}

El Cuadro~\ref{tab:pubfund} recoge las publicaciones que fundan cada una de
las cuatro propiedades exigidas en la sección~\ref{sec:alcance}. La selección
es estricta: sólo entran los trabajos que introducen una propiedad, no los que
la aplican.

\begin{table}[htbp]
  \centering
  \footnotesize
  \setlength{\tabcolsep}{4pt}
  \caption{Publicaciones que fundan las propiedades de la tecnología.}
  \label{tab:pubfund}
  \begin{tabularx}{\textwidth}{@{}ll>{\raggedright\arraybackslash}X@{}}
    \toprule
    \textbf{Año} & \textbf{Referencia} & \textbf{Qué funda} \\
    \midrule
    1995 & \textcite{burgess1995} & La especificación externalizada y la
      convergencia: la configuración vive en un fichero central y el sistema
      se repara hacia ella sin intervención. \\
    1998 & \textcite{traugott1998} & La infraestructura como un único sistema
      distribuido y no como un conjunto de máquinas. Es el origen conceptual
      del término. \\
    2002 & \textcite{traugott2002} & La reproducibilidad: demuestra que el
      orden de aplicación importa y que sin orden determinista no hay estado
      reproducible. \\
    2010 & \textcite{delaet2010} & El primer marco comparativo sistemático de
      herramientas de configuración, sobre once productos. \\
    2013 & \textcite{hummer2013} & La idempotencia como propiedad verificable
      mediante pruebas basadas en modelo, y no como promesa del fabricante. \\
    2016 & \textcite{burns2016} & La conciliación continua como arquitectura
      de control, en la genealogía que va de Borg a Kubernetes. \\
    2018 & \textcite{forsgren2018} & Las cuatro figuras de mérito de entrega,
      con evidencia estadística de su relación con el rendimiento
      organizativo. \\
    2021 & \textcite{morris2021} & La síntesis de referencia de la práctica,
      y la fuente de las cuatro propiedades adoptadas en este documento. \\
    \bottomrule
  \end{tabularx}
\end{table}

\subsec{02. Estado de la literatura empírica}

El Cuadro~\ref{tab:pubemp} recoge los estudios que miden la práctica, y su
lectura conjunta es más informativa que cualquiera de ellos por separado.

\begin{table}[htbp]
  \centering
  \footnotesize
  \setlength{\tabcolsep}{4pt}
  \caption{Estudios empíricos sobre la práctica de la infraestructura como
  código.}
  \label{tab:pubemp}
  \begin{tabularx}{\textwidth}{@{}ll>{\raggedright\arraybackslash}X@{}}
    \toprule
    \textbf{Año} & \textbf{Referencia} & \textbf{Objeto y resultado} \\
    \midrule
    2016 & \textcite{sharma2016} & Olores de diseño en código de
      configuración, sobre repositorios públicos. Establece que el código de
      infraestructura se degrada como cualquier otro código. \\
    2019 & \textcite{rahman2019} & Mapeo sistemático de la investigación
      publicada. Concluye que la literatura se concentra en la calidad y los
      defectos del código, y deja fuera casi todo lo demás. \\
    2019 & \textcite{rahman2019smells} & Siete categorías de olores de
      seguridad recurrentes, con secretos en claro a la cabeza. Es el trabajo
      que justifica que la gestión de secretos aparezca en la segunda ola de
      la secuencia de la sección~\ref{sec:dsm}. \\
    2019 & \textcite{guerriero2019} & Cuarenta y cuatro entrevistas en
      empresa. El mantenimiento del propio código de infraestructura es el
      problema declarado con más frecuencia. \\
    2024 & \textcite{dora2024} & Serie anual de percentiles de las cuatro
      figuras de entrega, que es la única fuente pública comparable entre años
      y la que ancla las series de la sección~\ref{sec:evol}. \\
    \bottomrule
  \end{tabularx}
\end{table}

De aquí se extraen dos conclusiones que afectan a la hoja de ruta.

\textbf{La literatura confirma el diagnóstico de saturación.} Ninguno de los
trabajos posteriores a 2016 introduce una propiedad nueva de la tecnología:
todos miden, comparan o critican la práctica de propiedades ya establecidas
entre 1995 y 2013. El mapeo sistemático lo dice de forma expresa al concluir
que la investigación se concentra en la calidad del código y no en las
capacidades del sistema \parencite{rahman2019}. No hay frente abierto de
investigación básica del que pueda venir un salto que invalide esta cartera.

\textbf{El problema que la literatura señala como principal es el que esta
cartera atiende.} El estudio de campo con cuarenta y cuatro empresas
identifica el mantenimiento del código de infraestructura como la dificultad
declarada con más frecuencia \parencite{guerriero2019}, y es exactamente lo
que atacan R3, el catálogo de módulos certificados, y R5, la malla de pruebas.
La coincidencia no valida la cartera, pero sí indica que no se ha construido
sobre un problema imaginado.

\subsec{03. Patentes}

El Cuadro~\ref{tab:patentes} recoge tres patentes concedidas o publicadas
cuyas reivindicaciones alcanzan de forma directa a los componentes de esta
hoja de ruta. Se han verificado individualmente en el registro, y se citan con
su número, su titular y su fecha para que puedan comprobarse.

\begin{table}[htbp]
  \centering
  \footnotesize
  \setlength{\tabcolsep}{4pt}
  \caption{Patentes verificadas con reivindicaciones sobre componentes de esta
  hoja de ruta.}
  \label{tab:patentes}
  \begin{tabularx}{\textwidth}{@{}l>{\raggedright\arraybackslash}p{0.20\textwidth}
    >{\raggedright\arraybackslash}X l@{}}
    \toprule
    \textbf{Número} & \textbf{Titular y fecha} & \textbf{Objeto de la
      reivindicación} & \textbf{Afecta a} \\
    \midrule
    US 11.372.626 B2 \parencite{us11372626} & JPMorgan Chase Bank, junio de
      2022 & Marco de creación, empaquetado, versionado y despliegue de
      plantillas de infraestructura como código a través de un repositorio
      central, con actualización y verificación automatizadas del paquete. &
      R3 \\
    US 11.349.958 B1 \parencite{us11349958} & Salesforce, mayo de 2022 &
      Lenguaje específico de dominio que declara un centro de datos y compila
      a instrucciones de cada proveedor de nube, con canalizaciones
      unificadas. & R6 \\
    WO 2024/023685 A1 \parencite{wo2024023685} & Dazz, febrero de 2024 &
      Detección de configuraciones incorrectas transformando las definiciones
      propias de cada proveedor en definiciones universales sobre las que se
      evalúan políticas antes del despliegue. & R4 \\
    \bottomrule
  \end{tabularx}
\end{table}

\subsec{04. Lectura del panorama de patentes}

Tres observaciones, y la tercera es la que importa para la decisión.

\textbf{Los titulares no son los fabricantes de herramientas.} Dos de las tres
patentes pertenecen a empresas usuarias, un banco y un proveedor de programas
de gestión, que patentaron la manera de organizar su propia plataforma
interna. Es un indicio del estado de la tecnología: cuando el valor se
desplaza del motor al gobierno de lo que el motor aplica, quien patenta es
quien opera y no quien vende.

\textbf{Las tres reivindicaciones caen sobre proyectos de esta cartera y
ninguna sobre el motor.} Afectan al catálogo empaquetado (R3), a la
abstracción sobre varios proveedores (R6) y a la evaluación previa de
políticas (R4). El motor de aprovisionamiento, que es la pieza sobre la que
gira el debate técnico habitual, no aparece: está en dominio público y
saturado, lo que coincide una vez más con el diagnóstico de la
sección~\ref{sec:evol}.

\textbf{Ninguna de las tres constituye un impedimento, pero la tercera exige
vigilancia.} Las dos primeras describen arquitecturas cuyas piezas están
disponibles como programas libres de uso extendido y anterior, de modo que la
implantación descrita en esta hoja de ruta no incorpora nada que dependa de
ellas. La solicitud internacional de 2024 es distinta: reivindica la
normalización de definiciones de distintos proveedores a un formato universal
sobre el que evaluar políticas, y eso es aproximadamente lo que R4 hace si el
alcance se extiende a más de un proveedor de nube. La recomendación
consiguiente, que se traslada a la sección~\ref{sec:estrategia}, es doble:
construir R4 sobre el motor de políticas de código abierto ya citado
\parencite{opa}, cuya anterioridad es pública y documentada, y encargar una
revisión de libertad de operación antes de ampliar R4 a un segundo proveedor,
no antes de construirlo.

\subsec{05. Plan de publicación y presentación}

Este elemento del índice no se agota en leer lo publicado por otros: pide
también qué publicará y presentará la empresa. El Cuadro~\ref{tab:planpub} lo
concreta, con un criterio de selección explícito: se publica lo que la empresa
mide y nadie más publica, y no se publica lo que otros ya han publicado mejor.

\begin{table}[htbp]
  \centering
  \footnotesize
  \setlength{\tabcolsep}{4pt}
  \caption{Plan de publicación y presentación asociado a la hoja de ruta.}
  \label{tab:planpub}
  \begin{tabularx}{\textwidth}{@{}ll>{\raggedright\arraybackslash}X@{}}
    \toprule
    \textbf{Momento} & \textbf{Foro} & \textbf{Contenido y por qué es
      publicable} \\
    \midrule
    2027 T1 & Interno, comité de tecnología & Resultado de P1: la fracción
      real de recursos que admite importación automática. Ninguna fuente
      pública la publica, y es el dato que decide la puerta G1. \\
    2027 T4 & Interno, todos los equipos & Contrato de interfaz del catálogo
      de módulos. Su publicación interna es condición de la adopción, no un
      acto de difusión. \\
    2028 T4 & Externo, conferencia de ingeniería de plataforma & Medida de
      $\sigma$ con diez equipos (P4): qué fracción de las peticiones de
      infraestructura se resuelve por autoservicio y qué la limita. Es el
      parámetro que domina el modelo técnico y no existe medida publicada. \\
    2029 T2 & Externo, artículo técnico & Serie completa de las ocho figuras
      de mérito medidas durante tres años en una misma organización, con el
      procedimiento del Cuadro~\ref{tab:fom}. Cierra el hueco que el mapeo
      sistemático de \textcite{rahman2019} identifica: hay literatura sobre la
      calidad del código y casi ninguna sobre las capacidades medidas del
      sistema. \\
    \midrule
    \multicolumn{3}{@{}l}{\textit{No se solicitarán patentes.} Los tres
      componentes patentables de la cartera (catálogo, abstracción} \\
    \multicolumn{3}{@{}l}{multiproveedor y política previa) tienen ya
      anterioridad de terceros, y una patente defensiva no aporta} \\
    \multicolumn{3}{@{}l}{capacidad de negociación a una organización que no
      vende la tecnología.} \\
    \bottomrule
  \end{tabularx}
\end{table}

La publicación de 2029 T2 es la única con valor externo real, y su valor
procede de una asimetría: la literatura mide código y la práctica necesita
medir capacidades. Una serie de tres años de las ocho figuras de mérito en una
misma organización, con procedimiento declarado, es un dato que hoy no existe
publicado, y la empresa lo tendrá por el simple hecho de haber instrumentado
R8 en 2027.

% ==========================================================================
\section{Declaración de estrategia tecnológica}
\label{sec:estrategia}
% ==========================================================================

El elemento 11 del índice pide la declaración de estrategia tecnológica: el
documento breve que la dirección aprueba y contra el que se juzgan después las
decisiones. Todo lo anterior es su justificación; esta sección es la
declaración misma, y su valor depende de que sea corta, comprometida y
falsable. Se redacta por eso en tres piezas: la postura, los compromisos
numéricos y las renuncias explícitas.

\subsec{01. La postura}

\begin{quote}
\small
La infraestructura como código es para esta empresa una \textbf{capacidad de
paridad, no de diferenciación}. Su curva S pasó el punto de inflexión en 2012
y ha recorrido el 92 \% de su límite: la frontera de la tecnología avanzará
tres puntos de cobertura en los seis años del plan, mientras que la distancia
que separa a la empresa de esa frontera es de treinta y siete. La estrategia
se sigue de esa asimetría, y consiste en \textbf{cerrar la distancia con la
tecnología ya disponible, no en apostar por la que viene}.

En consecuencia, la empresa \textbf{no competirá en la elección de motor de
aprovisionamiento}, que es la parte saturada y en dominio público de la
tecnología, y que explica 1,9 de los 46,9 puntos de mejora disponibles.
Concentrará la inversión en las tres capacidades propias que explican los
otros 45: \textbf{la cobertura de codificación del parque, el catálogo de
módulos certificados y la interfaz de autoservicio}. Las tres son
independientes del proveedor y ninguna se puede comprar.
\end{quote}

\subsec{02. Los compromisos}

La declaración se compromete con las ocho cifras del
Cuadro~\ref{tab:anclas}, y de forma señalada con cinco de ellas, que son las
que la dirección revisará. El Cuadro~\ref{tab:compromisos} las recoge con su
responsable y su instrumento de verificación, sin los cuales un compromiso no
es más que una intención.

\begin{table}[htbp]
  \centering
  \footnotesize
  \setlength{\tabcolsep}{4pt}
  \caption{Compromisos numéricos de la declaración de estrategia. La columna
  de verificación indica con qué instrumento se comprobará cada cifra, que en
  todos los casos es el proyecto R8.}
  \label{tab:compromisos}
  \begin{tabularx}{\textwidth}{@{}l>{\raggedright\arraybackslash}p{0.19\textwidth}rrl
    >{\raggedright\arraybackslash}X@{}}
    \toprule
    \textbf{N.\textsuperscript{o}} & \textbf{Compromiso} & \textbf{2026} &
      \textbf{2029} & \textbf{Motiv.} & \textbf{Verificación} \\
    \midrule
    C1 & Cobertura de codificación (\pct) & 55 & 86 & E1, E2, E3, E4, E5 &
      Inventario del proveedor contra estado del motor, semanal. \\
    C2 & Tiempo de reconstrucción (h) & 40 & 11 & E4, E5 & Ejercicio
      cronometrado en región secundaria, semestral (P5). \\
    C3 & Tiempo de ciclo del cambio (h) & 26 & 8 & E1, E6 & Marca de tiempo
      del repositorio a la aplicación, continua. \\
    C4 & Cumplimiento de política (\pct) & 62 & 94 & E3 & Recuento de planes
      con excepción manual, continuo. \\
    C5 & Tasa de deriva (rec./1000 sem.) & 38 & 8 & E2, E3, E6 & Ejecuciones
      de sólo lectura, semanal. \\
    \bottomrule
  \end{tabularx}
\end{table}

C1 es el compromiso principal, y el orden no es casual. La cobertura de
codificación es indicador principal de cinco de las seis motivaciones
estratégicas (Cuadro~\ref{tab:motivfom}), la variable que más mueve el valor
actual neto del plan (Cuadro~\ref{tab:sensfin}), el 90 \% del tiempo de
reconstrucción de 2026 (sección~\ref{sec:tecnico}) y la figura de mérito en la
que la empresa está más lejos de la frontera (sección~\ref{sec:evol}). Cuatro
análisis independientes la señalan, y la declaración de estrategia se limita a
reconocerlo: \textbf{si sólo puede cumplirse un compromiso de los cinco, es
C1}.

\subsec{03. Decisiones de construir, comprar o asociarse}

El Cuadro~\ref{tab:makebuy} aplica la postura a cada capa del sistema. El
criterio se deriva del diagnóstico de la curva S y es único: se compra o se
adopta lo que está saturado y en dominio público, y se construye lo que
depende del contexto propio de la empresa y por tanto nadie puede vender.

\begin{table}[htbp]
  \centering
  \footnotesize
  \setlength{\tabcolsep}{4pt}
  \caption{Decisión de construir, comprar o adoptar por componente, y su
  justificación desde el diagnóstico de la sección~\ref{sec:evol}.}
  \label{tab:makebuy}
  \begin{tabularx}{\textwidth}{@{}ll>{\raggedright\arraybackslash}X@{}}
    \toprule
    \textbf{Componente} & \textbf{Decisión} & \textbf{Por qué} \\
    \midrule
    Motor de aprovisionamiento & Adoptar & Saturado y en dominio público. Se
      adopta la bifurcación abierta \parencite{opentofuga} por la motivación
      E5, no por capacidad técnica. \\
    Registro de estado & Comprar & Servicio gestionado. Es riesgo puro sin
      diferenciación posible. \\
    Motor de políticas & Adoptar & Programa libre con anterioridad pública
      documentada \parencite{opa}, lo que además atiende la reserva de la
      sección~\ref{sec:publica} sobre la solicitud WO 2024/023685. \\
    Reglas de política & \textbf{Construir} & Son las normas de esta
      organización. No existen fuera de ella. \\
    Catálogo de módulos & \textbf{Construir} & Codifica las decisiones de
      arquitectura propias. Es la pieza que atrapa el 68,8 \% de los modos de
      fallo (sección~\ref{sec:tecnico}). \\
    Malla de pruebas & Adoptar y construir & Herramienta adoptada, casos de
      prueba propios. \\
    Interfaz de autoservicio & Adoptar y construir & Armazón de código
      abierto \parencite{backstage}, plantillas propias sobre el catálogo. \\
    Observabilidad & Comprar y construir & Plataforma comprada,
      instrumentación de las ocho figuras de mérito construida. Nadie la
      vende. \\
    \bottomrule
  \end{tabularx}
\end{table}

El cuadro se reduce a una frase: se construye el catálogo, las reglas, las
plantillas y las medidas; se adopta o se compra todo lo demás. Las cuatro
piezas que se construyen tienen en común que son las que expresan
conocimiento de esta organización, y ninguna de ellas es un motor.

\subsec{04. Las renuncias}

Una estrategia sin renuncias declaradas no es una estrategia. Estas cuatro
quedan asumidas de forma expresa, con la consecuencia que cada una acarrea.

\begin{enumerate}
  \item \textbf{No se buscará la frontera en las figuras de entrega.} El
    objetivo de 8 horas de tiempo de ciclo está por detrás de las 4 horas
    alcanzables (Cuadro~\ref{tab:anclas}). Los últimos puntos de valor cuestan
    catorce veces más que los anteriores (Cuadro~\ref{tab:pareto}), y la
    empresa se detiene en el codo del frente de Pareto de forma deliberada.
  \item \textbf{No se adoptará un lenguaje de propósito general.} El
    jugador C ofrece 3,4 puntos más de valor por 48 puntos más de coste. Es
    una posición eficiente del frente y se descarta igualmente, porque el
    canje no compensa.
  \item \textbf{No se solicitarán patentes}, por las razones del apartado 05
    de la sección~\ref{sec:publica}: hay anterioridad de terceros en los tres
    componentes patentables y una patente defensiva no aporta capacidad de
    negociación a quien no vende la tecnología.
  \item \textbf{No se compromete un segundo proveedor de nube.} La empresa
    conserva la \emph{capacidad} de abrirlo, que es lo que la motivación E5
    exige, pero no lo planifica. La sensibilidad del
    Cuadro~\ref{tab:sensib} muestra que la ruta propia sólo pierde 1,9 puntos
    de valor si esa cartera cambia, de modo que la opción se mantiene abierta
    a coste casi nulo. Antes de ejercerla habrá que encargar la revisión de
    libertad de operación sobre R4 que recomienda la
    sección~\ref{sec:publica}.
\end{enumerate}

\subsec{05. Condiciones que invalidarían esta estrategia}

La declaración es falsable, y estas son las tres observaciones que obligarían
a rehacerla. Se enuncian con su umbral para que no dependan de interpretación.

\begin{enumerate}
  \item \textbf{Que la curva S se reanude.} Si apareciera una tecnología que
    llevara el tiempo de reconstrucción por debajo de una hora, esto es, un
    orden de magnitud por debajo de la asíntota de 3,6 horas estimada en la
    sección~\ref{sec:evol}, el diagnóstico de saturación quedaría refutado y
    con él la postura de paridad. La vigilancia de esta condición es el
    propósito de mantener el elemento 10 al día.
  \item \textbf{Que $\sigma$ no llegue a 0,60 en la puerta G3.} Sería la
    señal de que el valor no está donde este documento lo sitúa. La
    consecuencia está preparada y cuantificada en la
    sección~\ref{sec:portfolio}: R7 se cancela y el tiempo de ciclo objetivo
    pasa de 8 a 14 horas.
  \item \textbf{Que la cobertura se estanque por debajo del 74 \%.} Es el
    escenario que destruye 434,6 miles de euros de valor actual neto
    (Cuadro~\ref{tab:sensfin}) y el único que pone en cuestión el caso de
    inversión completo. La puerta G1, en el cuarto trimestre de 2026, existe
    para detectarlo en el primer año y no en el tercero.
\end{enumerate}

% ==========================================================================
\section{Evaluación del nivel de madurez de la hoja de ruta}
\label{sec:madurez}
% ==========================================================================

El elemento 12 cierra el índice pidiendo que la hoja de ruta se evalúe a sí
misma. Es el elemento más fácil de convertir en autoelogio, y para evitarlo se
puntúa con un criterio fijado antes de mirar el documento y se separa en dos
dimensiones que en este caso dan resultados muy distintos.

\subsec{01. Criterio de evaluación}

Cada uno de los doce elementos se puntúa de 0 a 4 en dos dimensiones
independientes:

\begin{description}[leftmargin=1.6em, style=nextline]
  \item[Completitud] Si el elemento dispone del instrumento que el método
    exige, y completo. 0, ausente; 1, mencionado; 2, con el instrumento
    incompleto; 3, con el instrumento completo pero con huecos declarados; 4,
    completo.
  \item[Verificabilidad] De dónde proceden sus cifras. 0, sin cifras; 1,
    supuestos del autor sin fuente; 2, estimaciones razonadas con supuestos
    declarados; 3, cifras derivadas de fuentes publicadas; 4, cifras
    procedentes de medida instrumentada o de fuente verificada
    individualmente.
\end{description}

La separación es lo que hace útil el ejercicio: un elemento puede estar
completo y no ser verificable, que es la situación dominante en este
documento.

\begin{table}[htbp]
  \centering
  \footnotesize
  \setlength{\tabcolsep}{4pt}
  \caption{Evaluación de los doce elementos en las dos dimensiones. C,
  completitud; V, verificabilidad.}
  \label{tab:madurez}
  \begin{tabularx}{\textwidth}{@{}llcc>{\raggedright\arraybackslash}X@{}}
    \toprule
    \textbf{N.\textsuperscript{o}} & \textbf{Elemento} & \textbf{C} &
      \textbf{V} & \textbf{Justificación de la puntuación} \\
    \midrule
    1 & Descripción & 4 & 3 & Siete secciones, cuatro propiedades y frontera
      declarada, con apoyo en la literatura de referencia. \\
    2 & DSM & 4 & 2 & Dos matrices y particionado completo, pero las
      dependencias las declara el autor y no proceden de medida. \\
    3 & OPM & 3 & 3 & Tres diagramas y su OPL, conformes al vocabulario de la
      norma. Falta la composición con la herramienta oficial, que se entrega
      aparte. \\
    4 & Figuras de mérito & 4 & 2 & Ocho figuras con unidad, procedimiento,
      anclas, ajustes de curva S y tasas de mejora. Cuatro de las ocho series
      históricas son estimaciones de ingeniería. \\
    5 & Motivaciones estratégicas & 4 & 1 & Matriz completa y trazabilidad
      hasta los proyectos, pero las seis motivaciones las enuncia el autor sin
      documento de dirección que las respalde. \\
    6 & Competencia & 4 & 2 & Gráfico vectorial, tres competidores, análisis
      de sensibilidad y frente de Pareto. Coordenadas estimadas. \\
    7 & Modelo técnico & 3 & 1 & Cuatro ecuaciones calibradas que reproducen
      los objetivos. No modela dos de las cuatro figuras de entrega, y sus
      coeficientes se ajustan sobre extremos estimados. \\
    8 & Modelo financiero & 3 & 1 & Flujo, valor actual neto, tasa interna de
      retorno y siete escenarios. Los tres precios unitarios son supuestos y
      falta el valor terminal. \\
    9 & Cartera y prototipos & 4 & 2 & Ocho proyectos con madurez de partida,
      cinco prototipos con el parámetro que miden y cuatro puertas con
      criterio numérico. \\
    10 & Publicaciones y patentes & 4 & 4 & Trece publicaciones y tres
      patentes verificadas una a una en el registro, con número, titular y
      fecha. \\
    11 & Estrategia & 4 & 2 & Postura, cinco compromisos, ocho decisiones de
      construir o comprar, cuatro renuncias y tres condiciones de
      invalidación. \\
    12 & Madurez & 3 & 2 & Esta sección. Es una autoevaluación sin revisión
      externa, que es su límite principal. \\
    \midrule
    \multicolumn{2}{@{}l}{\textbf{Media}} & \textbf{3{,}67} &
      \textbf{2{,}08} & \\
    \bottomrule
  \end{tabularx}
\end{table}

\begin{figure}[ht]
  \centering
  \begin{tikzpicture}
    \begin{axis}[
      width=15.2cm, height=6.4cm,
      ybar=1pt, bar width=8.5pt,
      ymin=0, ymax=4.9,
      ylabel={\footnotesize Puntuación (0 a 4)},
      ylabel style={font=\footnotesize},
      xlabel={\footnotesize Elemento del índice},
      xlabel style={font=\footnotesize},
      symbolic x coords={1,2,3,4,5,6,7,8,9,10,11,12},
      xtick={1,2,3,4,5,6,7,8,9,10,11,12},
      ytick={0,1,2,3,4},
      xticklabel style={font=\scriptsize},
      yticklabel style={font=\scriptsize},
      ymajorgrids, grid style={gray!20}, axis lines=left,
      enlarge x limits=0.05,
      legend style={font=\scriptsize, at={(0.5,1.03)}, anchor=south,
                    legend columns=2, draw=none, column sep=1.4ex},
    ]
      \addplot[fill=capa!35, draw=capa] coordinates {
        (1,4) (2,4) (3,3) (4,4) (5,4) (6,4) (7,3) (8,3) (9,4) (10,4)
        (11,4) (12,3)};
      \addlegendentry{Completitud (media 3,67)}
      \addplot[fill=acento!30, draw=acento] coordinates {
        (1,3) (2,2) (3,3) (4,2) (5,1) (6,2) (7,1) (8,1) (9,2) (10,4)
        (11,2) (12,2)};
      \addlegendentry{Verificabilidad (media 2,08)}
      \draw[capa, thick, densely dashed]
        ({rel axis cs:0,0}|-{axis cs:1,3.67}) --
        ({rel axis cs:1,0}|-{axis cs:1,3.67});
      \draw[acento, thick, densely dashed]
        ({rel axis cs:0,0}|-{axis cs:1,2.08}) --
        ({rel axis cs:1,0}|-{axis cs:1,2.08});
    \end{axis}
  \end{tikzpicture}
  \caption{Evaluación de madurez por elemento. La brecha constante entre las
  dos series, y no su altura, es el resultado: el documento está casi completo
  en su forma y débilmente sustentado en sus cifras. El único elemento sin
  brecha es el 10, cuyas fuentes se verificaron una a una. Elaboración
  propia.}
  \label{fig:madurez}
\end{figure}

\subsec{02. Nivel alcanzado}

Sobre la escala de cinco niveles del Cuadro~\ref{tab:niveles}, la hoja de ruta
está en el nivel 3 de 5, y no alcanza el 4.

\begin{table}[htbp]
  \centering
  \footnotesize
  \caption{Niveles de madurez de una hoja de ruta tecnológica y situación de
  esta.}
  \label{tab:niveles}
  \begin{tabularx}{\textwidth}{@{}ll>{\raggedright\arraybackslash}Xl@{}}
    \toprule
    \textbf{Niv.} & \textbf{Nombre} & \textbf{Qué exige} & \textbf{Estado} \\
    \midrule
    1 & Inicial & Elementos dispersos, sin índice común. & Superado \\
    2 & Estructurado & Los doce elementos presentes, cada uno con el
      instrumento que le corresponde. & Superado \\
    3 & Cuantificado & Cifras trazables en todos los elementos, con supuestos
      declarados y sensibilidad calculada. & \textbf{Alcanzado} \\
    4 & Medido & Las cifras proceden de instrumentación y no de estimación.
      Las series históricas son medidas. & No alcanzado \\
    5 & Gobernado & Revisión periódica contra la medida, con registro de las
      decisiones que la desviación provoca. & No alcanzado \\
    \bottomrule
  \end{tabularx}
\end{table}

El nivel 3 está alcanzado sin reservas: los doce elementos tienen cifras,
todas declaran sus supuestos y los tres análisis de sensibilidad
(Figura~\ref{fig:tornado}, Cuadros~\ref{tab:sensib} y~\ref{tab:sensfin})
acotan el efecto de equivocarse en ellos.

El nivel 4 no lo está, y la razón es única: ninguna de las ocho figuras de
mérito está instrumentada hoy. La columna de 2026 del
Cuadro~\ref{tab:anclas} es tan estimada como la de 2029. Esto explica la
brecha de la Figura~\ref{fig:madurez} y la explica \emph{entera}: los ocho
elementos con verificabilidad de 1 o 2 la tienen baja por el mismo motivo, no
por ocho motivos distintos.

\subsec{03. Lo que falta para subir de nivel}

De aquí se sigue la conclusión más útil de esta sección, y es una conclusión
sobre la secuencia de la cartera y no sobre el documento.

\textbf{Un solo proyecto separa el nivel 3 del nivel 4: R8, la
observabilidad de la plataforma.} En cuanto las ocho figuras de mérito se
midan de forma continua, la verificabilidad de los elementos 4, 6, 7 y 8 pasa
de 1 o 2 a 4, porque las cuatro dependen de las mismas series. La media de
verificabilidad subiría de 2,08 a aproximadamente 3,1 con un solo proyecto de
ocho meses-persona. Es la justificación definitiva de una decisión que la
sección~\ref{sec:dsm} había tomado por razones de dependencia estructural:
situar R8 en la tercera ola, en 2027, y no al final del plan. La hoja de ruta
no puede verificarse hasta que su instrumento de medida exista, y ese
instrumento es un proyecto de la propia hoja de ruta.

El nivel 5 exige algo que no es un proyecto sino una práctica: revisar los
cinco compromisos del Cuadro~\ref{tab:compromisos} contra la medida con una
cadencia fija y registrar qué decisión se toma ante cada desviación. Las
cuatro puertas del Cuadro~\ref{tab:puertas} son el embrión de esa práctica,
pero cubren sólo los años del plan y no su régimen permanente.

\subsec{04. Las tres debilidades que no arregla R8}

Conviene terminar por lo que la medida no resuelve, porque es donde este
documento es más débil.

\begin{enumerate}
  \item \textbf{Las seis motivaciones estratégicas no tienen respaldo
    documental.} El elemento 5 es el único con verificabilidad 1 que R8 no
    mejora: las motivaciones no se miden, se aprueban. Mientras no exista un
    acta de dirección que las suscriba, toda la cadena de trazabilidad del
    Cuadro~\ref{tab:trazafom} cuelga de una premisa que el autor ha
    formulado. Es la debilidad más grave del documento.
  \item \textbf{E5 sigue sin indicador propio.} La capacidad de negociación
    con el proveedor no la vigila ninguna de las ocho figuras de mérito, por
    las razones expuestas en el apartado 04 de la sección~\ref{sec:motiv}. No
    hay solución dentro del marco elegido; hay que reconocerlo en lugar de
    inventar una métrica que no resistiría el criterio de admisión.
  \item \textbf{Esta evaluación no tiene revisión externa.} Es el motivo de
    que el propio elemento 12 se puntúe con completitud 3 y no 4. Una
    autoevaluación que se pone un 3,67 en completitud y un 2,08 en
    verificabilidad es más creíble que una que se ponga un 4 en las dos, pero
    sigue siendo una autoevaluación.
\end{enumerate}


% ==========================================================================
\section{Conclusiones}
\label{sec:conclusiones}
% ==========================================================================

\begin{enumerate}
  \item El funcionamiento de un sistema de infraestructura como código se
    reduce a un mecanismo único: la conciliación repetida entre un estado
    deseado, expresado en código versionado, y un estado real, observado a
    través de la interfaz de un proveedor. Todo lo demás, herramientas,
    formatos y flujos de trabajo, son variaciones sobre ese mecanismo.
  \item La pieza que hace posible el mecanismo es el estado registrado. No es
    un detalle de implementación sino el elemento que permite reconocer la
    propiedad de los recursos, detectar destrucciones y encadenar dependencias,
    y por eso concentra los requisitos de cifrado, versionado y bloqueo.
  \item Un sistema completo comprende siete capas, del sustrato del proveedor
    a la interfaz de autoservicio, y seis módulos transversales. Las capas
    intermedias compiten por el mismo territorio, de modo que la primera
    decisión de arquitectura consiste en asignar de forma explícita la
    propiedad de cada característica a una sola capa.
  \item Los módulos transversales, y en particular la política como código y
    las pruebas, son los que distinguen un sistema que automatiza de uno que
    además garantiza. Su ausencia no impide desplegar: hace que los errores se
    desplieguen igual de rápido que los aciertos.
  \item El tiempo de reconstrucción de un entorno completo desde el
    repositorio es la figura de mérito que mejor resume el estado de una
    implantación, porque verifica de una vez las cuatro propiedades que
    definen al sistema.
  \item La matriz de estructura de dependencia desplaza el orden de
    implantación que sugiere la intuición. Los concentradores de riesgo no son
    las capas visibles sino dos módulos transversales, el control de versiones
    con ocho salidas y la gestión de secretos con seis, y en la secuencia
    particionada la gestión de secretos aparece en la segunda ola, antes que
    cualquier capa intermedia. Tratarla como un endurecimiento posterior
    obliga a rehacer todo lo construido antes.
  \item La matriz de la cartera resulta enteramente triangular inferior, de
    modo que existe una secuencia de los ocho proyectos que no obliga a
    rehacer ninguno: custodia del estado, cierre de cobertura, catálogo y
    observabilidad, política y pruebas, sustitución del motor y, en último
    lugar, interfaz de autoservicio. Convierte además en constatación
    estructural lo que la sección~\ref{sec:comp} defiende como criterio: el
    catálogo antes que el motor no es una preferencia sino la única
    ordenación que no fuerza una iteración.
  \item El gráfico vectorial de la sección~\ref{sec:comp} desplaza la
    pregunta habitual.
    La comparación entre Terraform, Bicep y Pulumi ocupa la mayor parte del
    debate técnico y explica menos de la décima parte del valor que se puede
    ganar: en la trayectoria propuesta, el vector que sustituye el motor
    aporta 1,9 de los 46,9 puntos de mejora, y los otros 45 proceden del
    catálogo de módulos certificados y de la interfaz de autoservicio, que son
    inversiones propias e independientes del proveedor elegido.
  \item Lo que separa a los jugadores en el plano de las figuras de mérito no
    es la velocidad ni la fiabilidad, magnitudes en las que las cuatro
    posiciones avanzadas quedan a pocos puntos entre sí, sino la cobertura de
    codificación, esto es, cuánta infraestructura queda dentro del alcance del
    motor. La ventaja de coste del jugador B no es una propiedad de su
    tecnología sino del supuesto de proveedor único: bajo una cartera
    multinube pierde 22,5 puntos de valor y abandona la frontera de
    eficiencia.
  \item El modelado en OPM aporta algo que la descripción en prosa no aporta:
    obliga a declarar el tipo de cada relación. Al hacerlo aparecen con
    claridad la doble condición del registro de estado, la diferencia formal
    entre aplicar, conciliar y retirar, que sólo se sostiene si los enlaces
    parten de estados y no del objeto entero, y el alcance exacto de la
    intervención humana, que en este modelo se limita a dos de los siete
    subprocesos.
  \item Los elementos 2 y 3 del índice se sostienen mutuamente, y ese es el
    hallazgo metodológico del trabajo. La matriz de dependencias detecta un
    único bloque acoplado, el que forman el motor de aprovisionamiento y la
    política como código, y no dispone de instrumentos para romperlo; la
    ampliación en detalle del modelo OPM lo rompe al partir el proceso en
    planificación, validación y aplicación. La matriz encuentra el problema y
    el modelo aporta la partición que lo resuelve.
  \item Las tres formas de mirar la evolución de la tecnología dan
    respuestas divergentes, y su divergencia es el hallazgo de la
    sección~\ref{sec:evol}. Lo que la tecnología entrega, medido como razón o
    como tiempo, está saturando: la cobertura alcanzable pasó su punto de
    inflexión en 2012, ha recorrido el 92 \pct{} de su límite y el tiempo de
    reconstrucción se acerca a una asíntota de 3,6 horas con un tiempo de
    división por dos que se alarga de 4,3 a 5,1 años. Lo que la tecnología
    abarca, medido como recuento, sigue en régimen exponencial: la superficie
    declarable se duplica cada 1,8 a 2,2 años, indistinguible del ritmo de la
    ley de Moore. No hay contradicción, sino la misma magnitud vista desde sus
    dos lados: cada servicio nuevo aparece a la vez como recurso declarable y
    como recurso que hay que declarar, de modo que la razón entre ambos
    satura mientras el recuento explota.
  \item De ese diagnóstico se sigue la estrategia entera, y no de una
    preferencia. Con la frontera avanzando tres puntos de cobertura en seis
    años y la empresa treinta y siete puntos por detrás de ella, el
    competidor relevante es la plataforma propia de 2026 y no Terraform ni
    Pulumi. La infraestructura como código es una capacidad de paridad, no de
    diferenciación, y por eso la sección~\ref{sec:estrategia} renuncia
    explícitamente a competir en la elección de motor.
  \item El frente de Pareto añade dos resultados que el gráfico vectorial no
    hacía explícitos. Terraform gestionado está \emph{estrictamente dominado}
    por la ruta propia, que ofrece 1,4 puntos más de valor por 36 puntos menos
    de coste, de modo que ninguna ponderación de los objetivos lo hace
    preferible; y la plataforma actual está dominada por la misma razón, lo
    que equivale a decir que no actuar no es una posición eficiente. La ruta
    propia queda además en el codo del frente: los últimos 3,4 puntos de valor
    disponibles cuestan 48 puntos de coste, catorce veces más caros que los
    anteriores.
  \item Una misma figura de mérito, la cobertura de codificación, resulta
    señalada por cuatro análisis independientes: es indicador principal de
    cinco de las seis motivaciones estratégicas, la variable que más mueve el
    valor actual neto (434,6 miles de euros entre el 74 y el 86 \pct), el
    90 \pct{} del tiempo de reconstrucción de 2026 y la figura en la que la
    empresa está más lejos de la frontera. La convergencia de cuatro caminos
    distintos sobre la misma magnitud es el argumento más sólido del
    documento, y es lo que justifica que sea el único compromiso que la
    declaración de estrategia señala como irrenunciable.
  \item El modelo técnico produce dos resultados contrarios a la intuición
    técnica habitual. El primero: lo que domina el tiempo de ciclo no es el
    catálogo de módulos sino la espera en la cola de plataforma, de modo que
    el autoservicio aporta el 66,5 \pct{} de la mejora y el catálogo sólo el
    9,0 \pct. El segundo: el catálogo certificado atrapa el 68,8 \pct{} de los
    modos de fallo y la política como código apenas el 35,4 \pct, porque la
    política sólo detecta lo que alguien anticipó mientras un módulo validado
    impide de raíz lo que nadie previó. La política atrapa lo previsto; el
    catálogo elimina lo imprevisto.
  \item El mayor riesgo de una hoja de ruta de infraestructura no está en la
    infraestructura. El análisis de sensibilidad concentra la incertidumbre en
    un solo parámetro, la adopción del autoservicio, con una amplitud de 4,8
    horas sobre un tiempo de ciclo objetivo de 8,0, y ese parámetro mide
    comportamiento de los equipos de producto y no capacidad técnica. La
    respuesta de la cartera es proporcionada: un prototipo de tres
    meses-persona lo mide con diez equipos reales, y una puerta de decisión
    cancela el proyecto de 126,7 miles de euros si no alcanza el umbral.
  \item El plan es financiable y su rentabilidad no depende de acertar en los
    supuestos: 820,2 miles de euros de inversión, un valor actual neto de
    743,3 al 9 \pct, una tasa interna de retorno del 51,8 \pct{} y retorno en
    2030, con valor actual neto positivo en los siete escenarios de
    sensibilidad. La cifra que la dirección tiene que autorizar no es
    ninguna de ésas, sino la exposición máxima de tesorería: 335,6 miles de
    euros al cierre de 2028.
  \item Ningún proyecto de la cartera parte por debajo del nivel 6 de madurez
    tecnológica y seis de los ocho parten del 7 o superior. No hay
    investigación en esta hoja de ruta, y no es un defecto: es lo que cabe
    esperar de una tecnología cuya curva S pasó su inflexión hace catorce
    años. Las incertidumbres que quedan son de volumen, de diseño de interfaz,
    de coste y de adopción humana, no de física ni de algoritmo. El panorama
    de publicaciones lo confirma por otra vía: ningún trabajo posterior a 2016
    introduce una propiedad nueva de la tecnología.
  \item La hoja de ruta alcanza el nivel 3 de 5 de madurez, cuantificado, y no
    el 4, medido. La evaluación separa completitud de verificabilidad y
    obtiene 3,67 frente a 2,08: el documento está casi completo en su forma y
    débilmente sustentado en sus cifras, porque ninguna de las ocho figuras de
    mérito está instrumentada hoy y la columna de 2026 es tan estimada como la
    de 2029. Un solo proyecto de ocho meses-persona, R8, cierra esa brecha en
    cuatro elementos a la vez, lo que confirma por una tercera vía la decisión
    de situarlo en la tercera ola y no al final del plan: la hoja de ruta no
    puede verificarse hasta que exista su instrumento de medida, y ese
    instrumento es un proyecto de la propia hoja de ruta.
  \item La debilidad más grave del documento no la resuelve ninguna medida.
    Las seis motivaciones estratégicas de la sección~\ref{sec:motiv} las
    enuncia el autor, y de ellas cuelga toda la cadena de trazabilidad que
    llega hasta los ocho proyectos. Mientras no exista un acta de dirección
    que las suscriba, la hoja de ruta es técnicamente sólida y
    estratégicamente huérfana.
\end{enumerate}

\printbibliography[heading=bibintoc,title={Referencias bibliográficas}]

\end{document}
